\section[Smooth morphisms, unramified (or net) morphisms étale morphisms]{Smooth morphisms, unramified (or net) morphisms\\étale morphisms}
\label{section:IV.17}

In the present section, we take up again the notions studied in (0, 19), expressed in the geometric language of schemes and from the global point of view, for preschemes locally of finite presentation over a given base prescheme. Most of the results (with the exception of n\textsuperscript{os} 17.7, 17.8, 17.9, 17.13 and 17.16) moreover reduce to variants of properties already encountered in (0, 19). For more special results on étale morphisms, the reader should consult § 18.

\subsection{Formally smooth morphisms, formally unramified morphisms, formally étale morphisms}
\EGAReaderTarget{ega4.c04.subsec.17-1}{subsection:IV.17.1}

\begin{definition}[17.1.1]\label{IV.17.1.1}
\EGAReaderTarget{ega4.c04.def.17-1-1}{IV.17.1.1}
Let $f:X\to Y$ be a morphism of preschemes. We say that $f$ is \emph{formally smooth} (resp. \emph{formally unramified} or \emph{formally net}, resp. \emph{formally étale}) if, for every affine scheme $Y'$, every closed subscheme $Y'_0$ of $Y'$ defined by a nilpotent ideal $\sh{J}$ of $\sh{O}_{Y'}$, and every morphism $Y'\to Y$, the map
\[
\label{IV.17.1.1.1}
  \Hom_Y(Y',X)\longrightarrow\Hom_Y(Y'_0,X)
  \tag{17.1.1.1}
\]
induced by the canonical injection $Y'_0\to Y'$, is \emph{surjective} (resp. \emph{injective}, resp. \emph{bijective}).
\end{definition}

We then also say that $X$ is \emph{formally smooth} (resp. \emph{formally unramified} or \emph{formally net}, resp. \emph{formally étale}) over $Y$.

It is clear that to say that $f$ is formally étale means that it is \emph{both} formally smooth and formally unramified.

\begin{remarks}[17.1.2]
\label{IV.17.1.2}
\begin{enumerate}
  \item[(i)] Suppose that $Y=\Spec(A)$ and $X=\Spec(B)$ are affine, so that $f$ comes from a ring homomorphism $\varphi:A\to B$. By virtue of (0, 19.3.1 and 0, 19.10.1), to say that $f$ is formally smooth (resp. formally unramified, resp. formally étale) means that $\varphi$ makes $B$ a \emph{formally smooth} (resp. \emph{formally unramified}, resp. \emph{formally étale}) $A$-algebra for the \emph{discrete} topologies on $A$ and $B$.
  \item[(ii)] To verify that $f$ is formally smooth (resp. formally unramified, resp. formally étale), we may, in Definition (17.1.1), restrict to the case where $\sh{J}^2=0$. Indeed, if $f$ satisfies the corresponding condition of Definition (17.1.1) in this particular case, and if $\sh{J}^n=0$, consider the closed subscheme $Y'_j$ of $Y'$ defined by the ideal $\sh{J}^{j+1}$ for $0\leq j\leq n-1$, so that $Y'_j$ is a closed subscheme of $Y'_{j+1}$ defined by an ideal of square zero; the hypothesis implies that each of the maps
  \[
    \Hom_Y(Y'_{j+1},X)\longrightarrow\Hom_Y(Y'_j,X)qquad(0\leq j\leq n-1)
  \]
  \oldpage[IV-4]{57}
  is surjective (resp. injective, resp. bijective); by composition, we conclude that the same holds for (17.1.1.1).
  \item[(iii)] Note that the properties of the morphism $f$ defined in (17.1.1) are properties of the \emph{representable functor} $(0_{\mathrm{III}}, 8.1.8)$
  \[
    Y'\longmapsto\Hom_Y(Y',X)
  \]
  from the category of $Y$-preschemes to the category of sets; they retain a meaning for \emph{any} contravariant functor having the same source and target categories, whether representable or not.
  \item[(iv)] Suppose that the morphism $f$ is formally unramified (resp. formally étale); consider an \emph{arbitrary} $Y$-prescheme $Z$ and a closed sub-prescheme $Z_0$ of $Z$ defined by a \emph{locally nilpotent} ideal $\sh{J}$ of $\sh{O}_Z$. Then the map
  \[
  \label{IV.17.1.2.1}
    \Hom_Y(Z,X)\longrightarrow\Hom_Y(Z_0,X)
    \tag{17.1.2.1}
  \]
  induced by the canonical injection $Z_0\to Z$, is again injective (resp. bijective). Indeed, let $(U_\alpha)$ be an affine open covering of $Z$ such that the ideals $\sh{J}|U_\alpha$ are nilpotent, and for every $\alpha$, let $U^0_\alpha$ be the inverse image of $U_\alpha$ in $Z_0$, which is the closed subscheme of $U_\alpha$ defined by $\sh{J}|U_\alpha$. Let $f_0:Z_0\to X$ be a $Y$-morphism; by hypothesis, for every $\alpha$, there is at most one (resp. one and only one) $Y$-morphism $f_\alpha:U_\alpha\to X$ whose restriction to $U^0_\alpha$ is equal to $f_0|U^0_\alpha$. We conclude immediately that if $f_\alpha$ and $f_\beta$ are defined, then, for every affine open subset $V\subset U_\alpha\cap U_\beta$, we have $f_\alpha|V=f_\beta|V$, since the restrictions of these morphisms to the inverse image $V_0$ of $V$ in $Z_0$ coincide. There is therefore at most one (resp. one and only one) $Y$-morphism $f:Z\to X$ whose restriction to $Z_0$ coincides with $f_0$.
\end{enumerate}
\end{remarks}

\begin{proposition}[17.1.3]
\label{IV.17.1.3}
\begin{enumerate}
  \item[(i)] A monomorphism of preschemes is formally unramified; an open immersion is formally étale.
  \item[(ii)] The composite of two formally smooth (resp. formally unramified, resp. formally étale) morphisms is formally smooth (resp. formally unramified, resp. formally étale).
  \item[(iii)] If $f:X\to Y$ is a formally smooth (resp. formally unramified, resp. formally étale) $S$-morphism, then so is $f_{(S')}:X_{(S')}\to Y_{(S')}$ for every extension $S'\to S$ of the base prescheme.
  \item[(iv)] If $f:X\to X'$ and $g:Y\to Y'$ are two formally smooth (resp. formally unramified, resp. formally étale) $S$-morphisms, then so is $f\times_S g:X\times_S Y\to X'\times_S Y'$.
  \item[(v)] Let $f:X\to Y$, $g:Y\to Z$ be two morphisms; if $g\circ f$ is formally unramified, then so is $f$.
  \item[(vi)] If $f:X\to Y$ is a formally unramified morphism, then so is $f_{\mathrm{red}}:X_{\mathrm{red}}\to Y_{\mathrm{red}}$.
\end{enumerate}
\end{proposition}

By virtue of (I, 5.5.12), it suffices to prove (i), (ii) and (iii). The two assertions of (i) are trivial. To prove (ii), consider two morphisms $f:X\to Y$, $g:Y\to Z$, an affine scheme $Z'$, a closed subscheme $Z'_0$ of $Z'$ defined by a nilpotent ideal, and a morphism $Z'\to Z$. Suppose $f$ and $g$ formally smooth, and consider a
\oldpage[IV-4]{58}
$Z$-morphism $u_0:Z'_0\to X$; the hypothesis on $g$ implies that there exists a $Z$-morphism $v:Z'\to Y$ such that $f\circ u_0=v\circ j$ (where $j:Z'_0\to Z'$ is the canonical injection); the hypothesis on $f$ then implies that there exists a morphism $u:Z'\to X$ such that $f\circ u=v$ and $u\circ j=u_0$, so $(g\circ f)\circ u$ is equal to the given morphism $Z'\to Z$ and $u\circ j=u_0$, which proves that $g\circ f$ is formally smooth; we argue in the same way when we suppose $f$ and $g$ formally unramified.

Finally, to prove (iii), set $X'=X_{(S')}$, $Y'=Y_{(S')}$, $f'=f_{(S')}$; consider an affine scheme $Y''$, a closed subscheme $Y''_0$ of $Y''$ defined by a nilpotent ideal, and a morphism $g:Y''\to Y'$ making $Y''$ a $Y'$-prescheme; we know then (I, 3.3.8) that $\Hom_{Y'}(Y'',X')$ is canonically identified with $\Hom_Y(Y'',X)$ and $\Hom_{Y'}(Y''_0,X')$ with $\Hom_Y(Y''_0,X)$, and the conclusion then follows immediately from Definition (17.1.1).

Note that a \emph{closed immersion} is not necessarily a formally smooth morphism.

\begin{proposition}[17.1.4]
\label{IV.17.1.4}
Let $f:X\to Y$, $g:Y\to Z$ be two morphisms, and suppose $g$ formally unramified. Then, if $g\circ f$ is formally smooth (resp. formally étale), so is $f$.
\end{proposition}

Indeed, let $Y'$ be an affine scheme, $Y'_0$ a closed subscheme of $Y'$ defined by a nilpotent ideal, $h:Y'\to Y$ a morphism, $j:Y'_0\to Y'$ the canonical injection, and $u_0:Y'_0\to X$ a $Y$-morphism, hence such that $f\circ u_0=h\circ j$. Suppose $g\circ f$ formally smooth; then there exists a morphism $u:Y'\to X$ such that $u\circ j=u_0$ and $(g\circ f)\circ u=g\circ h$. But these relations imply that $f\circ u$ and $h$ are two $Z$-morphisms from $Y'$ to $Y$ such that $(f\circ u)\circ j=h\circ j$; by virtue of the hypothesis that $g$ is formally unramified, we deduce that $f\circ u=h$, in other words $u$ is a $Y$-morphism; hence $f$ is formally smooth. Taking into account (17.1.3, (v)), this proves the proposition.

\begin{corollary}[17.1.5]
\label{IV.17.1.5}
Suppose $g$ formally étale; then, for $g\circ f$ to be formally smooth (resp. formally unramified, resp. formally étale), it is necessary and sufficient that $f$ be so.
\end{corollary}

This follows from (17.1.4) and from (17.1.3, (ii) and (v)).

\begin{proposition}[17.1.6]
\label{IV.17.1.6}
Let $f:X\to Y$ be a morphism of preschemes.
\begin{enumerate}
  \item[(i)] Let $(U_\alpha)$ be an open covering of $X$ and, for every $\alpha$, let $i_\alpha:U_\alpha\to X$ be the canonical injection. For $f$ to be formally smooth (resp. formally unramified, resp. formally étale), it is necessary and sufficient that each of the morphisms $f\circ i_\alpha$ be so.
  \item[(ii)] Let $(V_\lambda)$ be an open covering of $Y$. For $f$ to be formally smooth (resp. formally unramified, resp. formally étale), it is necessary and sufficient that each of the restrictions $f^{-1}(V_\lambda)\to V_\lambda$ of $f$ be so.
\end{enumerate}
\end{proposition}

Let us note first that (ii) is a consequence of (i): indeed, if $j_\lambda:V_\lambda\to Y$ and $i_\lambda:f^{-1}(V_\lambda)\to X$ are the canonical injections, the restriction $f_\lambda:f^{-1}(V_\lambda)\to V_\lambda$ of $f$ is such that $j_\lambda\circ f_\lambda=f\circ i_\lambda$; if $f$ is formally smooth (resp. formally unramified), then so is $f\circ i_\lambda$, since $i_\lambda$ is formally étale (17.1.3); but as $j_\lambda$ is formally étale, this implies that $f_\lambda$ is formally smooth (resp. formally unramified) by virtue of (17.1.5). Conversely, if all the $f_\lambda$ are formally smooth (resp. formally unramified), then so are the $j_\lambda\circ f_\lambda$ (17.1.3), and hence so is $f$, by virtue of (i).

\oldpage[IV-4]{59}
Taking into account that the $i_\alpha$ are formally étale, everything therefore comes down to proving that if the $f\circ i_\alpha$ are formally smooth (resp. formally unramified), then so is $f$.

So let $Y'$ be an affine scheme, $Y'_0$ a closed subscheme of $Y'$ defined by a nilpotent ideal $\sh{J}$, which we may suppose to be such that $\sh{J}^2=0$ (17.1.2, (ii)), and finally let $g:Y'\to Y$ be a morphism. Suppose given a $Y$-morphism $u_0:Y'_0\to X$; denote by $W_\alpha$ (resp. $W^0_\alpha$) the prescheme induced by $Y'$ (resp. $Y'_0$) on the open subset $u_0^{-1}(U_\alpha)$ (recall that $Y'$ and $Y'_0$ have \emph{the same underlying topological space}). Suppose first that the $f\circ i_\alpha$ are \emph{formally unramified}, and let us show that, if $u'$ and $u''$ are two $Y$-morphisms from $Y'$ to $X$ whose restrictions to $Y'_0$ coincide, then $u'=u''$. Indeed, taking into account (17.1.2, (iv)), the hypothesis that the $f\circ i_\alpha$ are unramified implies that for every $\alpha$, we have $u'|W_\alpha=u''|W_\alpha$, since the restrictions of these two $Y$-morphisms to $W^0_\alpha$ coincide. Whence the conclusion in this case.

Now suppose all the $f\circ i_\alpha$ \emph{formally smooth}, and let us prove that there exists a $Y$-morphism $u:Y'\to X$ whose restriction to $Y'_0$ is $u_0$. Now, since $Y'$ is an \emph{affine scheme}, we can apply (16.5.17), whose hypotheses are satisfied and whose conclusion proves precisely the existence of $u$.

We can therefore say that the notions introduced in (17.1.1) are \emph{local} on $X$ and on $Y$, which always makes it possible, by virtue of (17.1.2, (i)), to reduce to the study of formally smooth (resp. formally unramified, resp. formally étale) \emph{algebras}.

\subsection{General differential properties}
\EGAReaderTarget{ega4.c04.subsec.17-2}{subsection:IV.17.2}

\begin{proposition}[17.2.1]
\label{IV.17.2.1}
For a morphism $f:X\to Y$ to be formally unramified, it is necessary and sufficient that $\Omega_f^1=0$ (which is also written $\Omega_{X/Y}^1=0$ (16.3.1)).
\end{proposition}

Taking into account (17.1.6), we are reduced to the case where $Y=\Spec(A)$ and $X=\Spec(B)$ are affine, and the conclusion then follows from (0, 20.7.4) and from the interpretation of $\Omega_{X/Y}^1$ in this case (16.3.7).

\begin{corollary}[17.2.2]
\label{IV.17.2.2}
Let $f:X\to Y$, $g:Y\to Z$ be two morphisms. For $f$ to be formally unramified, it is necessary and sufficient that the canonical homomorphism (16.4.19)
\[
  f^*(\Omega_{Y/Z}^1)\longrightarrow\Omega_{X/Z}^1
\]
be surjective.
\end{corollary}

This is an immediate consequence of (17.2.1) and of the exact sequence (16.4.19.1).

\begin{proposition}[17.2.3]
\label{IV.17.2.3}
Let $f:X\to Y$ be a formally smooth morphism.
\begin{enumerate}
  \item[(i)] The $\sh{O}_X$-module $\Omega_{X/Y}^1$ is locally projective (16.10.1). If $f$ is locally of finite type, $\Omega_{X/Y}^1$ is locally free of finite type.
  \item[(ii)] For every morphism $g:Y\to Z$, the sequence (16.4.19) of $\sh{O}_X$-modules
  \[
  \label{IV.17.2.3.1}
    0\longrightarrow f^*(\Omega_{Y/Z}^1)\longrightarrow\Omega_{X/Z}^1\longrightarrow\Omega_{X/Y}^1\longrightarrow0
    \tag{17.2.3.1}
  \]
  is exact; moreover, for every $x\in X$, there exists an open neighbourhood $U$ of $x$ such that the restrictions to $U$ of the homomorphisms of (17.2.3.1) form a split exact sequence.
\end{enumerate}
\end{proposition}

\oldpage[IV-4]{60}
\begin{enumerate}
  \item[(i)] We know (16.3.9) that if $f$ is locally of finite type, $\Omega_f^1$ is an $\sh{O}_X$-module of finite type. To prove that, in all cases, it is locally projective, we may restrict, by virtue of (17.1.6), to the case where $Y=\Spec(A)$ and $X=\Spec(B)$ are affine, and this follows from the hypothesis on $f$ and from (0, 20.4.9 and 0, 19.2.1).
  \item[(ii)] Here again, we may restrict to the case where $X$, $Y$ and $Z$ are affine (17.1.6), and the conclusion follows in this case from the interpretation of the modules appearing in the sequence (17.2.3.1) and from (0, 20.5.7).
\end{enumerate}

\begin{corollary}[17.2.4]
\label{IV.17.2.4}
If $f:X\to Y$ is a formally étale morphism, then, for every morphism $g:Y\to Z$, the canonical homomorphism of $\sh{O}_X$-modules
\[
  f^*(\Omega_{Y/Z}^1)\longrightarrow\Omega_{X/Z}^1
\]
is bijective.
\end{corollary}

This follows from the exactness of the sequence (17.2.3.1) and from the fact that we then have $\Omega_{X/Y}^1=0$ (17.2.1).

\begin{proposition}[17.2.5]
\label{IV.17.2.5}
Let $f:X\to Y$ be a morphism, and $X'$ a sub-prescheme of $X$ such that the composite morphism $X'\xrightarrow{j}X\to Y$ (where $j$ is the canonical injection) is formally smooth. Then the sequence of $\sh{O}_{X'}$-modules (16.4.21)
\[
\label{IV.17.2.5.1}
  0\longrightarrow\mathcal{N}_{X'/X}\longrightarrow\Omega_{X/Y}^1\otimes_{\sh{O}_X}\sh{O}_{X'}\longrightarrow\Omega_{X'/Y}^1\longrightarrow0
  \tag{17.2.5.1}
\]
is exact; moreover, for every $x\in X$, there exists an open neighbourhood $U$ of $x$ such that the restrictions to $U$ of the homomorphisms of (17.2.5.1) form a split exact sequence.
\end{proposition}

Again by virtue of (17.1.6), we may restrict to the case where $Y=\Spec(A)$ and $X=\Spec(B)$ are affine, and $X'=\Spec(B/\mathfrak{J})$, where $\mathfrak{J}$ is an ideal of $B$. Then the conormal sheaf $\mathcal{N}_{X'/X}$ corresponds to the $B$-module $\mathfrak{J}/\mathfrak{J}^2$ (16.1.3), and the conclusion follows from (0, 20.5.14).

\begin{proposition}[17.2.6]
\label{IV.17.2.6}
Let $X$, $Y$ be two preschemes, and $f:X\to Y$ a morphism locally of finite type. The following conditions are equivalent:
\begin{enumerate}
  \item[a)] $f$ is a monomorphism.
  \item[b)] $f$ is radicial and formally unramified.
  \item[c)] For every $y\in Y$, the fibre $f^{-1}(y)$ is empty or $k(y)$-isomorphic to $\Spec(k(y))$ (in other words, it consists of a single point $z$ such that $k(y)\to\sh{O}_z/\mathfrak{m}_y\sh{O}_z$ is an isomorphism).
\end{enumerate}
\end{proposition}

The fact that a) implies c) follows from (8.11.5.1). It is clear that c) implies that $f$ is radicial; let us show that it also follows from c) that $\Omega_{X/Y}^1=0$, which will prove that c) implies b) (17.2.1). Note that the $\sh{O}_X$-module $\Omega_{X/Y}^1$ is quasi-coherent of finite type (16.3.9). It therefore follows from (I, 9.1.13.1) that, for $(\Omega_{X/Y}^1)_x=0$, it is necessary and sufficient that, if we set $Y_1=\Spec(k(y))$, $X_1=f^{-1}(y)=X\times_Y Y_1$, we have $(\Omega_{X_1/Y_1}^1)_x=0$; but as the morphism $f_1:X_1\to Y_1$ deduced from $f$ is formally unramified by virtue of hypothesis c) (17.1.3), the conclusion follows from (17.2.1). Finally, let us prove that b) implies a); for this, consider the diagonal morphism $g=\Delta_f:X\to X\times_Y X$; since $f$ is radicial, $g$ is surjective (I, 8.7.1); on the other hand, $\Omega_{X/Y}^1$ is by definition the conormal sheaf $\mathop{Gr}\nolimits_1(g)$ of the immersion $g$ (16.3.1), and to say that $f$ is formally unramified therefore means that

\oldpage[IV-4]{61}
$\mathop{Gr}\nolimits_1(g)=0$ (17.2.1). Moreover, $g$ is locally of finite presentation (I, 4.3.1); hence the hypothesis $\mathop{Gr}\nolimits_1(g)=0$ implies that $g$ is an open immersion (16.1.10); being surjective, this immersion is an isomorphism, hence $f$ is a monomorphism (I, 5.3.8).

\subsection{Smooth morphisms, unramified morphisms, étale morphisms}
\EGAReaderTarget{ega4.c04.subsec.17-3}{subsection:IV.17.3}

\begin{definition}[17.3.1]\label{IV.17.3.1}
\EGAReaderTarget{ega4.c04.def.17-3-1}{IV.17.3.1}
We say that a morphism $f:X\to Y$ is \emph{smooth} (resp. \emph{unramified}, or \emph{net}\footnote{The words ``net'' and ``formally net'' seem much preferable to the established terminology of ``unramified'' (resp. ``formally unramified'') and will be used almost exclusively from Chapter V on. In this chapter, we have kept the old terminology so as not to come into conflict with (0, 19.10).}, resp. \emph{étale}) if it is locally of finite presentation and formally smooth (resp. formally unramified, resp. formally étale).
\end{definition}

We then also say that $X$ is \emph{smooth} (resp. \emph{unramified} or \emph{net}, resp. \emph{étale}) over $Y$.

We will see later (17.5.2) that this definition of a smooth morphism coincides with the one already given in (6.8.1); until then, it is the definition of (17.3.1) that we will use exclusively.

It is clear that to say that $f$ is étale means that it is \emph{both} smooth and unramified.

\begin{remarks}[17.3.2]\label{IV.17.3.2}
\EGAReaderTarget{ega4.c04.subsec.17-3-2}{IV.17.3.2}
\begin{enumerate}
  \item[(i)] Note that Definition (17.3.1) can be expressed solely in terms of the functor
  \[
    Y'\longmapsto\Hom_Y(Y',X)
  \]
  considered in (17.1.2, (iii)), since to say that $f$ is locally of finite presentation is equivalent to saying that the preceding functor \emph{commutes with projective limits of affine schemes} (8.14.2).
  \item[(ii)] Let $A$ be a ring, and $B$ an $A$-algebra. We say that $B$ is a \emph{smooth} (resp. \emph{unramified}, resp. \emph{étale}) $A$-algebra if the corresponding morphism $\Spec(B)\to\Spec(A)$ is smooth (resp. unramified, resp. étale). It amounts to the same to say that $B$ is an $A$-algebra \emph{of finite presentation} (I, 4.4.6) and formally smooth (resp. formally unramified, resp. formally étale) for the discrete topologies.
  \item[(iii)] It follows from (17.1.6) and from the definition of a morphism locally of finite presentation (I, 4.4.2) that the notions of smooth, unramified and étale morphism are \emph{local on $X$ and on $Y$}.
\end{enumerate}
\end{remarks}

\begin{proposition}[17.3.3]
\label{IV.17.3.3}
\begin{enumerate}
  \item[(i)] An open immersion is étale. For an immersion to be unramified, it is necessary and sufficient that it be locally of finite presentation.
  \item[(ii)] The composite of two smooth (resp. unramified, resp. étale) morphisms is smooth (resp. unramified, resp. étale).
  \item[(iii)] If $f:X\to Y$ is a smooth (resp. unramified, resp. étale) $S$-morphism, then so is $f_{(S')}:X_{(S')}\to Y_{(S')}$ for every extension $S'\to S$ of the base prescheme.
  \item[(iv)] If $f:X\to X'$ and $g:Y\to Y'$ are two smooth (resp. unramified, resp. étale) $S$-morphisms, then so is $f\times_S g:X\times_S Y\to X'\times_S Y'$.
\end{enumerate}

\oldpage[IV-4]{62}

\begin{enumerate}
  \item[(v)] Let $f:X\to Y$, $g:Y\to Z$ be two morphisms; if $g$ is locally of finite type and $g\circ f$ is unramified, then $f$ is unramified.
\end{enumerate}
\end{proposition}

This follows immediately from (I, 4.3) and from (17.1.3).

\begin{proposition}[17.3.4]
\label{IV.17.3.4}
Let $f:X\to Y$, $g:Y\to Z$ be two morphisms, and suppose $g$ unramified. Then, if $g\circ f$ is smooth (resp. unramified, resp. étale), so is $f$.
\end{proposition}

Indeed, as $g$ and $g\circ f$ are locally of finite presentation, so is $f$ (I, 4.4.3, (v)); the conclusion therefore follows from (17.1.4) and (17.1.3, (v)).

\begin{corollary}[17.3.5]
\label{IV.17.3.5}
Suppose $g$ étale; then, for $f$ to be smooth (resp. unramified, resp. étale), it is necessary and sufficient that $g\circ f$ be so.
\end{corollary}

This follows from (17.3.4) and from (17.3.3, (iii)).

\begin{proposition}[17.3.6]
\label{IV.17.3.6}
Let $g:Y\to S$, $h:X\to S$ be two morphisms locally of finite presentation. For an $S$-morphism $f:X\to Y$ to be unramified, it is necessary and sufficient that the canonical homomorphism (16.4.19)
\[
  f^*(\Omega_{Y/S}^1)\longrightarrow\Omega_{X/S}^1
\]
be surjective.
\end{proposition}

As $f$ is then locally of finite presentation (I, 4.4.3, (v)), the proposition follows immediately from (17.2.2).

\begin{definition}[17.3.7]\label{IV.17.3.7}
\EGAReaderTarget{ega4.c04.def.17-3-7}{IV.17.3.7}
Let $f:X\to Y$ be a morphism. We say that $f$ is \emph{smooth} (resp. \emph{unramified}, resp. \emph{étale}) \emph{at a point} $x\in X$ if there exists an open neighbourhood $U$ of $x$ in $X$ such that the restriction $f|U$ is a smooth (resp. unramified, resp. étale) morphism from $U$ to $Y$.
\end{definition}

We then also say that $X$ is smooth (resp. unramified, resp. étale) over $Y$ at the point $x$.

Taking into account Remark (17.3.2, (iii)), it amounts to the same to say that $f$ is a smooth (resp. unramified, resp. étale) morphism or that it is smooth (resp. unramified, resp. étale) at every point of $X$.

It is clear that the set of points of $X$ where a morphism $f:X\to Y$ is smooth (resp. unramified, resp. étale) is \emph{open} in $X$.

\begin{proposition}[17.3.8]
\label{IV.17.3.8}
For every prescheme $Y$ and every locally free $\sh{O}_Y$-module $\sh{E}$ of finite type, the vector bundle prescheme $\mathbf{V}(\sh{E})$ (II, 1.7.8) is a smooth $Y$-prescheme.
\end{proposition}

Indeed (17.3.2, (iii)), we may restrict to the case where $Y=\Spec(A)$ is affine and $\mathbf{V}(\sh{E})=\Spec(A[T_1,\ldots,T_r])$; as $A[T_1,\ldots,T_r]$ is an $A$-algebra that is formally smooth for the discrete topologies (0, 19.3.2) and of finite presentation, this proves the proposition (17.3.2, (ii)).

\begin{corollary}[17.3.9]
\label{IV.17.3.9}
Under the hypotheses of (17.3.8), the projective bundle prescheme $\mathbf{P}(\sh{E})$ (II, 4.1.1) is a smooth $Y$-prescheme.
\end{corollary}

We may again restrict to the case where $Y=\Spec(A)$ is affine and $\mathbf{P}(\sh{E})=\mathbf{P}^r_Y$. We know then (II, 2.3.14) that we obtain a finite open covering of $\mathbf{P}^r_A$ by taking the $D_+(T_i)$ ($0\leq i\leq r$), which are respectively equal to the spectra of the rings $S_{(f)}$, where $S$ is replaced by $A[T_0,T_1,\ldots,T_r]$ and $f$ by $T_i$; but it follows immediately from the definition of $S_{(f)}$ (II, 2.2.1) that this ring, in the case under consideration, is isomorphic to $A[T_0,\ldots,T_{i-1},T_{i+1},\ldots,T_r]$; hence the corollary follows from (17.3.8).

\oldpage[IV-4]{63}

\subsection{Characterizations of unramified morphisms}
\EGAReaderTarget{ega4.c04.subsec.17-4}{subsection:IV.17.4}

\begin{theorem}[17.4.1]
\label{IV.17.4.1}
Let $f:X\to Y$ be a morphism locally of finite presentation, and $x$ a point of $X$. The following properties are equivalent:
\begin{enumerate}
  \item[a)] $f$ is unramified at the point $x$.
  \item[b)] The diagonal morphism $\Delta_f:X\to X\times_Y X$ is a local isomorphism at the point $x$.
  \item[b')] If we set $\Delta_f=(\psi,\theta)$, $Z=X\times_Y X$ and $z=\psi(x)$, the homomorphism $\theta_z^\#:\sh{O}_{Z,z}\to\sh{O}_{X,x}$ is bijective.
  \item[b'')] For every morphism $g:Y'\to Y$, and every point $y'\in Y'$ above $y=f(x)$, every $Y'$-section $s'$ of $X'=X\times_Y Y'$ such that $x'=s'(y')$ is above $x$ is a local isomorphism at the point $y'$.
  \item[c)] We have $(\Omega^1_{X/Y})_x=0$.
  \item[d)] The $k(y)$-prescheme $f^{-1}(y)$ is unramified over $k(y)$ at the point $x$.
  \item[d')] The point $x$ is isolated in $X_y=f^{-1}(y)$ (in other words (Err$_{\mathrm{III}}$, 20), the morphism $f$ is quasi-finite at the point $x$) and the ring $\sh{O}_{X_y,x}$ is a field which is a separable extension of $k(y)$.
  \item[d'')] The ring $\sh{O}_{X_y,x}=\sh{O}_{X,x}/\mathfrak{m}_y\sh{O}_{X,x}$ is a field which is a finite separable extension of $k(y)$.
  \item[e)] The ring $\sh{O}_{X,x}$ is a formally unramified $\sh{O}_{Y,y}$-algebra for the discrete topologies.
\end{enumerate}
\end{theorem}

As $f$ is locally of finite type, the $\sh{O}_X$-module $\Omega^1_{X/Y}$ is of finite type (16.3.9), so it amounts to the same to say that $(\Omega^1_{X/Y})_x=0$ or that there exists an open neighbourhood $U$ of $x$ such that $\Omega^1_{X/Y}|U=0$. Taking into account (17.2.1), this proves the equivalence of a) and c). On the other hand, if we set $A=\sh{O}_{Y,y}$, $B=\sh{O}_{X,x}$, we have $(\Omega^1_{X/Y})_x=\Omega^1_{B/A}$ (16.4.15), and the equivalence of c) and e) therefore follows from (0, 20.7.4).

As d') involves only properties of the morphism $X_y\to\Spec(k(y))$, the equivalence of a) and d') will \emph{ipso facto} imply that of d) and d'). On the other hand, d') and d'') are equivalent, because it amounts to the same to say that $\sh{O}_{X_y,x}$ is a \emph{finite} $k(y)$-algebra or that $x$ is an isolated point of $X_y$, as $X_y$ is a $k(y)$-prescheme locally of finite type (I, 6.4.4).

Let us now prove the equivalence of b) and b'). We may restrict to the case where $Y=\Spec(R)$ and $X=\Spec(S)$ are affine and $f$ is of finite presentation; then we have $Z=\Spec(S\otimes_R S)$ and $\Delta_f$ corresponds to the canonical surjective homomorphism $S\otimes_R S\to S$, whose kernel $\mathfrak{J}$ is, as we know, an ideal of finite type (0, 20.4.4). If we set $\sh{J}=\widetilde{\mathfrak{J}}$, the $\sh{O}_Z$-module $\psi_*(\sh{O}_X)=\sh{O}_Z/\sh{J}$ is therefore of finite presentation, and the hypothesis that the homomorphism $\theta_z^\#: \sh{O}_{Z,z}\to(\psi_*(\sh{O}_X))_z$ is bijective implies that, on replacing $X$ by an open neighbourhood of $x$ if necessary, the homomorphism $\theta:\sh{O}_Z\to\psi_*(\sh{O}_X)$ is itself bijective ($0_{\mathrm I}$, 5.2.7). This therefore shows that b') implies b); the converse is evident.

On the other hand, the equivalence of b) and b'') follows from (I, 5.3.7) without any finiteness hypothesis on $f$: indeed, to give a $Y'$-section $s':Y'\to X'$ is equivalent to giving a

\oldpage[IV-4]{64}
$Y$-morphism $h=g'\circ s':Y'\to X$ (where $g':X'\to X$ is the canonical projection), so that $s'=(1_{Y'},h)_X$, and then the diagram
\[
\tag{17.4.1.1}
\label{IV.17.4.1.1}
\xymatrix{
Y' \ar[r]^{s'} \ar[d]_h & X'=Y'\times_Y X \ar[d]^{h\times 1_X}\\
X \ar[r]_{\Delta_f} & X\times_Y X
}
\]
identifies $Y'$ with the product of the $(X\times_Y X)$-preschemes $X$ and $X'$. Consequently (I, 4.3.2), if $\Delta_f$ is a local isomorphism at the point $x$, $s'$ is a local isomorphism at the point $y'$ (since $x=h(y')$), which proves that b) implies b''). The converse is obtained by applying b'') to the case where we take $Y'=X$, $y'=x$, $g=f$ and $s'=\Delta_f$.

To complete the proof of (17.4.1), it suffices to prove the implications
\[
d'')\Rightarrow c)\Rightarrow b)\Rightarrow d'').
\]

$d'')\Rightarrow c)$: As $\Omega^1_{X/Y}$ is an $\sh{O}_X$-module of finite type, it follows from Nakayama's lemma that condition c) is equivalent to $(\Omega^1_{X/Y})_x/\mathfrak{m}_y(\Omega^1_{X/Y})_x=0$, that is (16.4.5), to $(\Omega^1_{X_y/\Spec(k(y))})_x=0$. We are therefore reduced to the case where $Y$ is the spectrum of a field $k$ and $X$ an algebraic $k$-prescheme. The hypothesis that $\sh{O}_{X,x}$ is a field $k'$ which is a finite extension of $k$ implies first that $x$ is \emph{closed} in $X$ (I, 6.4.2), and then that $x$ is a \emph{maximal} point of the Noetherian prescheme $X$, hence is an isolated point of $X$. Replacing $X$ by the open subset $\{x\}$ of $X$, we may therefore suppose that $X=\Spec(k')$; but then the hypothesis that $k'$ is a \emph{finite separable} extension of $k$ implies $\Omega^1_{k'/k}=0$ (0, 20.6.20), which proves c).

$c)\Rightarrow b)$: We have seen above that we then have $\Omega^1_{X/Y}|U=0$ for an open neighbourhood $U$ of $x$ in $X$; assertion b) then follows from the definition of $\Omega^1_{X/Y}$ (16.3.1) and from (16.1.9).

$b)\Rightarrow d'')$: Replacing $X$ by an open neighbourhood of $x$, we may suppose that $\Delta_f$ is an open immersion; if we denote by $f_y:X_y\to\Spec(k(y))$ the morphism deduced from $f$ by base change, $\Delta_{f_y}$ is then also an open immersion (I, 5.3.4), and as condition d'') concerns only the prescheme $X_y$, we see that we may restrict to the case where $Y$ is the spectrum of a field $k$ and $X$ the spectrum of a $k$-algebra $A$ of finite type; property d'') will be established if we prove that $A$ is a \emph{finite and separable} $k$-algebra, such an algebra being a direct product of finite separable extensions of $k$. If $K$ is an algebraically closed extension of $k$, it amounts to the same to say that $A\otimes_k K$ is a finite and separable $K$-algebra (4.6.1), so we see that we may restrict to the case where $k$ is algebraically closed. Let us first show that $A$ is a \emph{finite} $k$-algebra: it will suffice to show that every \emph{closed} point $x$ of $X$ is isolated, for then the set of these points is open in $X$ and discrete, hence finite since it is quasi-compact ($X$ being Noetherian), which will establish our assertion by virtue of (I, 6.4.4). Now, we then have $k(x)=k$ since $k$ is algebraically closed (I, 6.4.3), so there is a $Y$-section $s$ of $X$ such that $s(Y)=\{x\}$, and by virtue of (17.4.1.1), $\{x\}$ is the inverse image of the diagonal $\Delta_X(X)$ under a morphism $X\to X\times_Y X$, hence $\{x\}$ is open in $X$ by virtue of hypothesis b). We have thus shown that $A$ is a

\oldpage[IV-4]{65}
\emph{finite} $k$-algebra, a direct product of finite local $k$-algebras. To express that $\Delta_f$ is an open immersion, we may therefore restrict to the case where $A$ is a finite local $k$-algebra, $X=\Spec(A)$ therefore consisting of a single point; the residue field of $A$, being a finite extension of $k$, is necessarily identical to $k$, and consequently (I, 3.4.9) $X\times_k X$ consists of a single point and $\Delta_f$ is therefore necessarily an \emph{isomorphism}. Now, since $A$ is a $k$-algebra, the canonical homomorphism $A\otimes_k A\to A$ can be bijective only if $A=k$. Q.E.D.

\begin{remark}[17.4.1.2]
\label{IV.17.4.1.2}
Suppose only that $f$ is \emph{locally of finite type}. Then $\Omega^1_{X/Y}$ is again an $\sh{O}_X$-module of finite type (16.3.9), and $\Delta_f$ a morphism locally of finite presentation (I, 4.3.1). The whole proof of (17.4.1) is therefore valid, provided that a) is replaced by: \emph{the restriction of $f$ to a suitable neighbourhood of $x$ is a formally unramified morphism}. We see moreover that in this case the restriction of $f$ to a suitable neighbourhood of $x$ is a \emph{locally quasi-finite} morphism.
\end{remark}

\begin{corollary}[17.4.2]
\label{IV.17.4.2}
Let $f:X\to Y$ be a morphism locally of finite presentation. The following properties are equivalent:
\begin{enumerate}
  \item[a)] $f$ is unramified.
  \item[b)] The diagonal morphism $\Delta_f:X\to X\times_Y X$ is an open immersion.
  \item[b')] For every morphism $Y'\to Y$, every $Y'$-section of $X'=X\times_Y Y'$ is an open immersion.
  \item[c)] We have $\Omega^1_{X/Y}=0$.
  \item[d)] For every $y\in Y$, the $k(y)$-prescheme $f^{-1}(y)$ is unramified over $k(y)$.
  \item[d')] For every $y\in Y$, the $k(y)$-prescheme $f^{-1}(y)$ is isomorphic to a prescheme of the form $\coprod_{\lambda\in L}\Spec(K_\lambda)$, where, for every $\lambda\in L$, $K_\lambda$ is a finite and separable extension of $k(y)$.
  \item[e)] For every $x\in X$, the ring $\sh{O}_{X,x}$ is a formally unramified $\sh{O}_{Y,f(x)}$-algebra for the discrete topologies.
\end{enumerate}
\end{corollary}

\begin{corollary}[17.4.3]
\label{IV.17.4.3}
If $f:X\to Y$ is unramified, then $f$ is locally quasi-finite (Err$_{\mathrm{III}}$, 20).
\end{corollary}

\begin{proposition}[17.4.4]
\label{IV.17.4.4}
Let $Y$ be a locally Noetherian prescheme, $f:X\to Y$ a morphism locally of finite type, $x$ a point of $X$, and $y=f(x)$. Set $A=\sh{O}_{Y,y}$, $B=\sh{O}_{X,x}$, which are Noetherian local rings, and let $k$ be the residue field of $A$. Then the equivalent conditions a) to e) of Theorem (17.4.1) are also equivalent to each of the following:
\begin{enumerate}
  \item[f)] $\widehat{B}\otimes_{\widehat{A}}k$ is a field which is a finite separable extension of $k$ (which implies that $\widehat{B}$ is a finite $\widehat{A}$-algebra).
  \item[f')] $\widehat{B}$ is a formally unramified $\widehat{A}$-algebra for the adic topologies.
\end{enumerate}
If moreover $k(x)=k(y)$, or if $k$ is separably closed, these conditions are also equivalent to:
\begin{enumerate}
  \item[f'')] The homomorphism $\widehat{A}\to\widehat{B}$ is surjective.
\end{enumerate}
\end{proposition}

Let us note first, by the same reasoning as in (0, 19.3.6), that it amounts to the same to say that $B$ is a formally unramified $A$-algebra for the preadic topologies, or that $\widehat{B}$ is a formally unramified $\widehat{A}$-algebra for the adic topologies. On the other hand, the hypothesis that $f$ is locally of finite type implies that $\Omega^1_{B/A}$

\oldpage[IV-4]{66}
is a $B$-module of finite type (16.3.9), hence \emph{separated} for the $\mathfrak n$-preadic topology (where $\mathfrak n$ is the maximal ideal of $B$) ($0_{\mathrm I}$, 7.3.5); consequently it amounts to the same to say that $\Omega^1_{B/A}=0$ or that $\widehat{\Omega}^1_{B/A}=0$; hence (0, 20.7.4), it amounts to the same to say that $B$ is a formally unramified $A$-algebra for the \emph{discrete} topologies, or for the \emph{preadic} topologies. This proves the equivalence of conditions e) and f'). If $\mathfrak m$ is the maximal ideal of $A$, we have
\[
k=A/\mathfrak m=\widehat A/\mathfrak m\widehat A,
\qquad
\widehat B\otimes_{\widehat A}k=\widehat B/\mathfrak m\widehat B
=\widehat B\otimes_B(B/\mathfrak m B),
\]
and consequently ($0_{\mathrm I}$, 7.3.5) $\widehat B/\mathfrak m\widehat B$ is the completion of $B/\mathfrak m B=B\otimes_A k$ for the $\mathfrak n$-preadic topology; this proves the equivalence of d'') and f). Finally, when $k(x)=k(y)$ or when $k$ is separably closed, condition f) implies that the homomorphism $\widehat A/\mathfrak m\widehat A\to\widehat B/\mathfrak m\widehat B$ is bijective; condition f) implies on the other hand that $\widehat B$ is a \emph{quasi-finite} $\widehat A$-algebra ($0_{\mathrm I}$, 7.4.4), hence \emph{finite} since $\widehat A$ is complete and $\widehat B$ is separated for the $\mathfrak m$-preadic topology, $\mathfrak m\widehat B$ being an ideal of definition of $\widehat B$ ($0_{\mathrm I}$, 7.4.1). The homomorphism $\widehat A\to\widehat B$ is therefore surjective by virtue of Nakayama's lemma. Hence f) then implies f''), and the converse is evident.

\begin{env}[17.4.5]\label{IV.17.4.5}
\EGAReaderTarget{ega4.c04.statement.17-4-5}{IV.17.4.5}
Given an $S$-prescheme $Y$ and two $S$-morphisms $f:X\to Y$, $g:X\to Y$, we canonically deduce from them an $S$-morphism $(f,g)_S:X\to Y\times_S Y$. We will call the \emph{prescheme of coincidences} of $f$ and $g$ the inverse image under $(f,g)_S$ of the diagonal $\Delta_{Y|S}$; it is therefore a sub-prescheme of $X$, which is \emph{closed} when $Y$ is an $S$-scheme (I, 5.4.1).
\end{env}

\begin{proposition}[17.4.6]
\label{IV.17.4.6}
Let $h:Y\to S$ be an unramified morphism, and let $f:X\to Y$, $g:X\to Y$ be two $S$-morphisms. Then the prescheme $C$ of coincidences of $f$ and $g$ is a sub-prescheme induced on an open subset of $X$; if moreover $Y$ is an $S$-scheme (I, 5.4.1), $C$ is a closed sub-prescheme of $X$.
\end{proposition}

Indeed, since $\Delta_h:Y\to Y\times_S Y$ is an open immersion (17.4.2), the inverse image under $(f,g)_S$ of $\Delta_{Y|S}$ is a sub-prescheme induced on an open subset of $X$ (I, 4.4.1). The last assertion follows from (17.4.5).

\begin{corollary}[17.4.7]
\label{IV.17.4.7}
Under the hypotheses of (17.4.6), let $x$ be a point of $X$ such that the two composite morphisms $\Spec(k(x))\to X\xrightarrow{f}Y$ and $\Spec(k(x))\to X\xrightarrow{g}Y$ are equal. Then there exists an open neighbourhood $U$ of $x$ such that $f|U=g|U$. If moreover $Y$ is an $S$-scheme, there exists an open and closed neighbourhood $X'$ of $x$ in $X$ such that $f|X'=g|X'$. If finally we suppose in addition that $X$ is connected, we have $f=g$.
\end{corollary}

This follows from (17.4.6) and from (I, 5.3.17).

\begin{corollary}[17.4.8]
\label{IV.17.4.8}
Under the hypotheses of (17.4.6), suppose that the structural morphism $\varphi=h\circ f=h\circ g$ from $X$ to $S$ is closed. Let $s$ be a point of $S$; let $X_s$ be the $k(s)$-prescheme $\varphi^{-1}(s)$, and suppose that the two composite morphisms $X_s\to X\xrightarrow{f}Y$ and $X_s\to X\xrightarrow{g}Y$ are equal. Then there exists an open neighbourhood $V$ of $s$ in $S$ such that $f|\varphi^{-1}(V)=g|\varphi^{-1}(V)$. If moreover $Y$ is an $S$-scheme and $\varphi$ is open, we may take $V$ open and closed. If finally we suppose moreover that $S$ is connected, we have $f=g$.
\end{corollary}

It follows from (17.4.7) that the prescheme $C$ of coincidences of $f$ and $g$ is induced on an open subset of $X$ and contains $X_s$. As $\varphi$ is closed, there exists an open neighbourhood $V$ of $s$ such that $\varphi^{-1}(V)\subset C$. If moreover $Y$ is an $S$-scheme, $C$ is closed, so $\varphi(X-C)$ is both

\oldpage[IV-4]{67}
open and closed in $S$, and its complement $V$ in $S$ is therefore an open and closed neighbourhood of $s$ such that $\varphi^{-1}(V)\subset C$.

\begin{proposition}[17.4.9]
\label{IV.17.4.9}
Let $Y$ be a connected prescheme, and $f:X\to Y$ an unramified and separated morphism. Then every $Y$-section $g$ of $X$ is an isomorphism of $Y$ onto an open connected component of $X$, and the map $g\mapsto g(Y)$ is a bijection of $\Gamma(X/Y)$ onto the set of connected components $Z$ of $X$ (necessarily open in $X$) such that the restriction of $f$ to $Z$ is an isomorphism of $Z$ onto $Y$. In particular, if $g'$ and $g''$ are two $Y$-sections of $X$ such that $g'(y)=g''(y)$ for some $y\in Y$, we have $g'=g''$.
\end{proposition}

Indeed, it follows from (17.4.1, b'')) that a $Y$-section $s$ of $X$ is an open immersion, and as $X$ is a $Y$-scheme, $s$ is also a closed immersion (I, 5.4.7); it follows that $s$ is an isomorphism of $Y$ onto a sub-prescheme of $X$ induced on an open and closed subset of $X$, and as $s(Y)$ is connected, it is necessarily a connected component of $X$. The rest of the proposition is immediate.

\begin{remark}[17.4.10]
\label{IV.17.4.10}
Taking into account Remark (17.4.1.2), we see that, in statements (17.4.6) to (17.4.9), we can replace the word ``unramified'' everywhere by ``formally unramified and locally of finite type''.
\end{remark}


\subsection{Characterizations of smooth morphisms}
\EGAReaderTarget{ega4.c04.subsec.17-5}{subsection:IV.17.5}

\begin{theorem}[17.5.1]
\label{IV.17.5.1}
Let $f:X\to Y$ be a morphism locally of finite presentation, $x$ a point of $X$, and $y=f(x)$. The following conditions are equivalent:
\begin{enumerate}
  \item[a)] $f$ is smooth at the point $x$.
  \item[b)] $f$ is flat at the point $x$ and the $k(y)$-prescheme $f^{-1}(y)$ is smooth over $k(y)$ at the point $x$.
  \item[b')] $f$ is regular at the point $x$ (6.8.1).
  \item[c)] The ring $\sh{O}_{X,x}$ is a formally smooth $\sh{O}_{Y,y}$-algebra for the discrete topologies.
\end{enumerate}
\end{theorem}

We may restrict to the case where $Y=\Spec(A)$ and $X=\Spec(C)$, where $C=B/\mathfrak{J}$, $B=A[T_1,\ldots,T_n]$ being a polynomial algebra and $\mathfrak{J}$ an ideal of $B$ of finite type. The equivalence of a) and c) then follows from the equivalence of a) and c) in (0, 22.6.4). On the other hand, applying this result to the morphism locally of finite type $f^{-1}(y)\to\Spec(k(y))$, we see that the equivalence of b) and b') follows from the equivalence of a) and b) in (6.8.6). It therefore remains to prove the equivalence of a) and b).

Let us first show that a) implies b); denote by $\mathfrak p$ the prime ideal $\mathfrak i_x$ in $C$, and by $\mathfrak r$ the prime ideal $\mathfrak i_y$ in $A$; we have $\mathfrak p=\mathfrak q/\mathfrak J$, where $\mathfrak q$ is a prime ideal of $B$, and $\mathfrak r$ is the inverse image of $\mathfrak q$ in $A$. Hypothesis a) implies first that $f^{-1}(y)$ is smooth over $k(y)$ at the point $x$, by (17.3.3), and we must show moreover that $C_{\mathfrak p}=\sh{O}_{X,x}$ is a \emph{flat} $A_{\mathfrak r}$-module. As $C_{\mathfrak p}$ is a formally smooth $A_{\mathfrak r}$-algebra and $B$ is a formally smooth $A$-algebra (for the discrete topologies), the Jacobian criterion (0, 22.6.4), together with (0, 19.1.12), implies that there exist in $\mathfrak J$ a system of $r$ polynomials $u_i$ ($1\leq i\leq r$) and $r$ indices $j_i$ ($1\leq i\leq r$) such that the images of the $u_i$ in $\mathfrak J_{\mathfrak q}/\mathfrak J_{\mathfrak q}^2$ generate this $B_{\mathfrak q}$-module and such that
\[
\label{IV.17.5.1.1}
\det(\partial u_i/\partial T_{j_k})\notin\mathfrak q.
\tag{17.5.1.1}
\]

\oldpage[IV-4]{68}
Let us now note that if $g:\Spec(B)\to\Spec(A)$ is the structural morphism, the fibre $g^{-1}(y)$ is the spectrum of the regular ring $k(y)[T_1,\ldots,T_n]$ (0, 17.3.7), so the Noetherian local ring $B_{\mathfrak q}/\mathfrak rB_{\mathfrak q}$ at a point of this fibre is regular. Now, condition (17.5.1.1) implies that the canonical images $v_i$ of the $u_i$ in the maximal ideal $\mathfrak m$ of $B_{\mathfrak q}/\mathfrak rB_{\mathfrak q}$ are \emph{linearly independent} mod. $\mathfrak m^2$: otherwise, indeed, there would exist polynomials $w_i\in B$ ($1\leq i\leq r$), not all belonging to $\mathfrak q$, such that $\sum_{i=1}^r w_i u_i\in\mathfrak q^2$. Differentiating with respect to the $T_{j_k}$, we would conclude that $\sum_{i=1}^r w_i(\partial u_i/\partial T_{j_k})\in\mathfrak q$ for $1\leq k\leq r$, which would contradict (17.5.1.1) since $\mathfrak q$ is prime. We therefore conclude from (0, 17.1.7) that $(v_i)$ is a \emph{regular sequence} in $B_{\mathfrak q}/\mathfrak rB_{\mathfrak q}$. But as the morphism $g$ is locally of finite presentation and $B$ is a flat $A$-module, it follows from (11.3.8) that the canonical images $u'_i$ of the $u_i$ in $B_{\mathfrak q}$ also form a \emph{regular sequence} and that $B_{\mathfrak q}/(\sum_i u'_iB_{\mathfrak q})$ is a \emph{flat} $A_{\mathfrak r}$-module. As the images of the $u'_i$ in $\mathfrak J_{\mathfrak q}/\mathfrak J_{\mathfrak q}^2$ generate this $B_{\mathfrak q}$-module, it follows from Nakayama's lemma that $\sum_i u'_iB_{\mathfrak q}=\mathfrak J_{\mathfrak q}$, and $C_{\mathfrak p}=B_{\mathfrak q}/\mathfrak J_{\mathfrak q}$ is therefore indeed a flat $A_{\mathfrak r}$-module.

Finally, let us prove that b) implies a). With the same notation, the hypothesis that $B_{\mathfrak q}/\mathfrak J_{\mathfrak q}$ is a flat $A_{\mathfrak r}$-module implies that the canonical homomorphism $\mathfrak J_{\mathfrak q}/\mathfrak r\mathfrak J_{\mathfrak q}\to B_{\mathfrak q}/\mathfrak rB_{\mathfrak q}$ is \emph{injective} ($0_{\mathrm I}$, 6.1.2), so that $\mathfrak J_{\mathfrak q}/\mathfrak r\mathfrak J_{\mathfrak q}$ is identified with an ideal of $B_{\mathfrak q}/\mathfrak rB_{\mathfrak q}$. As $B_{\mathfrak q}/\mathfrak rB_{\mathfrak q}$ is a formally smooth $A_{\mathfrak r}$-algebra for the discrete topologies, we can apply the Jacobian criterion (0, 22.6.4) to $C_{\mathfrak p}=(B_{\mathfrak q}/\mathfrak rB_{\mathfrak q})/(\mathfrak J_{\mathfrak q}/\mathfrak r\mathfrak J_{\mathfrak q})$; together with (0, 19.1.12), this criterion proves, by virtue of hypothesis b), the existence of $r$ polynomials $v_i\in k(y)[T_1,\ldots,T_n]$ such that their images in $(\mathfrak J_{\mathfrak q}/\mathfrak r\mathfrak J_{\mathfrak q})/(\mathfrak J_{\mathfrak q}/\mathfrak r\mathfrak J_{\mathfrak q})^2$ generate this $(B_{\mathfrak q}/\mathfrak rB_{\mathfrak q})$-module and such that
\[
\label{IV.17.5.1.2}
\det(\partial v_i/\partial T_{j_k})\notin\mathfrak qB_{\mathfrak q}/\mathfrak rB_{\mathfrak q}.
\tag{17.5.1.2}
\]
If, for every $i$, we then denote by $u_i$ an element of $\mathfrak J$ whose canonical image is $v_i$, it follows from (17.5.1.2) that the $u_i$ satisfy condition (17.5.1.1); on the other hand, by virtue of Nakayama's lemma, the images $u'_i$ of the $u_i$ in $\mathfrak J_{\mathfrak q}$ generate this $B_{\mathfrak q}$-module. The Jacobian criterion (0, 22.6.4), together with (0, 19.1.12), then proves that $C_{\mathfrak p}=B_{\mathfrak q}/\mathfrak J_{\mathfrak q}$ is a formally smooth $A_{\mathfrak r}$-algebra for the discrete topologies. Q.E.D.

\begin{corollary}[17.5.2]
\label{IV.17.5.2}
Let $f:X\to Y$ be a morphism locally of finite presentation. For $f$ to be smooth (in the sense of (17.3.1)), it is necessary and sufficient that $f$ be regular (6.8.1), in other words that $f$ be flat and that for every $y\in Y$, $f^{-1}(y)$ be a geometrically regular $k(y)$-prescheme (6.7.6).
\end{corollary}

We have thus established the equivalence of the two definitions of ``smooth morphism'' given in (6.8.1) and (17.3.1).

\begin{proposition}[17.5.3]
\label{IV.17.5.3}
Let $Y$ be a locally Noetherian prescheme, $f:X\to Y$ a morphism locally of finite type, $x$ a point of $X$, and $y=f(x)$. Set $A=\sh{O}_{Y,y}$, $B=\sh{O}_{X,x}$, which are Noetherian local rings. Then the equivalent conditions a) to c) of (17.5.1) are also equivalent to each of the following:
\begin{enumerate}
  \item[d)] $B$ is a formally smooth $A$-algebra for the preadic topologies.
  \item[d')] $\widehat B$ is a formally smooth $\widehat A$-algebra for the adic topologies.
\end{enumerate}
If moreover $k(x)=k(y)$, these conditions are also equivalent to:
\begin{enumerate}
  \item[d'')] $\widehat B$ is an $\widehat A$-algebra isomorphic to an algebra of formal power series $\widehat A[[T_1,\ldots,T_n]]$.
\end{enumerate}
\end{proposition}

\oldpage[IV-4]{69}
The equivalence of condition c) of (17.5.1) and of d) follows from the equivalence of a) and d) in the Jacobian criterion (0, 22.6.4), and the equivalence of d) and d') follows from (0, 19.3.6). On the other hand, d'') implies d') without any hypothesis on the residue fields (0, 19.3.4). Finally, if $\mathfrak m$ denotes the maximal ideal of $\widehat A$, hypothesis d') implies that $\widehat B/\mathfrak m\widehat B$ is a complete Noetherian local $k(y)$-algebra which is formally smooth for its adic topology (0, 19.3.5); the hypothesis $k(y)=k(x)$ therefore implies that $\widehat B/\mathfrak m\widehat B$ is $k(y)$-isomorphic to an algebra of formal power series $k(y)[[T_1,\ldots,T_n]]$ (0, 19.6.4). As, on the other hand, $\widehat A[[T_1,\ldots,T_n]]$ is a flat $\widehat A$-module and a complete Noetherian local $\widehat A$-algebra, we conclude from (0, 19.7.1.5) that this algebra is isomorphic to $\widehat B$. Hence d') implies d'') under the additional hypothesis $k(x)=k(y)$.

\begin{remark}[17.5.4]
\label{IV.17.5.4}
Suppose that $Y$ is a locally Noetherian prescheme, and $f:X\to Y$ a morphism locally of finite type. The criterion (17.5.3, d)), together with (0, 22.1.4), shows that to prove that $f$ is smooth, we can apply Definition (17.1.1) while restricting to the case where the affine scheme $Y'$ is the spectrum of an Artinian local ring.
\end{remark}

\begin{proposition}[17.5.5]
\label{IV.17.5.5}
Let $Y$ be a locally Noetherian prescheme, $f:X\to Y$ a morphism locally of finite type, $x$ a point of $X$, and $y=f(x)$. Suppose that $Y$ is reduced at the point $y$. Then, for $f$ to be smooth at the point $x$, it is necessary and sufficient that $f$ be universally open in a neighbourhood of $x$ in $f^{-1}(y)$ and that $f^{-1}(y)$ be a $k(y)$-prescheme geometrically regular at the point $x$.
\end{proposition}

Taking into account (17.5.1), everything comes down to seeing that, if $f^{-1}(y)$ is a $k(y)$-prescheme geometrically regular at the point $x$, it is equivalent to say that $f$ is flat at the point $x$ or that it is universally open in a neighbourhood of $x$ in $f^{-1}(y)$. Now, if $f$ is flat at the point $x$, it is flat in a neighbourhood of $x$ in $X$ (11.1.1), and consequently it is universally open in this neighbourhood (2.4.6). Conversely, the hypothesis that $f$ is universally open in a neighbourhood of $x$ in $f^{-1}(y)$ and that $f^{-1}(y)$ is a $k(y)$-prescheme geometrically regular at the point $x$ implies that $f$ is flat at the point $x$, since $\sh{O}_{Y,y}$ is reduced (15.2.2).

\begin{corollary}[17.5.6]
\label{IV.17.5.6}
Let $Y$ be a locally Noetherian prescheme, $f:X\to Y$ a morphism locally of finite type, $x$ a point of $X$, and $y=f(x)$. Suppose that $Y$ is reduced and geometrically unibranch (6.15.1) at the point $y$. Then, for $f$ to be smooth at the point $x$, it is necessary and sufficient that $f$ be equidimensional at the point $x$ and that $f^{-1}(y)$ be a $k(y)$-prescheme geometrically regular at the point $x$.
\end{corollary}

Noting that the set of points where $f$ is equidimensional is open (13.3.2), we see that the corollary follows from (17.5.5) and from Chevalley's criterion (14.4.4).

The fact that $f$ is a morphism smooth at a point implies in particular that $f$ satisfies

\oldpage[IV-4]{70}
at this point all the properties defined in (6.8.1). We therefore have the following properties, which we recall for convenience of reference:

\begin{proposition}[17.5.7]
\label{IV.17.5.7}
Let $f:X\to Y$ be a morphism locally of finite presentation, smooth at a point $x\in X$; set $y=f(x)$. Then, for the ring $\sh{O}_{X,x}$ to be reduced (resp. integrally closed, resp. geometrically unibranch), it is necessary and sufficient that $\sh{O}_{Y,y}$ be so.
\end{proposition}

Indeed, this was proved in (11.3.13) and (11.3.14), supplemented by Err$_{\mathrm{IV}}$, 30.

\begin{proposition}[17.5.8]
\label{IV.17.5.8}
Let $Y$ be a locally Noetherian prescheme, and $f:X\to Y$ a morphism locally of finite type, smooth at a point $x\in X$; set $y=f(x)$. Then:
\begin{enumerate}
  \item[(i)] We have $\dim(\sh{O}_{X,x})=\dim(\sh{O}_{Y,y})+\dim(\sh{O}_{X,x}\otimes_{\sh{O}_{Y,y}}k(y))$.
  \item[(ii)] We have $\operatorname{coprof}(\sh{O}_{X,x})=\operatorname{coprof}(\sh{O}_{Y,y})$.
  \item[(iii)] For the ring $\sh{O}_{X,x}$ to have property $(S_n)$ (5.7.2) (resp. $(R_n)$ (5.8.2)), it is necessary and sufficient that the ring $\sh{O}_{Y,y}$ have it. In particular, for $\sh{O}_{X,x}$ to be regular, it is necessary and sufficient that $\sh{O}_{Y,y}$ be so.
\end{enumerate}
\end{proposition}

These are particular cases of (6.1.2), (6.3.2), (6.4.1) and (6.5.3).

\subsection{Characterizations of étale morphisms}
\EGAReaderTarget{ega4.c04.subsec.17-6}{subsection:IV.17.6}

\begin{theorem}[17.6.1]
\label{IV.17.6.1}
Let $f:X\to Y$ be a morphism locally of finite presentation, $x$ a point of $X$, and $y=f(x)$. The following conditions are equivalent:
\begin{enumerate}
  \item[a)] $f$ is étale at the point $x$.
  \item[a')] $f$ is smooth at the point $x$ and unramified at the point $x$.
  \item[b)] $f$ is smooth at the point $x$ and quasi-finite at the point $x$ (Err$_{\mathrm{III}}$, 20).
  \item[c)] $f$ is flat at the point $x$ and unramified at the point $x$.
  \item[c')] $f$ is flat at the point $x$ and the ring $\sh{O}_{X,x}/\mathfrak m_y\sh{O}_{X,x}$ is a field which is a finite separable extension of $k(y)$.
  \item[d)] The ring $\sh{O}_{X,x}$ is a formally étale $\sh{O}_{Y,y}$-algebra for the discrete topologies.
\end{enumerate}
\end{theorem}

The equivalence of a) and a') follows immediately from the definitions; that of a) and d) follows from the equivalence of a) and e) in (17.4.1) and from the equivalence of a) and d) in (17.5.1). The equivalence of c) and c') follows from the equivalence of a) and d') in (17.4.1). The fact that a') implies c') follows from (17.5.1); conversely, if c') is satisfied, $f$ is regular (hence smooth by (17.5.1)) at the point $x$, because if $K$ is a finite separable extension of a field $k$, then, for every extension $k'$ of $k$, $\Spec(K\otimes_k k')$ is regular, being a sum of finitely many spectra of fields. The fact that a') implies b) follows from (17.4.1, d')) and from (17.5.1, b)). It therefore remains to see that b) implies c), and as we already know that $f$ is flat at the point $x$, by (17.5.1), it suffices to show that $f^{-1}(y)$ is a $k(y)$-prescheme unramified over $k(y)$; in other words, we are reduced to proving that b) implies c) when $Y=\Spec(k)$ is the spectrum of a field $k$. As the question is local on $X$, we may restrict to the case where $X=\Spec(A)$, where $A$ is a finite and local $k$-algebra ($0_{\mathrm I}$, 7.4.1). By virtue of hypothesis b), $A$ is a formally smooth $k$-algebra for the discrete topologies, which here coincide with the preadic topologies; hence (0, 19.6.5), $A$ is a regular local ring, hence

\oldpage[IV-4]{71}
a field since it is Artinian, and it then follows from (0, 19.6.5.1) that $A$ must be a finite and separable extension of $k$, which completes the proof (17.4.1).

\begin{corollary}[17.6.2]
\label{IV.17.6.2}
Let $f:X\to Y$ be a morphism locally of finite presentation. The following conditions are equivalent:
\begin{enumerate}
  \item[a)] $f$ is étale.
  \item[a')] $f$ is smooth and unramified.
  \item[b)] $f$ is smooth and locally quasi-finite (Err$_{\mathrm{III}}$, 20).
  \item[c)] $f$ is flat and unramified.
  \item[c')] $f$ is flat, and every fibre $f^{-1}(y)$ is a $k(y)$-scheme which is a sum of spectra of fields that are finite and separable extensions of $k(y)$.
  \item[c'')] $f$ is flat, and for every $y\in Y$ and every algebraically closed extension $k'$ of $k(y)$, the ``geometric fibre'' $f^{-1}(y)\otimes_{k(y)}k'$ is a sum of spectra of fields isomorphic to $k'$.
\end{enumerate}
\end{corollary}

The only point that remains to be proved is the equivalence of c') and c''). It is clear that c') implies c'') by base change (17.3.3). On the other hand, as the projection morphism $f^{-1}(y)\otimes_{k(y)}k'\to f^{-1}(y)$ is open (2.4.10), hypothesis c'') implies that the space $f^{-1}(y)$ is discrete; hence, for every point $x\in X_y=f^{-1}(y)$, the local ring $\sh{O}_{X_y,x}=A$ is a finite $k(y)$-algebra, hence an Artinian local ring; moreover it follows from c'') that $\Spec(A\otimes_{k(y)}k')$ is a sum of spectra of fields isomorphic to $k'$, which is possible only if $A$ is a field which is a finite separable extension of $k(y)$ (4.6.1).

\begin{proposition}[17.6.3]
\label{IV.17.6.3}
Let $Y$ be a locally Noetherian prescheme, $f:X\to Y$ a morphism locally of finite type, $x$ a point of $X$, and $y=f(x)$. Set $A=\sh{O}_{Y,y}$, $B=\sh{O}_{X,x}$, which are Noetherian local rings, and let $k$ be the residue field of $A$. Then the equivalent conditions a) to d) of (17.6.1) are also equivalent to each of the following:
\begin{enumerate}
  \item[e)] $\widehat B$ is a formally étale $\widehat A$-algebra for the adic topologies.
  \item[e')] $\widehat B$ is a free $\widehat A$-module and $\widehat B\otimes_{\widehat A}k$ is a field which is a finite and separable extension of $k$ (which implies that $\widehat B$ is a finite $\widehat A$-algebra).
\end{enumerate}
If moreover $k(x)=k(y)$, or if $k$ is separably closed, these conditions are also equivalent to:
\begin{enumerate}
  \item[e'')] The canonical homomorphism $\widehat A\to\widehat B$ is bijective.
\end{enumerate}
\end{proposition}

The equivalence of e) and each of the conditions of (17.6.1) follows immediately from (17.4.4, f')) and (17.5.3, d')). The fact that e) implies e') follows from (17.4.4, f)) and from (0, 19.7.1), taking into account that $\widehat B$ is then a finite $\widehat A$-algebra (17.4.4), so that it amounts to the same to say that $\widehat B$ is a flat $\widehat A$-module or a free $\widehat A$-module ($0_{\mathrm{III}}$, 10.1.3). Conversely, the fact that e') implies e) follows from (17.4.4) and from (0, 19.7.1). Finally, e') implies that the homomorphism $\widehat A\to\widehat B$ is injective, and if $k(x)=k(y)$ or if $k$ is separably closed, this homomorphism is surjective by (17.4.4). The converse is immediate.

\begin{proposition}[17.6.4]
\label{IV.17.6.4}
Under the hypotheses of (17.6.3), if $f$ is étale at the point $x$, we have $\dim(\sh{O}_{X,x})=\dim(\sh{O}_{Y,y})$.
\end{proposition}

This is a particular case of (17.5.8, (i)), since $x$ is isolated in its fibre $f^{-1}(y)$.

\oldpage[IV-4]{72}

\subsection{Properties of descent, of passage to the limit and of constructibility}
\EGAReaderTarget{ega4.c04.subsec.17-7}{subsection:IV.17.7}
\label{IV.17.7}

\begin{proposition}[17.7.1]
\label{IV.17.7.1}
Let $f:X\to Y$ be a morphism locally of finite presentation, $g:Y'\to Y$ a morphism, $X'=X\times_Y Y'$, and $f'=f_{(Y')}:X'\to Y'$ and $g':X'\to X$ the canonical projections. Let $x'$ be a point of $X'$, and set $x=g'(x')$, $y'=f'(x')$.
\begin{enumerate}
  \item[(i)] If $f'$ is unramified at the point $x'$, then $f$ is unramified at the point $x$.
  \item[(ii)] Suppose moreover that $g$ is flat at the point $y'$. Then, if $f'$ is smooth (resp. étale) at the point $x'$, $f$ is smooth (resp. étale) at the point $x$.
\end{enumerate}
\end{proposition}

Set $y=f(x)=g(y')$, so that we have $f'^{-1}(y')=f^{-1}(y)\otimes_{k(y)}k(y')$; note that $f'$ is locally of finite presentation.

(i) As the property, for a morphism (locally of finite presentation), of being unramified at a point involves only the fibre of the morphism at this point (17.4.1, $d$)), we may restrict to the case where $Y$ and $Y'$ are spectra of fields. But then it amounts to the same to say that $f$ (resp. $f'$) is unramified at $x$ (resp. $x'$) or that it is étale at this point (17.6.1, $c$)), so (i) is a consequence of (ii).

(ii) As $f'$ is flat at the point $x'$ (17.5.1), the hypothesis that $g$ is flat at the point $y'$ implies that $f$ is flat at the point $x$, because the projection $g':X'\to X$ is a morphism flat at the point $x'$, and $f\circ g'=g\circ f'$ is a morphism flat at the point $X'$, whence the conclusion (2.2.11, (iv)). The fact that $f$ is smooth (resp. étale) at the point $x$ then involves only the fibre $f^{-1}(y)$ (17.5.1 and 17.6.1), and we are therefore again reduced to the case where $Y=\Spec(k)$ and $Y'=\Spec(k')$ are spectra of fields. To say that $f'$ is smooth at the point $x'$ then means (17.5.1) that $X'$ is a $k'$-prescheme geometrically regular at the point $x'$, and this implies (6.7.8) that $X$ is a $k$-prescheme geometrically regular at the point $x$, hence that $f$ is smooth at the point $x$. Suppose moreover that $f'$ is étale at the point $x'$, so that $x'$ is isolated in $f'^{-1}(y')$ (17.6.1); as the projection $f'^{-1}(y')\to f^{-1}(y)$ is an open morphism (2.4.10), $x$ is isolated in $f^{-1}(y)$; as we already know that $f$ is smooth at the point $x$, it is étale at this point (17.6.1).

\begin{corollary}[17.7.2]
\label{IV.17.7.2}
\begin{enumerate}
  \item[(i)] With the notation of (17.7.1), let $U$ (resp. $U'$) be the set of points of $X$ (resp. $X'$) where $f$ (resp. $f'$) is unramified; then we have $U'=g'^{-1}(U)$.
  \item[(ii)] Suppose moreover that $g$ is flat, and let $V$ (resp. $V'$) be the set of points of $X$ where $f$ (resp. $f'$) is smooth; then $V'=g'^{-1}(V)$.
\end{enumerate}
\end{corollary}

\begin{corollary}[17.7.3]
\label{IV.17.7.3}
\begin{enumerate}
  \item[(i)] Suppose $g$ surjective; then, for $f$ to be unramified, it is necessary and sufficient that $f'$ be so.
  \item[(ii)] Let $f:X\to Y$ be an $S$-morphism, and $g:S'\to S$ a faithfully flat morphism. Suppose that $f$ is locally of finite presentation, or that $g$ is quasi-compact. Then, for $f$ to be smooth (resp. unramified, resp. étale), it is necessary and sufficient that $f'$ be so.
\end{enumerate}
\end{corollary}

In (ii), the case where $f$ is locally of finite presentation follows from (17.7.1). If $g$ is quasi-compact and $f'$ is smooth (resp. unramified, resp. étale), $f'$ is locally of finite presentation by (2.7.1, (iv)), and we are reduced to the first case.

\begin{proposition}[17.7.4]
\label{IV.17.7.4}
Let $f:X\to Y$ be a morphism, $g:Y'\to Y$ a flat morphism locally of finite presentation, $X'=X\times_Y Y'$, and $f'=f_{(Y')}:X'\to Y'$ and $g':X'\to X$ the

\oldpage[IV-4]{73}

canonical projections. Let $V$ (resp. $V'$) be the set of $x\in X$ (resp. $x'\in X'$) at which $f$ has one of the following properties (resp. at which $f'$ has the same property): being:
\begin{enumerate}
  \item[(i)] locally of finite type;
  \item[(ii)] locally of finite presentation;
  \item[(iii)] flat;
  \item[(iv)] unramified;
  \item[(v)] smooth;
  \item[(vi)] étale.
\end{enumerate}
Then we have $V'=g'^{-1}(V)$ (in other words, for $f'$ to have the property under consideration at a point $x'$, it is necessary and sufficient that $f$ have this property at the point $x$).
\end{proposition}

Property (iii) is included only for the record, and does not require the hypothesis that $g$ be locally of finite presentation ((2.2.11, (iv)), taking into account the fact that the projection $g':X'\to X$ is a flat morphism). By virtue of (17.7.2), the assertions relating to properties (iv), (v) and (vi) are consequences of the assertion relating to (ii). It therefore suffices to consider cases (i) and (ii). It is clear that $V'\supset g'^{-1}(V)$; it therefore remains to show that $V'\subset g'^{-1}(V)$; note that the sets $V$ and $V'$ are open, and $W=g'(V')$ is also open in $X$ by virtue of (2.4.6). We must therefore prove that the morphism $f|W:W\to Y$ is locally of finite type (resp. locally of finite presentation); as by hypothesis the composite
\[
V'\xrightarrow{g''}W\xrightarrow{f|W}Y
\]
(where $g''$ is the restriction of $g'$), which is equal to $V'\xrightarrow{f'}Y'\xrightarrow{g}Y$, is locally of finite type (resp. locally of finite presentation) and $g''$ is surjective (hence faithfully flat), we are reduced to proving the following lemma, which improves (11.3.16):

\begin{lemma}[17.7.5]
\label{IV.17.7.5}
Let $f:X\to Y$ be a faithfully flat morphism locally of finite presentation, and $g:Y\to Z$ a morphism such that $g\circ f:X\to Z$ has one of the following properties: being:
\begin{enumerate}
  \item[(i)] locally of finite type;
  \item[(ii)] locally of finite presentation;
  \item[(iii)] of finite type.
\end{enumerate}
Then $g$ has the same property. If moreover $f$ is quasi-compact or $g$ is quasi-separated, the same conclusion is valid for the property:
\begin{enumerate}
  \item[(iv)] being of finite presentation.
\end{enumerate}
\end{lemma}

In cases (i) and (ii), we must see that for every $y\in Y$, there are an affine open neighbourhood $V$ of $y$ in $Y$ and an affine open neighbourhood $W$ of $z=g(y)$ in $Z$ containing $g(V)$, such that the morphism $V\to W$, restriction of $g$, is of finite type (resp. of finite presentation). Now, by hypothesis there exist an $x\in X$ such that $f(x)=y$, an affine open neighbourhood $U$ of $x$ in $X$, an affine open neighbourhood $V'$ of $y$ in $Y$ containing $f(U)$, and an affine open neighbourhood $W$ of $z$ in $Z$ containing $g(V')$, such that the morphism $f_1:U\to V'$, restriction of $f$, is flat and of finite presentation, and the morphism $g_1:V'\to W$, restriction of $g$, is such that $g_1\circ f_1$ is of finite type (resp. of finite presentation). Then $f(U)$ is open in $Y$ (2.4.6), and if $V\subset f(U)$ is an affine neighbourhood of $y$, the morphism $f_2:f_1^{-1}(V)\to V$,

\oldpage[IV-4]{74}

restriction of $f_1$, is again of finite presentation and is moreover faithfully flat; furthermore, if $g_2=g_1|V$, $g_2\circ f_2$ is of finite type (resp. of finite presentation), the open subset $f_1^{-1}(V)$ being quasi-compact (resp. quasi-compact and quasi-separated). We are then reduced to the hypotheses of (11.3.16), from which we conclude that $g_2$ is a morphism of finite type (resp. of finite presentation).

In case (iii), the question is local on $Z$, so we may suppose $Z$ affine, and it then follows that $X$ is quasi-compact, hence so is $Y=f(X)$. Case (iii) is therefore a consequence of (i).

In case (iv) (with the additional hypotheses on $f$ or $g$), we may also suppose $Z$ affine, hence $X$ and $Y$ quasi-compact; moreover, we already know that $g$ is locally of finite presentation and quasi-compact (1.1.3), so everything comes down to seeing that $g$ is quasi-separated, and it therefore suffices to show that this property holds when $f$ is supposed quasi-compact. Now, since $g\circ f$ is quasi-separated, so is $f$ (1.2.2, (v)); as $f$ is quasi-compact and locally of finite presentation, it is of finite presentation (1.6.1); it then suffices to repeat the reasoning of the first paragraph of the proof of (11.3.16).

\begin{remark}[17.7.6]
\label{IV.17.7.6}
In the assertion relating to property (iv), we cannot drop the hypothesis that $f$ is quasi-compact; otherwise, this would imply that every quasi-compact morphism $g:Y\to Z$ locally of finite presentation would be of finite presentation, a conclusion which we know to be false (1.6.4). Indeed, we may restrict to the case where $Z$ is affine, hence $Y$ quasi-compact; consequently there is a finite covering $(U_i)$ of $Y$ by affine open subsets such that the restrictions $g|U_i$ are of finite presentation; it would then suffice to take for $X$ the prescheme \emph{sum} of the $U_i$, and for $f:X\to Y$ the canonical morphism, which is evidently faithfully flat and locally of finite presentation (1.4.3); $g\circ f$ would be of finite presentation by virtue of the choice of the $U_i$ and of (1.6.5), whence our assertion.
\end{remark}

\begin{proposition}[17.7.7]
\label{IV.17.7.7}
Let $f:Y\to S$ and $h:X\to S$ be two morphisms locally of finite presentation, $g:X\to Y$ an $S$-morphism, $x$ a point of $X$, and $y=g(x)$. Suppose that $g$ is flat at the point $x$. Then, if $h$ is smooth (resp. unramified, resp. étale) at the point $x$, $f$ is smooth (resp. unramified, resp. étale) at the point $y$.
\end{proposition}

Set $s=h(x)=f(y)$. To say that $h$ is unramified at the point $x$ (resp. that $f$ is unramified at the point $y$) is equivalent to saying that $h^{-1}(s)$ is étale over $k(s)$ at the point $x$ (resp. that $f^{-1}(s)$ is étale over $k(s)$ at the point $y$) (17.4.1 and 17.6.1); as the morphism $g_s:h^{-1}(s)\to f^{-1}(s)$ deduced from $g$ is flat at the point $x$, we see that we may restrict to proving the proposition when $h$ is smooth or étale at the point $x$. Moreover, as $h$ is then flat at the point $x$ (17.5.1), $f$ is flat at the point $y$, as follows from (2.2.11, (iv)). It therefore amounts to the same (17.5.1) to say that $f$ is smooth (resp. étale) at the point $y$, or that $f^{-1}(s)$ is smooth (resp. étale) over $k(s)$ at the point $x$. We are thus reduced to the case where $S=\Spec(k)$ is the spectrum of a field.

(i) \emph{Case of smooth morphisms}. As $g$ is a morphism locally of finite presentation (1.4.3, (v)), there is an open neighbourhood $U$ of $x$ in $X$ in which $g$ is flat (11.3.1) and $h$ is smooth. Moreover, $g(U)$ is an open neighbourhood of $y$ in $Y$ (2.4.6); replacing $X$ by $U$ and $Y$ by $g(U)$, we may therefore suppose that $g$ is \emph{faithfully flat} and that $h$ is smooth, and we are reduced to proving that $f$ is then smooth. If $k'$ is an algebraically

\oldpage[IV-4]{75}

closed extension of $k$, $X\otimes_k k'$ is then smooth over $k'$, and by virtue of (17.7.3, (ii)), it suffices to prove that $Y\otimes_k k'$ is smooth over $k'$ (since $\Spec(k')\to\Spec(k)$ is faithfully flat); we may therefore restrict to the case where $k$ is \emph{algebraically closed}. As the set of points of $Y$ rational over $k$ is then very dense in $Y$ (10.4.8), and as the set of points of $Y$ where $Y$ is smooth over $k$ is open, we see that it suffices to prove that $Y$ is smooth over $k$ at every point rational over $k$. But at such a point $y$, to say that $Y$ is smooth over $k$ at this point is equivalent to saying that $Y$ is \emph{regular} at $y$ (17.5.1 and 6.7.8). Now we have $y=g(x)$ for some $x\in X$, and by hypothesis (17.5.1) $X$ is regular at the point $x$; as $X$ and $Y$ are then locally Noetherian and $g$ is flat, $Y$ is indeed regular at the point $y$ (6.5.1, (i)).

(ii) \emph{Case of étale morphisms}. By (i) we already know that $f$ is smooth at the point $y$; by virtue of (17.6.1), it therefore suffices to show that $f$ is quasi-finite at the point $y$, or equivalently that $\sh{O}_{Y,y}$ is a finite $k$-algebra. Now, as $\sh{O}_{X,x}$ is a faithfully flat $\sh{O}_{Y,y}$-module ($0_{\mathrm I}$, 6.6.2), $\sh{O}_{Y,y}$ is identified with a sub-$k$-module of $\sh{O}_{X,x}$, and as $\sh{O}_{X,x}$ is by hypothesis a finite $k$-algebra, so is $\sh{O}_{Y,y}$.

\begin{proposition}[17.7.8]
\label{IV.17.7.8}
The notation being that of (8.8.1), suppose $X_\alpha$ and $Y_\alpha$ locally of finite presentation over $S_\alpha$. Let $f_\alpha:X_\alpha\to Y_\alpha$ be an $S_\alpha$-morphism, and $f:X\to Y$ the corresponding $S$-morphism.
\begin{enumerate}
  \item[(i)] Let $x$ be a point of $X$, and $x_\lambda$ its canonical projection in $X_\lambda$. For $f$ to be smooth (resp. unramified, resp. étale) at the point $x$, it is necessary and sufficient that there exist $\lambda\geq\alpha$ such that $f_\lambda$ is smooth (resp. unramified, resp. étale) at the point $x_\lambda$.
  \item[(ii)] Suppose moreover $X_\alpha$ quasi-compact. For $f$ to be smooth (resp. unramified, resp. étale), it is necessary and sufficient that there exist $\lambda\geq\alpha$ such that $f_\lambda$ is smooth (resp. unramified, resp. étale).
\end{enumerate}
\end{proposition}

(i) If $y=f(x)$, $y_\lambda=f_\lambda(x_\lambda)$ is the canonical projection of $y$ in $Y_\lambda$, and we have $f^{-1}(y)=f_\lambda^{-1}(y_\lambda)\otimes_{k(y_\lambda)}k(y)$; the part of the statement concerning unramified morphisms thus follows from (17.7.1, (i)), and it therefore suffices to consider the case of smooth morphisms. As $f$ and $f_\lambda$ are locally of finite presentation, it amounts to the same to say that $f^{-1}(y)$ is geometrically regular at the point $x$ or that $f_\lambda^{-1}(y_\lambda)$ is geometrically regular at the point $x_\lambda$ (6.7.8). The proposition therefore follows from (17.5.1) and from (11.2.6).

(ii) For every $\lambda$, let $U_\lambda$ be the open set of $x_\lambda\in X_\lambda$ such that $f_\lambda$ is smooth (resp. unramified, resp. étale) at the point $x_\lambda$, and let $V_\lambda$ be its inverse image in $X$. As, by hypothesis, for every $x\in X$ there exists, by virtue of (i), a $\lambda$ such that $f_\lambda$ is smooth (resp. unramified, resp. étale) at the point $x_\lambda$, $X$ is the union of the $V_\lambda$. Moreover (17.3.3), for $\lambda\leq\mu$ we have $V_\lambda\subset V_\mu$, so, as $X$ is quasi-compact, there exists an index $\mu$ such that $X=V_\mu$. As the $X_\lambda$ are quasi-compact, it then follows from (8.3.4) that there exists an index $\nu\geq\mu$ such that the inverse image of $U_\mu$ in $X_\nu$ is the whole of $X_\nu$, which means that $f_\nu$ is smooth (resp. unramified, resp. étale) by (17.3.3).

\begin{corollary}[17.7.9]
\label{IV.17.7.9}
Let $S=\Spec(A)$ be an affine scheme, and $f:X\to S$ a morphism. The following conditions are equivalent:
\begin{enumerate}
  \item[a)] $f$ is a morphism of finite presentation and smooth (resp. unramified, resp. étale).

  \oldpage[IV-4]{76}

  \item[b)] There exist a Noetherian affine scheme $S_0=\Spec(A_0)$, a morphism of finite type $f_0:X_0\to S_0$, and a morphism $S\to S_0$ such that the $S$-prescheme $X_0\otimes_{S_0}S$ is $S$-isomorphic to $X$ and $f_0$ is smooth (resp. unramified, resp. étale).
  \item[c)] The conditions of b) are satisfied, and moreover $A_0$ is a sub-$\mathbb{Z}$-algebra of finite type of $A$, the morphism $S\to S_0$ corresponding to the canonical injection $A_0\to A$.
\end{enumerate}
\end{corollary}

The proof from (17.7.8) is the same as that of (11.2.7) from (11.2.6).

\begin{proposition}[17.7.10]
\label{IV.17.7.10}
Let $f:Y\to S$, $h:X\to S$ be two morphisms locally of finite presentation, $g:X\to Y$ an $S$-morphism, $x$ a point of $X$, and $y=g(x)$. Suppose that $h$ is flat at the point $x$ and that $f$ is unramified at the point $y$. Then $f$ is étale at the point $y$ and $g$ is flat at the point $x$.
\end{proposition}

(We will see later (18.4.9) that we can in fact dispense with the hypothesis that $h$ is locally of finite presentation).

The question being local on $S$, $X$ and $Y$, we may suppose that $S$, $X$ and $Y$ are affine, $f,g,h$ morphisms of finite presentation, $h$ flat and $f$ unramified. Taking into account (11.2.7) and (17.7.9), we may suppose moreover that $S$, $X$ and $Y$ are Noetherian; finally, we may restrict to the case where $S=\Spec(\sh{O}_{S,s})$ (where $s=f(y)=h(x)$). The ring $A=\sh{O}_{S,s}$ being a Noetherian local ring, there exist a complete Noetherian local ring $B$ with algebraically closed residue field and a local homomorphism $A\to B$ making $B$ a faithfully flat $A$-module ($0_{\mathrm{III}}$, 10.3.1). Replacing $X$ and $Y$ by $X\times_S\Spec(B)$ and $Y\times_S\Spec(B)$, we conclude by (2.5.1) and (17.7.1) that we may restrict to the case where $A$ is complete and has an algebraically closed residue field. The hypothesis that $f$ is unramified at the point $y$ then implies (17.4.4) that $\sh{O}_{Y,y}$ is a finite $\sh{O}_{S,s}$-algebra and a complete Noetherian local ring ($0_{\mathrm I}$, 7.4.2), and as the residue field of $\sh{O}_{S,s}$ is algebraically closed, the homomorphism $\sh{O}_{S,s}\to\sh{O}_{Y,y}$ is surjective (17.4.4). But, on the other hand, the hypothesis that $h$ is flat at the point $x$ implies that the composite homomorphism $\sh{O}_{S,s}\to\sh{O}_{Y,y}\to\sh{O}_{X,x}$ is injective ($0_{\mathrm I}$, 6.5.1), hence the homomorphism $\sh{O}_{S,s}\to\sh{O}_{Y,y}$ is bijective, which shows that $f$ is étale at the point $y$ (17.6.3); moreover, $\sh{O}_{X,x}$ is a flat $\sh{O}_{Y,y}$-module, hence $g$ is flat at the point $x$.

\begin{proposition}[17.7.11]
\label{IV.17.7.11}
Let $S$ be a prescheme, $X$, $Y$ two $S$-preschemes locally of finite presentation over $S$, and $f:X\to Y$ an $S$-morphism. For every $s\in S$, we denote by $X_s$, $Y_s$, $f_s$ the preschemes and the morphism deduced from $X$, $Y$, $f$ by the base change $\Spec(k(s))\to S$. Then:
\begin{enumerate}
  \item[(i)] The set of $x\in X$ such that, if $s$ is the image of $x$ in $S$, $f_s:X_s\to Y_s$ is smooth (resp. unramified, resp. étale, resp. differentially smooth) at the point $x$, is locally constructible.
  \item[(ii)] Suppose $f$ of finite presentation. The set of $y\in Y$ such that, if $s$ is the image of $y$ in $S$, $f_s$ is smooth (resp. unramified, resp. étale, resp. differentially smooth) at all the points of $f_s^{-1}(y)$, is locally constructible.
  \item[(iii)] Suppose $X$ and $Y$ of finite presentation over $S$. The set of $s\in S$ such that $f_s$ is smooth (resp. unramified, resp. étale, resp. differentially smooth) is locally constructible.
\end{enumerate}
\end{proposition}

Let $E$ be the set of $x\in X$ satisfying the property considered in (i). Then the set $F$ of $y\in Y$ satisfying the corresponding property considered in (ii) is

\oldpage[IV-4]{77}

none other than $Y-f(X-E)$, so (ii) follows from (i) and from Chevalley's theorem when $f$ is of finite presentation (1.8.4). Likewise, if $h:Y\to S$ is the structural morphism, the set of $s\in S$ satisfying the corresponding property considered in (iii) is $S-h(Y-F)$, so (iii) again follows from (ii) and from Chevalley's theorem when $f$ and $h$ are of finite presentation. It therefore suffices to prove the assertions of (i).

Let us first prove (i) for the property of being \emph{smooth}. The question being local on $X$, we may restrict to the case where $S=\Spec(A)$, $X=\Spec(B)$ and $Y=\Spec(C)$ are affine, $B$ and $C$ being $A$-algebras of finite presentation. Arguing as at the beginning of (9.9.1) and using (17.7.2, (ii)), we reduce to the case where $A$ is Noetherian. By virtue of ($0_{\mathrm{III}}$, 9.2.3), we are reduced to seeing that if $x\in E$ (resp. if $x\notin E$), there exists a neighbourhood $V$ of $x$ in $\overline{\{x\}}$ contained in $E$ (resp. in $X-E$). Denoting by $g:X\to S$ and $h:Y\to S$ the structural morphisms, we can first replace $S$ by the reduced sub-prescheme $S'$ of $S$ having $\overline{\{g(x)\}}$ as underlying space, $X$ by $g^{-1}(S')$, and $Y$ by $h^{-1}(S')$, the fibres of $X$ and $X'$ (resp. $Y$ and $Y'$) at the points of $S'$ being the same. In other words, we may restrict to the case where $S$ is \emph{integral} and $\eta=g(x)=h(y)$ (where $y=f(x)$) is its generic point.

$1^\circ$ Suppose first that $x\in E$. The local rings $\sh{O}_{X_\eta,x}$ and $\sh{O}_{Y_\eta,y}$ are respectively equal to $\sh{O}_{X,x}$ and $\sh{O}_{Y,y}$; as the property of smoothness of a morphism of finite presentation at a point depends only on the local ring of this point and on the local ring of its image (17.5.1), we see that the hypothesis $x\in E$ amounts to saying that the morphism $f$ is smooth at the point $x$; it then still has this property at the points of an open neighbourhood of $x$ in $X$, and it suffices to apply (17.3.3, (iii)) to obtain the conclusion.

$2^\circ$ Suppose secondly that $x\in X-E$, and that the morphism $f_\eta$ is not flat at the point $x$. The conclusion then follows from the following lemma, which sharpens (11.2.8):

\begin{lemma}[17.7.11.1]
\label{IV.17.7.11.1}
Let $g:X\to S$, $h:Y\to S$ be two morphisms locally of finite presentation, $f:X\to Y$ an $S$-morphism, and $\sh{F}$ a quasi-coherent $\sh{O}_X$-module of finite presentation. Then the set $E$ of $x\in X$ such that $\sh{F}_{g(x)}$ is $f_{g(x)}$-flat at the point $x$ is locally constructible.
\end{lemma}

Arguing again as at the beginning of (9.9.1) and using (2.5.1), we reduce to the case where $S$, $X$ and $Y$ are Noetherian; then we reduce as above to the case where $S$ is integral and $\eta=g(x)=h(f(x))$ is its generic point, and we must show that if $x\in E$ (resp. $x\notin E$) there is a neighbourhood $V$ of $x$ in $\overline{\{x\}}$ contained in $E$ (resp. in $X-E$). The case where $x\in E$ is an immediate consequence of (11.1.1). To deal with the case where $x\notin E$, we argue as in (9.4.7.1), whose notation we keep, so that we may suppose, replacing $Y$ and $X$ if necessary by neighbourhoods of $f(x)$ and $x$ respectively, that there exist two coherent $\sh{O}_Y$-modules $\sh{G}$, $\sh{H}$ and an $\sh{O}_Y$-homomorphism $u:\sh{G}\to\sh{H}$ such that for every $s\in S$, $u_s:\sh{G}_s\to\sh{H}_s$ is injective, but the homomorphism $1\otimes u_\eta:\sh{F}_\eta\otimes_{\sh{O}_{Y_\eta}}\sh{G}_\eta\to\sh{F}_\eta\otimes_{\sh{O}_{Y_\eta}}\sh{H}_\eta$ is not injective at the point $x$, in other words $x\in\operatorname{Supp}(\operatorname{Ker}(1\otimes u_\eta))$; if we set $T=\overline{\{x\}}$, we therefore have $\operatorname{Supp}(\operatorname{Ker}(1\otimes u_\eta))\supset T_\eta$. But we may suppose that for every $s\in S$, we have $\operatorname{Ker}(1\otimes u_s)=(\operatorname{Ker}(1\otimes u))_s$ (9.4.2), hence

\oldpage[IV-4]{78}

$\operatorname{Supp}(\operatorname{Ker}(1\otimes u_s))=(\operatorname{Supp}(\operatorname{Ker}(1\otimes u)))_s$ (I, 9.1.13.1); finally, it follows from (9.5.2) that for $s$ in a neighbourhood of $\eta$, we have $(\operatorname{Supp}(\operatorname{Ker}(1\otimes u)))_s\supset T_s$, which establishes the lemma.

$3^\circ$ Suppose now that $x\in X-E$, that the morphism $f_\eta$ is flat at the point $x$, but that $f_\eta$ is not smooth at the point $x$. Note that to say that $f_\eta$ is flat at the point $x$ is equivalent to saying that $f$ itself is flat at the point $x$, and by replacing $X$ by a neighbourhood of $x$, we may suppose that $f$ is flat (11.1.1); we conclude that the same holds for $f_s$ for every $s\in S$, and as for every $y\in Y$, $f^{-1}(y)=f_{h(y)}^{-1}(y)$, it amounts to the same to say that $f_{g(x')}$ is smooth at the point $x'$ or to say that $f$ is smooth at the point $x'$. But the set of $x'\in X$ where $f$ is smooth is \emph{open} in $X$ (12.1.7), so the set of $x'\in X$ where $f$ is not smooth is closed, and since it contains $x$ by hypothesis, it also contains $\overline{\{x\}}$, which completes the proof for the first property considered in (i).

Secondly, let us prove (i) for the property of being \emph{étale}. For this, note that this property for $f_s$ at the point $x$ is equivalent to saying that $f_s$ is both smooth and quasi-finite at the point $x$ (17.6.1). Now, it amounts to the same to say that $f_{g(x)}$ is quasi-finite at the point $x$ or that $f$ itself is quasi-finite at this point; it therefore follows from (13.1.4) that the set of points $x$ such that $f_{g(x)}$ is quasi-finite at the point $x$ is open in $X$, and \emph{a fortiori} locally constructible; the conclusion therefore follows from the fact that the set of $x$ such that $f_{g(x)}$ is smooth at the point $x$ is also locally constructible.

Let us pass to the proof of (i) for the property of being \emph{differentially smooth}. Let $p:X\times_Y X\to X$ be the second canonical projection; the second projection $X_s\times_{Y_s}X_s\to X_s$ is none other than $p_s$ for every $s\in S$, and it therefore follows from (17.12.5.1)\footnote{The reader will verify that the result of (17.7.11) is not used in the remainder of § 17, and that there is therefore no vicious circle.} that for $f_{g(x)}$ to be differentially smooth at the point $x$, it is necessary and sufficient that $p_{g(x)}$ be smooth at the point $\Delta_f(x)$. As $p$ is locally of finite presentation, the set of points $z\in X\times_Y X$ such that $p_{g(p(z))}$ is smooth at the point $z$ is locally constructible, and the intersection of this set with the locally closed set $\Delta_f(X)$ is therefore also locally constructible in $\Delta_f(X)$ (1.8.2); the restriction of $p$ to $\Delta_f(X)$ being an isomorphism onto $X$, we see that the set of $x\in X$ such that $f_{g(x)}$ is differentially smooth at the point $x$ is locally constructible.

Finally, consider the property of being unramified; note that the diagonal morphism $\Delta_f:X\to X\times_Y X$ is an immersion locally of finite presentation (1.4.3.1) and that, for every $s\in S$, the diagonal morphism $\Delta_{f_s}:X_s\to X_s\times_{Y_s}X_s$ is none other than $(\Delta_f)_s$; to say that $f_{g(x)}$ is unramified at the point $x$ is equivalent to saying that $(\Delta_f)_{g(x)}$ is a local isomorphism at the point $x$ (17.4.1), and since it is an immersion locally of finite presentation, it amounts to the same (17.9.1)\footnote{The reader will verify that the result of (17.7.11) is not used in the rest of § 17, and that there is therefore no vicious circle.} to say that $(\Delta_f)_{g(x)}$ is étale at the point $x$; it therefore suffices to apply what was seen above for the property of being étale.

\oldpage[IV-4]{79}


\subsection{Criteria for smoothness and unramifiedness by fibres}
\EGAReaderTarget{ega4.c04.subsec.17-8}{subsection:IV.17.8}
\label{IV.17.8}

\begin{proposition}[17.8.1]
\label{IV.17.8.1}
Let $g:Y\to S$, $h:X\to S$ be two morphisms locally of finite presentation. For an $S$-morphism $f:X\to Y$ to be unramified, it is necessary and sufficient that for every $s\in S$, the morphism $f_s:h^{-1}(s)\to g^{-1}(s)$ deduced from $f$ by the base change $\Spec(k(s))\to S$ be unramified.
\end{proposition}

This follows immediately from the fact that, for a morphism locally of finite presentation, being unramified is a property of the \emph{fibres} of this morphism (17.4.1, $d$)), and from the fact that, for every $y\in Y$, we have $f^{-1}(y)=f_s^{-1}(y)$ if $s=g(y)$.

\begin{proposition}[17.8.2]
\label{IV.17.8.2}
Let $g:Y\to S$, $h:X\to S$ be two morphisms locally of finite presentation; suppose moreover that $h$ is flat. For an $S$-morphism $f:X\to Y$ to be smooth (resp. étale), it is necessary and sufficient that, for every $s\in S$, the morphism $f_s:h^{-1}(s)\to g^{-1}(s)$ deduced from $f$ by the base change $\Spec(k(s))\to S$ be smooth (resp. étale). When this is the case, the morphism $g$ is flat at the points of $f(X)$.
\end{proposition}

Indeed, we know (11.3.10) that for $f$ to be flat, it is necessary and sufficient that $f_s$ be flat for every $s\in S$, and that $g$ is then flat at the points of $f(X)$. But for a flat morphism locally of finite presentation, being smooth is a property of the fibres of this morphism (17.5.1, $b$)), and for every $y\in Y$, we have $f^{-1}(y)=f_s^{-1}(y)$ if $s=g(y)$.

\begin{remark}[17.8.3]
\label{IV.17.8.3}
The preceding proofs show (taking into account (11.3.10)) that if the hypotheses on $g$ and $h$ are the same as above, then, for $f$ to be unramified (resp. smooth, resp. étale) \emph{at a point} $x\in X$, it suffices that, if we set $s=h(x)$, $f_s$ be unramified (resp. smooth, resp. étale) at the point $x$.
\end{remark}

\subsection{Étale morphisms and open immersions}
\EGAReaderTarget{ega4.c04.subsec.17-9}{subsection:IV.17.9}
\label{IV.17.9}

\begin{theorem}[17.9.1]
\label{IV.17.9.1}
Let $f:X\to Y$ be a morphism. The following properties are equivalent:
\begin{enumerate}
  \item[a)] $f$ is an open immersion.
  \item[b)] $f$ is a flat monomorphism locally of finite presentation.
  \item[c)] $f$ is étale and radicial.
\end{enumerate}
\end{theorem}

It follows from (1.4.3, (i)) that a) implies b). Condition b) implies that for every $y\in Y$, the fibre $f^{-1}(y)$ is empty or isomorphic to $\Spec(k(y))$ (8.11.5.1), hence b) implies c), by virtue of (17.6.2, $c$)). It remains to see that c) implies a).

The question being local on $X$ and on $Y$ (since $f$ is injective), we may restrict to the case where $Y$ is affine and $f$ of finite presentation. Since $f$ is flat, it is an \emph{open} morphism (2.4.6), so, replacing $Y$ by $f(X)$, we may suppose that $f$ is \emph{surjective}. For every morphism $Y'\to Y$, $f'=f_{(Y')}:X_{(Y')}\to Y'$ is again étale, radicial, surjective and of finite presentation, hence open, and consequently a homeomorphism; in other words $f$ is a \emph{universal homeomorphism}, and, being of finite type and separated (1.8.7.1), $f$ is \emph{proper}. The hypothesis that $f$ is radicial and of finite type then implies that $f$ is quasi-finite; hence (8.11.1) $f$ is a \emph{finite} morphism. To prove that $f$ is an isomorphism, we may restrict to the case

\oldpage[IV-4]{80}

where $Y=\Spec(A)$, $A$ being a local ring. Since $f$ is of finite presentation, we have $X=\Spec(B)$, where $B$ is a flat $A$-module of finite presentation (1.4.7), hence \emph{free} (Bourbaki, \emph{Alg. comm.}, chap. II, § 5, n\textsuperscript{o} 2, cor. 2 of th. 1). Moreover, if $\mathfrak m$ is the maximal ideal of $A$ and $k$ its residue field, $B/\mathfrak mB$ is by hypothesis a field, at the same time a radicial extension and a finite separable extension of $k$, since $f$ is étale and radicial (17.6.1); hence $B/\mathfrak mB$ is isomorphic to $k$, and since $B$ is a free $A$-module, $B$ is isomorphic to $A$. Q.E.D.

\begin{corollary}[17.9.2]
\label{IV.17.9.2}
Let $X$ be a connected prescheme; if $f:X\to Y$ is an étale closed immersion, then $f$ is an isomorphism of $X$ onto an open connected component of $Y$.
\end{corollary}

Indeed $f$ is étale and radicial, hence an isomorphism of $X$ onto a prescheme induced on an open subset of $Y$; but by hypothesis $f(X)$ is closed in $Y$, hence both open and closed, and since $f(X)$ is connected, it is a connected component of $Y$.

\begin{corollary}[17.9.3]
\label{IV.17.9.3}
Let $f:X\to Y$ be an étale (resp. étale and separated) morphism. Then every $Y$-section $g:Y\to X$ of $X$ is an open (resp. open and closed) immersion; moreover the map $g\mapsto g(Y)$ is a bijection of the set $\Gamma(X/Y)$ of $Y$-sections of $X$ onto the set of open (resp. open and closed) subsets $Z$ of $X$ such that the restriction of $f$ to $Z$ is a surjective and radicial morphism of $Z$ onto $Y$.
\end{corollary}

Indeed, the fact that $g$ is an open immersion already follows from the fact that $f$ is unramified (17.4.1, $b''$)), and the restriction of $f$ to the open subset $g(Y)$ of $X$ is an isomorphism. Conversely, if $Z$ is an open subset of $X$ such that $f|Z$ is a surjective and radicial morphism of $Z$ onto $Y$, $f|Z$ is an isomorphism by virtue of (17.9.1), since it is étale. If moreover $f$ is separated, we know that $g$ is a closed immersion (1, 5.4.6), which completes the proof of the corollary.

\begin{corollary}[17.9.4]
\label{IV.17.9.4}
Let $Y$ be a connected prescheme, $f:X\to Y$ an étale and separated morphism. Then every $Y$-section $g$ of $X$ is an isomorphism of $Y$ onto an open connected component of $X$, and the map $g\mapsto g(Y)$ is a bijection of $\Gamma(X/Y)$ onto the set of open connected components $Z$ of $X$ such that the restriction of $f$ to $Z$ is a surjective and radicial morphism of $Z$ onto $Y$.
\end{corollary}

\begin{corollary}[17.9.5]
\label{IV.17.9.5}
Let $g:Y\to S$, $h:X\to S$ be two morphisms locally of finite presentation; suppose moreover that $h$ is flat. For an $S$-morphism $f:X\to Y$ to be an open immersion (resp. an isomorphism), it is necessary and sufficient that, for every $s\in S$, the morphism $f_s:h^{-1}(s)\to g^{-1}(s)$ deduced from $f$ by the base change $\Spec(k(s))\to S$ be an open immersion (resp. an isomorphism).
\end{corollary}

Indeed, if $f_s$ is an open immersion for every $s\in S$, it follows from (17.8.3) that $f$ is an étale morphism; since for every $y\in Y$ we have $f^{-1}(y)=f_s^{-1}(y)$ with $s=g(y)$, $f$ is radicial; hence $f$ is an open immersion by virtue of (17.9.1). If moreover $f_s$ is surjective for every $s\in S$, $f$ is surjective, hence an isomorphism.

The following proposition makes (10.4.11) more precise:

\begin{proposition}[17.9.6]\label{IV.17.9.6}
\EGAReaderTarget{ega4.c04.prop.17-9-6}{IV.17.9.6}
Let $S$ be a prescheme, $X$ an $S$-prescheme of finite presentation. Every $S$-endomorphism of $X$ which is a monomorphism is an automorphism of $X$.
\end{proposition}

Let $f:X\to S$ be the structural morphism and $g$ the $S$-endomorphism under consideration. The question is local on $S$, and we may consequently suppose that $S$ is affine. Using (8.9.1) and (8.10.5, (i \emph{bis})), we are reduced to the case where $S=\Spec(A)$, where $A$ is a $\mathbb{Z}$-algebra of finite type, and consequently $X$ is a $\mathbb{Z}$-prescheme of finite type. It already follows from (10.4.11) that $g$ is a bijective morphism, since a monomorphism is radicial (8.11.5.1); it will therefore suffice to show

\oldpage[IV-4]{81}

that $g$ is an open immersion, and since $g$ is radicial, it will suffice, by virtue of (17.9.1), to prove that $g$ is étale. Moreover, since the set of points where $g$ is étale is open and $X$ is a Jacobson prescheme (10.4.7), it will suffice to show that $g$ is étale at every \emph{closed} point $z$ of $X$ (10.3.1). Set $z'=g(z)$; it follows from (10.4.11.1, (i)) that $z'$ is also a closed point of $X$, and since $g$ is a monomorphism, the map $k(z')\to k(z)$ deduced from $g$ is an \emph{isomorphism}. For $g$ to be étale at the point $z$, it is therefore necessary and sufficient that the canonical homomorphism $\widehat{\sh{O}}_{X,z'}\to\widehat{\sh{O}}_{X,z}$ be bijective (17.6.3, $e''$)). For this we will prove that for every integer $n$, the canonical homomorphism $\sh{O}_{X,z'}/\mathfrak m_{z'}^{n+1}\to\sh{O}_{X,z}/\mathfrak m_z^{n+1}$ is bijective, from which the conclusion will follow immediately. We may suppose that $z$ belongs to the finite set $T_{p,d}$ of closed points $t$ of $X$ such that $k(t)$ is an extension of $\mathbb{F}_p$ whose degree divides $d$; in which case the same holds for $z'$. Denote by $\sh{J}$ the coherent ideal of $\sh{O}_X$ such that $\sh{J}_x=\sh{O}_x$ for $x\notin T_{p,d}$, and $\sh{J}_x=\mathfrak m_x^{n+1}$ for $x\in T_{p,d}$. Let $T_{p,d,n}$ be the closed sub-prescheme of $X$ defined by $\sh{J}$, and $j:T_{p,d,n}\to X$ the canonical injection; the composite morphism $g\circ j:T_{p,d,n}\to X$ maps the set $T_{p,d}$ into itself, and for every $x\in T_{p,d}$, the homomorphism $\sh{O}_{X,g(x)}\to\sh{O}_{X,x}/\mathfrak m_x^{n+1}$ factors as
\[
\sh{O}_{X,g(x)}\to\sh{O}_{X,g(x)}/\mathfrak m_{g(x)}^{n+1}\to\sh{O}_{X,x}/\mathfrak m_x^{n+1}.
\]
Hence (I, 4.9) there exists a unique endomorphism $g'$ of $T_{p,d,n}$ such that $j\circ g'=g\circ j$. Since $j$ is a monomorphism, and the same holds for $g$ by hypothesis, we deduce that $g'$ is also a monomorphism, and the proposition will be proved if we show that $g'$ is an \emph{automorphism} of $T_{p,d,n}$. Now, for every $x\in T_{p,d,n}$, we have seen that $k(x)$ is a finite field, and since $\sh{O}_{X,x}$ is Noetherian, each of the $\mathfrak m_x^h/\mathfrak m_x^{h+1}$ is a $k(x)$-vector space of finite rank; from which we conclude immediately that the local ring $\sh{O}_{X,x}/\mathfrak m_x^{n+1}$ of $T_{p,d,n}$ at the point $x$ has a \emph{finite} number of elements, and \emph{a fortiori} is a $\mathbb{Z}$-module of finite type. Since $T_{p,d,n}$ is the sum of finitely many preschemes $\Spec(\sh{O}_{X,x}/\mathfrak m_x^{n+1})$, it is a \emph{finite} $\mathbb{Z}$-prescheme; the endomorphism $g'$ is therefore also \emph{finite} (II, 6.1.5, (v)), hence proper. But a proper monomorphism is a \emph{closed immersion} (8.11.5); if $T_{p,d,n}=\Spec(B)$, $g'$ therefore corresponds to a surjective endomorphism $\varphi:B\to B$ of the ring $B$. Now the set $B$ is finite, hence $\varphi$ is necessarily bijective. Q.E.D.

\subsection{Relative dimension of a prescheme smooth over another}
\EGAReaderTarget{ega4.c04.subsec.17-10}{subsection:IV.17.10}

\begin{definition}[17.10.1]\label{IV.17.10.1}
\EGAReaderTarget{ega4.c04.def.17-10-1}{IV.17.10.1}
Let $f:X\to Y$ be a morphism locally of finite type. We call \emph{relative dimension of $f$ at the point $x\in X$} (or \emph{relative dimension of $X$ over $Y$ at the point $x$}), and denote by $\dim_x f$, the nonnegative integer $\dim_x(f^{-1}(f(x)))$.
\end{definition}

To say that $f$ is \emph{quasi-finite} at the point $x$ (Err$_{\mathrm{III}}$, 20) is therefore equivalent to saying that $\dim_x f=0$. We have seen (13.1.3) that the function $x\mapsto\dim_x f$ is upper semi-continuous. Note that, even when the morphism $f$ has property $(S_1)$ (in other words (6.8.1) is flat and such that its fibres have no embedded associated prime cycle), the function $x\mapsto\dim_x f$ is not necessarily continuous, as is shown by the example where $Y=\Spec(k)$, with $k$ a field, and $X=\Spec(k[U,V,W]/\mathfrak p\mathfrak q)$, where $\mathfrak p=(W)$ and $\mathfrak q=(U)+(V-W)$, prime ideals of $k[U,V,W]$ ($X$ thus being the union, in 3-dimensional space, of a plane and a line not parallel to this plane).

To say that a morphism $f:X\to Y$ locally of finite presentation is étale at the point $x\in X$ also means that $f$ is smooth at the point $x$ and that $\dim_x f=0$ (17.6.1).

\begin{proposition}[17.10.2]
\label{IV.17.10.2}
Let $f:X\to Y$ be a smooth morphism. For every $x\in X$, the locally free $\sh{O}_X$-module $\Omega_f^1$ (17.2.3) has rank at $x$ equal to $\dim_x f$ (which implies that $x\mapsto\dim_x f$ is a continuous function on $X$).
\end{proposition}

Indeed, if $y=f(x)$, $X_y=f^{-1}(y)$ is smooth over $k(y)$, and if $f_y:X_y\to\Spec(k(y))$ is the structural morphism, we have $\Omega_{f_y}^1=\Omega_f^1\otimes_{\sh{O}_X}\sh{O}_{X_y}$ (16.4.5); we may therefore restrict to the case where $Y=\Spec(k)$ is the spectrum of a field; moreover, by virtue of (16.4.5) and (17.7.1), we may replace $k$ by an algebraically closed extension, in other words

\oldpage[IV-4]{82}

suppose $k$ algebraically closed. The set of points of $X$ rational over $k$ then being dense in $X$ (10.4.8), we may restrict to the case where $x$ is rational over $k$. Now (16.4.12) $(\Omega^1_{X/k})_x\otimes_{\sh{O}_{X,x}}k(x)$ is then $k$-isomorphic to $\mathfrak m_x/\mathfrak m_x^2$, and since $\sh{O}_x$ is a regular local ring (17.5.1), $\operatorname{rg}_{k(x)}(\mathfrak m_x/\mathfrak m_x^2)=\dim(\sh{O}_x)$ (0, 17.1.1); hence $\operatorname{rg}_{\sh{O}_x}(\Omega^1_{X/k})_x=\dim(\sh{O}_x)$. But since $k(x)=k$, we have $\dim_x f=\dim(\sh{O}_x)$ (5.2.3). Q.E.D.

We will prove a converse of this result further on (17.15.5).

\begin{corollary}[17.10.3]
\label{IV.17.10.3}
Let $f:X\to Y$, $g:Y\to Z$ be two smooth morphisms. Then, for every $x\in X$, we have
\begin{equation}
\label{IV.17.10.3.1}
\tag{17.10.3.1}
\dim_x(g\circ f)=\dim_x f+\dim_{f(x)}g.
\end{equation}
\end{corollary}

Indeed, $g\circ f$ is smooth (17.3.3), hence the three $\sh{O}_X$-modules $\Omega^1_{X/Y}$, $\Omega^1_{X/Z}$ and $f^*(\Omega^1_{Y/Z})$ are locally free (17.2.3 and $0_{\mathrm I}$, 5.4.5); moreover the rank at $x$ of $f^*(\Omega^1_{Y/Z})$ is equal to the rank at $y=f(x)$ of $\Omega^1_{Y/Z}$. Equality (17.10.3.1) is therefore a consequence of (17.10.2) and of the exactness of the sequence (17.2.3.1).

\begin{corollary}[17.10.4]
\label{IV.17.10.4}
Let $f:X\to Y$ be a smooth morphism, $X'$ a sub-prescheme of $X$ such that the composite morphism $X'\xrightarrow{j}X\xrightarrow{f}Y$ (where $j$ is the canonical injection) is smooth. Then the conormal sheaf $\sh{N}_{X'/X}$ is a locally free $\sh{O}_{X'}$-module, and for every $x\in X'$, we have
\begin{equation}
\label{IV.17.10.4.1}
\tag{17.10.4.1}
\dim_x f=\dim_x(f\circ j)+\operatorname{rg}_{\sh{O}_{X',x}}(\sh{N}_{X'/X})_x.
\end{equation}
\end{corollary}

Indeed, $\Omega^1_{X/Y}\otimes_{\sh{O}_X}\sh{O}_{X'}$ and $\Omega^1_{X'/Y}$ are then locally free and the exact sequence (17.2.5.1) is split in a suitable neighbourhood of each point of $X'$, hence $\sh{N}_{X'/X}$ is locally free (Bourbaki, \emph{Alg. comm.}, chap. II, § 5, n\textsuperscript{o} 2, th. 1), and relation (17.10.4.1) follows immediately from the exactness of the sequence (17.2.5.1).

\subsection{Smooth morphisms of smooth preschemes}
\EGAReaderTarget{ega4.c04.subsec.17-11}{subsection:IV.17.11}

\begin{theorem}[17.11.1]
\label{IV.17.11.1}
Let $f:Y\to S$ and $h:X\to S$ be two morphisms locally of finite presentation, $g:X\to Y$ an $S$-morphism, $x$ a point of $X$; set $y=g(x)$, $s=f(y)=h(x)$. The following conditions are equivalent:
\begin{enumerate}
  \item[a)] $f$ is smooth at the point $y$ and $g$ is smooth at the point $x$.
  \item[b)] $g$ and $h$ are smooth at the point $x$.
  \item[c)] $h$ is smooth at the point $x$, and the canonical homomorphism (16.4.18)
  \begin{equation}
  \label{IV.17.11.1.1}
  \tag{17.11.1.1}
  (g^*(\Omega^1_{Y/S}))_x\to(\Omega^1_{X/S})_x
  \end{equation}
  is left invertible (in other words, is an isomorphism onto a \emph{direct factor} of $(\Omega^1_{X/S})_x$).
  \item[c')] $h$ is smooth at the point $x$, and the canonical homomorphism
  \begin{equation}
  \label{IV.17.11.1.2}
  \tag{17.11.1.2}
  ((\Omega^1_{Y/S})_y\otimes_{\sh{O}_{Y,y}}k(y))\otimes_{k(y)}k(x)\to(\Omega^1_{X/S})_x\otimes_{\sh{O}_{X,x}}k(x)
  \end{equation}
  is injective.
\end{enumerate}
Suppose moreover that the homomorphism $k(y)\to k(x)$ is bijective. Then the preceding conditions are also equivalent to the following:
\begin{enumerate}
  \item[d)] $h$ is smooth at the point $x$, and the canonical map $\mathrm{T}_{X/S}(x)\to\mathrm{T}_{Y/S}(y)$ from the tangent vector space to $X$ at $x$ to the tangent vector space to $Y$ at $y$ (16.5.12) is surjective.
\end{enumerate}
\end{theorem}

That a) implies b) follows trivially from (17.3.3, (iii)); b) implies c) by application of (17.2.3, (ii)). To see that c) is equivalent to c'), note that by virtue of (17.2.3, (i)), $\Omega^1_{X/S}$ is a locally free $\sh{O}_X$-module of finite type; it then suffices to apply (0, 19.1.12) to the local ring $\sh{O}_{X,x}$ and to the homomorphism (17.11.1.1) of $\sh{O}_{X,x}$-modules of finite type. When $k(x)=k(y)$, the tangent linear map $\mathrm{T}_{X/S}(x)\to\mathrm{T}_{Y/S}(y)$ is the transpose of (17.11.1.2), by (16.5.12), which establishes in this case the equivalence of c') and d).

\oldpage[IV-4]{83}

It therefore remains to prove that c) implies a). We may restrict to the case where $S=\Spec(A)$, $Y=\Spec(B)$, $X=\Spec(C)$ are affine. The hypothesis implies (17.2.3, (i)) that $(\Omega^1_{X/S})_x=\Omega^1_{C_{\mathfrak m}/A_{\mathfrak s}}$ is a free $C_{\mathfrak m}$-module: indeed (16.10.6) there exist elements $t_i$ ($1\leq i\leq r$) of $B_{\mathfrak p}=\sh{O}_{Y,y}$ such that the differentials $d_{B_{\mathfrak p}/A_{\mathfrak s}}(t_i)$ generate the $B_{\mathfrak p}$-module $\Omega^1_{B_{\mathfrak p}/A_{\mathfrak s}}$ and their images in $\Omega^1_{C_{\mathfrak m}/A_{\mathfrak s}}$ form part of a basis of this free $C_{\mathfrak m}$-module. Since $\Omega^1_{Y/S}$ (resp. $\Omega^1_{X/S}$) is an $\sh{O}_Y$-module (resp. an $\sh{O}_X$-module) of finite presentation (16.4.22), we may, replacing $X$ and $Y$ if necessary by suitable affine open neighbourhoods of $x$ and $y$ respectively, suppose that $\Omega^1_{C/A}$ is a free $C$-module and that the $t_i$ are the images of elements $s_i$ ($1\leq i\leq r$) of $B$ such that the $d_{B/A}(s_i)$ generate the $B$-module $\Omega^1_{B/A}$ and their images in $\Omega^1_{C/A}$ form part of a basis of this $C$-module (Bourbaki, \emph{Alg. comm.}, chap. II, § 5, n\textsuperscript{o} 1, prop. 2). Let $\varphi$ be the $A$-homomorphism from $B'=A[T_1,\ldots,T_r]$ to $B$ such that $\varphi(T_i)=s_i$ for every $i$; the corresponding di-homomorphism $\Omega^1_{B'/A}\to\Omega^1_{B/A}$ (0, 20.5.2) transforms the $d_{B'/A}(T_i)$, which form a basis of $\Omega^1_{B'/A}$ (0, 20.4.13), into the $d_{B/A}(s_i)$, and is consequently \emph{surjective}; if $Y'=\Spec(B')$, and if $u:Y\to Y'$ is the $S$-morphism corresponding to $\varphi$, we conclude from (17.2.2) that $u$ is \emph{unramified}. If we prove that the composite morphism $u\circ g:X\to Y'$ is smooth at the point $x$, it will therefore follow from (17.7.10) that $u$ is étale at the point $y$, and then from (17.3.5) that $g$ is smooth at the point $x$; finally, since the structural morphism $f':Y'\to S$ is smooth (17.3.8), $f=f'\circ u$ will be smooth at the point $y$. Since the canonical images of the $d_{B'/A}(T_i)$ in $\Omega^1_{C/A}$ are those of the $d_{B/A}(s_i)$, they form part of a basis of $\Omega^1_{C/A}$. We thus see that to prove that c) implies a), we may restrict to the case where $Y'=Y$, and consequently suppose that $f$ is a \emph{smooth} morphism.

By virtue of (17.8.2), we may then restrict to the case where $S=\Spec(k)$ is the spectrum of a field; furthermore, thanks to (17.7.1, (ii)), we may suppose that $k$ is algebraically closed; finally, replacing $X$ and $Y$ if necessary by open neighbourhoods of $x$ and $y$ respectively, we may suppose that $f$ and $h$ are both smooth, and that the canonical homomorphism
\[
g^*(\Omega^1_{Y/S})\to\Omega^1_{X/S}
\]
is left invertible (0, 19.1.12). Then the set of points of $X$ rational over $k$ is very dense in $X$ (10.4.8), hence, to prove that $g$ is smooth, it suffices to prove that $g$ is smooth at every point $x\in X$ rational over $k$, or equivalently that at such a point, $\sh{O}_{X,x}$ is a flat $\sh{O}_{Y,y}$-module and $\sh{O}_{X,x}/\mathfrak m_y\sh{O}_{X,x}$ a regular ring (17.5.1). Now $y=g(x)$ is

\oldpage[IV-4]{84}

\emph{a fortiori} rational over $k$, and since $\sh{O}_{Y,y}$ is then a formally smooth $k$-algebra for the discrete topologies and $\sh{O}_{Y,y}/\mathfrak m_y=k$, it follows from (0, 20.5.14) that the canonical homomorphism $\mathfrak m_y/\mathfrak m_y^2\to(\Omega^1_{Y/S})_y\otimes_{\sh{O}_{Y,y}}k(y)$ is bijective; likewise, the canonical homomorphism $\mathfrak m_x/\mathfrak m_x^2\to(\Omega^1_{X/S})_x\otimes_{\sh{O}_{X,x}}k(x)$ is bijective. Hypothesis c'), which is equivalent to c), therefore means here that the canonical homomorphism
\[
(\mathfrak m_y/\mathfrak m_y^2)\otimes_{k(y)}k(x)\to\mathfrak m_x/\mathfrak m_x^2
\]
is injective. Since the ring $\sh{O}_{X,x}$ is regular, the conclusion follows from (0, 17.3.3). Q.E.D.

\begin{corollary}[17.11.2]
\label{IV.17.11.2}
Under the general hypotheses of (17.11.1), the following conditions are equivalent:
\begin{enumerate}
  \item[a)] $f$ is smooth at the point $y$ and $g$ is étale at the point $x$.
  \item[b)] $h$ is smooth at the point $x$ and $g$ is étale at the point $x$.
  \item[c)] $h$ is smooth at the point $x$, and the canonical homomorphism
  \[
  (g^*(\Omega^1_{Y/S}))_x\to(\Omega^1_{X/S})_x
  \]
  is bijective.
  \item[c')] $h$ is smooth at the point $x$, and the canonical homomorphism
  \[
  ((\Omega^1_{Y/S})_y\otimes_{\sh{O}_{Y,y}}k(y))\otimes_{k(y)}k(x)\to(\Omega^1_{X/S})_x\otimes_{\sh{O}_{X,x}}k(x)
  \]
  is bijective.
\end{enumerate}
Suppose moreover that the homomorphism $k(y)\to k(x)$ is bijective. Then the preceding conditions are also equivalent to the following:
\begin{enumerate}
  \item[d)] $h$ is smooth at the point $x$, and the canonical map $\mathrm{T}_{X/S}(x)\to\mathrm{T}_{Y/S}(y)$ (16.5.12) is bijective.
\end{enumerate}
\end{corollary}

Each of the conditions a), b), c) of (17.11.2) is equivalent to the conjunction of the corresponding condition of (17.11.1) and of the fact that $g$ is unramified at the point $x$, taking into account (17.2.2) as regards condition c); whence the equivalence of a), b) and c). The equivalence of c) and c') follows from the equivalence of the corresponding conditions of (17.11.1) and from Nakayama's lemma; the equivalence of c') and d) when $k(x)=k(y)$ is immediate by transposition (16.5.12).

\begin{corollary}[17.11.3]
\label{IV.17.11.3}
Let $h:X\to S$ be a smooth morphism, $s_i$ ($1\leq i\leq r$) sections of $\sh{O}_X$ over $X$ (which are also sections of $h_*(\sh{O}_X)$ over $S$), and $g:X\to S[T_1,\ldots,T_r]=\mathbf{V}_S^r$ the $S$-morphism corresponding to the homomorphism of $\sh{O}_S$-modules $\sh{O}_S^r\to h_*(\sh{O}_X)$ defined by these sections (II, 1.2.7). For $g$ to be smooth (resp. étale) at a point $x\in X$, it is necessary and sufficient that the $(d_{X/S}(s_i))_x$ form part of a basis (resp. constitute a basis) of the $\sh{O}_{X,x}$-module $(\Omega^1_{X/S})_x$.
\end{corollary}

It suffices to apply (17.11.1) (resp. (17.11.2)), taking for $f$ the structural morphism $S[T_1,\ldots,T_r]\to S$.

\begin{corollary}[17.11.4]
\label{IV.17.11.4}
For a morphism $h:X\to S$ to be smooth at a point $x\in X$, it is necessary and sufficient that there exist an open neighbourhood $U$ of $x$, an integer $r$, and an étale $S$-morphism $U\to S[T_1,\ldots,T_r]$.
\end{corollary}

\oldpage[IV-4]{85}

The condition is evidently sufficient, since the structural morphism $S[T_1,\ldots,T_r]\to S$ is smooth (17.3.8). To show that it is necessary, note that since $h$ is smooth at the point $x$, there exists an open neighbourhood $U$ of $x$ such that $\Omega^1_{X/S}|U$ is locally free (17.2.3); it then suffices to use (16.10.6) to obtain (shrinking $U$ if necessary) sections $s_i$ of $\sh{O}_X$ over $U$ such that the $d_{X/S}(s_i)$ form a basis of $\Omega^1_{X/S}|U$; we conclude by applying (17.11.3).

\begin{proposition}[17.11.5]
\label{IV.17.11.5}
Let $f:Y\to S$, $h:X\to S$ be two smooth morphisms. For an $S$-morphism $g:X\to Y$ to be an open immersion, it is necessary and sufficient that $g$ be a monomorphism of preschemes and that for every $x\in X$, if we set $y=g(x)$, we have $\dim_y(f)=\dim_x(h)$ (for a generalization of this proposition, see (18.10.5)).
\end{proposition}

The condition is evidently necessary; let us prove that it is sufficient.

By virtue of (17.9.5), we are immediately reduced to the case where $S=\Spec(k)$ is the spectrum of a field, and by virtue of (2.7.1, (x)), we may suppose $k$ algebraically closed. Taking into account (17.9.1), it then suffices to prove that $g$ is étale, and since the set of points where $g$ is étale is open, it suffices to show that $g$ is étale at the \emph{closed} points (or equivalently, the points rational over $k$) of $X$ (10.4.8). So let $x$ be such a point, and set $y=g(x)$, which is also rational over $k$; by hypothesis, the ring $A=\sh{O}_{Y,y}$ is regular with residue field $k$; let $d=\dim(A)$ and let $(t_i)_{1\leq i\leq d}$ be a regular system of parameters for $A$. Set $B=\sh{O}_{X,x}$, $C=B/\mathfrak m_yB$. Since $g$ is a monomorphism, so is the morphism $\Spec(C)\to\Spec(k(y))=\Spec(k)$ deduced from $g$ by base change (I, 3.3.12); but this means that the corresponding homomorphism $u:k\to C$ is \emph{surjective} (hence bijective), because $u$ admits a left inverse $v:C\to k$, and $u\circ v$ and the identity of $C$, composed with $u$, give the same morphism $u:k\to C$. Since by hypothesis $B$ is a regular ring of dimension $d$, the images of the $t_i$ in $B$ form a regular system of parameters for $B$ (0, 17.1.7); condition (17.6.3, $e''$) is therefore satisfied by $g$ at the point $x$ (0, 17.1.1 and Bourbaki, \emph{Alg. comm.}, chap. III, § 2, n\textsuperscript{o} 8, cor. 3 of th. 1), which completes the proof.

\subsection{Smooth sub-preschemes of a smooth prescheme. Smooth morphisms and differentially smooth morphisms}
\EGAReaderTarget{ega4.c04.subsec.17-12}{subsection:IV.17.12}

\begin{theorem}[17.12.1]
\label{IV.17.12.1}
Let $f:X\to S$, $h:Y\to S$ be two morphisms locally of finite presentation, $j:Y\to X$ an immersion, $y$ a point of $Y$, $x=j(y)$. The following conditions are equivalent:
\begin{enumerate}
  \item[a)] $h$ is smooth at the point $y$ and $f$ is smooth at the point $x$.
  \item[b)] $f$ is smooth at the point $x$, and the canonical homomorphism (16.4.21)
  \[
  (\sh{N}_{Y/X})_y\to(j^*(\Omega^1_{X/S}))_y
  \]
  is left invertible.
  \item[c)] $h$ is smooth at the point $y$ and there exists an open neighbourhood $U$ of $y$ in $Y$ such that $j|U:U\to X$ is a regular immersion (16.9.2).

\oldpage[IV-4]{86}

  \item[c')] $h$ is smooth at the point $y$ and there exists an open neighbourhood $U$ of $y$ in $Y$ such that $j|U:U\to X$ is a quasi-regular immersion (in other words (16.9.8), there exists a neighbourhood $U$ of $y$ in $Y$ in which $\sh{N}_{Y/X}$ is a locally free $\sh{O}_Y$-module and the canonical homomorphism
  \[
  \mathbf{S}_{\sh{O}_Y}(\sh{N}_{Y/X})\to\operatorname{gr}^{\bullet}(j)
  \]
  is bijective).
\end{enumerate}
\end{theorem}

To prove the equivalence of a) and b), we may restrict to the case where $S=\Spec(A)$ and $X=\Spec(B)$ are affine, with $Y=\Spec(B/\mathfrak J)$, where $\mathfrak J$ is an ideal of finite type of $B$, and $B$ is a smooth $A$-algebra (17.3.2, (ii)). It then suffices to apply the Jacobian criterion (0, 22.6.1) as well as (0, 19.1.14), taking into account the fact that $\Omega^1_{X/S}$ is a locally free $\sh{O}_X$-module in a neighbourhood of $x$ and $\sh{N}_{Y/X}$ an $\sh{O}_Y$-module of finite type (16.1.6).

Secondly, let us prove that a) implies c'). We may again restrict to the case where $S$, $X$, $Y$ are affine, $B$ and $B/\mathfrak J$ smooth $A$-algebras, and it then follows from (17.7.9) and (8.10.5, (iv)) that there exist a Noetherian subring $A'$ of $A$, an $A'$-algebra of finite type $B'$ and an ideal $\mathfrak J'$ of $B'$ such that $B=B'\otimes_{A'}A$, $\mathfrak J=\mathfrak J'B$ and $B'$ and $B'/\mathfrak J'$ are smooth $A'$-algebras. Note that, by (0, 19.3.8), $B'$ is again a formally smooth $A'$-algebra when $A'$ is given the discrete topology and $B'$ the $\mathfrak J'$-preadic topology. Consequently (0, 19.5.4), $\mathfrak J'/\mathfrak J'^2$ is a projective $(B'/\mathfrak J')$-module and the canonical homomorphism
\[
\mathbf{S}_{B'/\mathfrak J'}(\mathfrak J'/\mathfrak J'^2)\to\operatorname{gr}_{\mathfrak J'}^{\bullet}(B')
\]
is bijective. In other words (16.9.8), the immersion $\Spec(B'/\mathfrak J')\to\Spec(B')$ is quasi-regular, hence regular since $B'$ is Noetherian (16.9.10). But since by hypothesis $B'/\mathfrak J'$ is a flat $A'$-module (17.5.1) and $B'$ an $A'$-algebra of finite presentation, we may apply (11.3.8), replacing $\sh{F}$ by $\widetilde{\mathfrak J'}$, and we therefore see (by base change) that $j$ is a regular immersion, and \emph{a fortiori} quasi-regular. The fact that $B$ is an $A$-algebra of finite presentation and a flat $A$-module shows moreover, by virtue of (11.3.8), that conditions c') and c) are equivalent, and are stable under base change.

Let us finally show that c) implies a). From the fact that, in (11.3.8), condition b) implies c), we already see that if c) is satisfied, $f$ is \emph{flat} at the point $x$; by virtue of (17.5.1), everything therefore comes down to seeing that under hypothesis c), $f^{-1}(s)$ is smooth over $k(s)$ at the point $x$, setting $s=h(x)$; since we have remarked that condition c) is stable under base change, we see that we are reduced to the case where $S=\Spec(k)$ is the spectrum of a field, and taking into account (17.7.1, (ii)) we may suppose that $k$ is \emph{algebraically closed}. It then suffices to prove that $\sh{O}_{X,x}$ is a \emph{regular} ring (17.5.1). Now, by hypothesis we have $\sh{O}_{Y,y}=\sh{O}_{X,x}/\mathfrak J_x$, where $\mathfrak J_x$ is an ideal generated by an $\sh{O}_{X,x}$-regular sequence $(t_i)$, and $\sh{O}_{Y,y}$ is a regular ring; since the $t_i$ form part of a system of parameters for $\sh{O}_{X,x}$ (0, 16.4.1), the conclusion follows from (0, 17.1.7).

\begin{corollary}[17.12.2]
\label{IV.17.12.2}
With the notation of (17.12.1), suppose that $f$ is smooth at the point $x$, and let $Y$ be a closed sub-prescheme of $X$ defined by an ideal $\sh{J}$ of $\sh{O}_X$, and $S$-smooth at the point $x$. Let $(g_i)_{1\leq i\leq r}$ be sections of $\sh{J}$ over $X$. The following conditions are equivalent:
\begin{enumerate}
  \item[a)] The $(g_i)_x$ form a system of generators of the $\sh{O}_{X,x}$-module $\sh{J}_x$ whose number of elements is the smallest possible.
  \item[b)] If $g_i'$ is the canonical image of $g_i$ in $\Gamma(X,\sh{J}/\sh{J}^2)$, the $(g_i')_x$ form a basis of the $(\sh{O}_{X,x}/\sh{J}_x)$-module $\sh{J}_x/\sh{J}_x^2$.

\oldpage[IV-4]{87}

  \item[c)] The images of the $d_{X/S}g_i$ in $(\Omega^1_{X/S})_x\otimes_{\sh{O}_{X,x}}k(x)$ are elements linearly independent over $k(x)$, and the $(g_i)_x$ generate $\sh{J}_x$.
  \item[d)] There exist an open neighbourhood $U$ of $x$ in $X$ and sections $g_j$ ($r+1\leq j\leq n$) of $\sh{O}_X$ over $U$ such that the $S$-morphism $u:U\to S[T_1,\ldots,T_n]$ corresponding to the homomorphism $\sh{O}_S^n\to f_*(\sh{O}_U)$ defined by the sections $g_i|U$ ($1\leq i\leq r$) and $g_j$ ($r+1\leq j\leq n$) (II, 1.2.7) is an étale morphism, and the inverse image of the sub-prescheme $Y'=S[T_{r+1},\ldots,T_n]$ under this morphism is the prescheme induced by $j(Y)$ on the open subset $U\cap j(Y)$.
\end{enumerate}
\end{corollary}

Since $(g_i')_x$ is the canonical image of $(g_i)_x$, the equivalence of a) and b) follows from Nakayama's lemma, $\sh{J}_x$ being of finite type and $\sh{J}_x/\sh{J}_x^2$ a free $(\sh{O}_{X,x}/\sh{J}_x)$-module (17.10.4) (Bourbaki, \emph{Alg. comm.}, chap. II, § 3, n\textsuperscript{o} 2, prop. 5). By virtue of (17.12.1, b)), $\sh{J}_x/\sh{J}_x^2$ is canonically identified with a direct factor of the free $(\sh{O}_{X,x}/\sh{J}_x)$-module of rank $n$, $(\Omega^1_{X/S})_x\otimes_{\sh{O}_{X,x}}(\sh{O}_{X,x}/\sh{J}_x)$, and the equivalence of b) and c) follows from Bourbaki, \emph{loc. cit.} Moreover, if a) is satisfied, $(g_i')_x$ is thus identified with $(d_{X/S}g_i)_x\otimes1$ for $1\leq i\leq r$; the $\sh{O}_{X,x}$-module $(\Omega^1_{X/S})_x$ being free of rank $n$, it again follows from Bourbaki, \emph{loc. cit.}, that, replacing $X$ if necessary by an open neighbourhood of $x$, we may suppose that there exist $n-r$ sections $g_i$ ($r+1\leq i\leq n$) of $\sh{O}_X$ over $X$ such that the $(d_{X/S}g_i)_x$ for $1\leq i\leq n$ form a basis of $(\Omega^1_{X/S})_x$; the fact that the corresponding morphism $u$ is étale at the point $x$ then follows from (17.11.3); for the same reason, the morphism $u^{-1}(Y')\to Y'$, restriction of $u$, is étale at the point $x$; replacing $X$ by a neighbourhood of $x$, we may suppose these two morphisms étale. Moreover, it is immediate that $Y$ (identified, for simplicity, with the closed sub-prescheme $j(Y)$ of $X$) is a closed sub-prescheme of $u^{-1}(Y')$, and by virtue of the choice of the $g_i$ for $i>r$ and of (17.2.5) and (17.11.2), the restriction of $u$ to $Y$ may also be supposed étale. We deduce that for every $y'\in Y'$, the immersion $Y\cap u^{-1}(y')\to u^{-1}(y')$ is open, from which we conclude with the help of (17.9.6) that $Y\to u^{-1}(Y')$ is an open immersion; this proves that a) implies d). The converse follows immediately from (17.11.3).

\begin{corollary}[17.12.3]
\label{IV.17.12.3}
Let $f:X\to S$ be a morphism locally of finite presentation, $u:S\to X$ an $S$-section of $X$. For $f$ to be smooth at a point $x\in X$, it is necessary and sufficient that $u$ be a quasi-regular immersion in a neighbourhood of $s=f(x)$.
\end{corollary}

The conclusion follows from (17.12.1), since $f\circ u=1_S$ is smooth.

\begin{proposition}[17.12.4]
\label{IV.17.12.4}
Every smooth morphism $f:X\to S$ is differentially smooth (in other words (16.10.5) $\Delta_f:X\to X\times_SX$ is a quasi-regular immersion). \emph{In particular, the $\sh{P}_f^n=\sh{P}_{X/S}^n$ and the $\operatorname{gr}_n(\sh{P}_{X/S})$ are locally free $\sh{O}_X$-modules of finite type.}
\end{proposition}

It suffices to remark that the structural morphism $p_1:X\times_SX\to X$ is smooth (17.3.3, (iii)) and that $\Delta_f$ is an $X$-section for $p_1$; the conclusion follows from (17.12.3).

\begin{proposition}[17.12.5]
\label{IV.17.12.5}
Let $f:X\to Y$ be a morphism locally of finite presentation. The following conditions are equivalent:
\begin{enumerate}
  \item[a)] $f$ is differentially smooth.
  \item[b)] For every morphism $g:Y'\to Y$, if we set $X'=X\times_Y Y'$ and $f'=f_{(Y')}:X'\to Y'$, then, for every $Y'$-section $s'$ of $X'$, $f'$ is smooth at all the points of $s'(Y')$.
  \item[c)] The second projection $p_2:X\times_YX\to X$ is smooth at all the points of the diagonal $\Delta_f(X)$.
\end{enumerate}
\end{proposition}

\oldpage[IV-4]{88}

Condition c) is a particular case of b): indeed it suffices to take $Y'=X$ and $g=f$ in b), since then $f'$ is none other than the second projection $p_2$, and $\Delta_f$ is an $X$-section of $X\times_YX$. Secondly, c) implies a), since $p_2$ is a morphism locally of finite presentation, and if $p_2$ is smooth at the points of $\Delta_f(X)$, $\Delta_f$ is a quasi-regular immersion (17.12.3), hence $f$ is differentially smooth. Finally, let us show that a) implies b); if $\Delta_f$ is a quasi-regular immersion, $p_2$ is smooth at all the points of $\Delta_f(X)$ (17.12.3).

Now, if $g':X'\to X$ is the canonical projection, and $v=(g',g')_Y:X'\to X\times_YX$, we verify immediately that the diagram
\[
\xymatrix{
  X\times_Y X \ar[d]_{p_2} & X' \ar[l]_-v \ar[d]^{f'}\\
  X & Y' \ar[l]^{h=g'\circ s'}
}
\]
is commutative and identifies $X'$ with the product $(X\times_YX)\times_XY'$ (I, 3.3.9); taking into account the fact that the diagram (17.4.1.1) identifies $Y'$ with the product of the $(X\times_YX)$-preschemes $X$ and $X'$, we conclude, from the fact that $p_2$ is smooth at every point of $\Delta_f(X)$, that $f'$ is smooth at every point of $s'(Y')$ (17.3.3, (iii)).

\begin{corollary}[17.12.6]
\label{IV.17.12.6}
The notation being that of (8.8.1), suppose $X_\alpha$ quasi-compact, and $f_\alpha:X_\alpha\to S_\alpha$ locally of finite presentation. For $f:X\to S$ to be differentially smooth, it is necessary and sufficient that there exist $\lambda\geq\alpha$ such that $f_\lambda:X_\lambda\to S_\lambda$ is differentially smooth (in which case $f_\mu:X_\mu\to S_\mu$ is differentially smooth for $\mu\geq\lambda$).
\end{corollary}

The sufficiency of the condition and the last assertion follow from (16.10.4). To prove that the condition is necessary, note that $\Delta_f:X\to X\times_SX$ and $p_2:X\times_SX\to X$ are deduced by base change from $\Delta_{f_\lambda}:X_\lambda\to X_\lambda\times_{S_\lambda}X_\lambda$ and $p_{2,\lambda}:X_\lambda\times_{S_\lambda}X_\lambda\to X_\lambda$. By virtue of (17.12.5), for every $z\in\Delta_f(X)$, there exists an affine open neighbourhood $U(z)$ of $z$ in $X\times_SX$ such that $p_2$ is smooth in $U(z)$; by virtue of (8.2.11) and (17.7.8), there exist an index $\lambda(z)$ and a neighbourhood $U_{\lambda(z)}$ of the projection of $z$ in $X_{\lambda(z)}\times_{S_{\lambda(z)}}X_{\lambda(z)}$ such that $U(z)$ is the inverse image of $U_{\lambda(z)}$ and $p_{2,\lambda(z)}$ is smooth in $U_{\lambda(z)}$. Since $X$ is quasi-compact, we can cover $\Delta_f(X)$ by finitely many neighbourhoods $U(z_i)$; by virtue of (8.3.4), there exists a $\lambda$ greater than all the $\lambda(z_i)$ such that the inverse images of the $U_{\lambda(z_i)}$ in $X_\lambda\times_{S_\lambda}X_\lambda$ form a covering of $\Delta_{f_\lambda}(X_\lambda)$; it then follows from (17.12.5) that this index $\lambda$ has the required property.

\begin{examples}[17.12.7]
\label{IV.17.12.7}
Let $S$ be a prescheme, $G$ an $S$-\emph{group prescheme}, locally of finite presentation over $S$. Then, for $G$ to be differentially smooth, it is necessary and sufficient that $G$ be smooth over $S$ at the points of the ``unit section'' $e$ of $G$. The condition is indeed necessary by (17.12.5). Conversely, suppose it is fulfilled; for every morphism $S'\to S$, $G'=G\times_SS'$ is then an $S'$-group prescheme locally of finite presentation over $S'$; we may therefore, by virtue of (17.12.5), restrict to proving that for every $S$-section $s$ of $G$, $f$ is smooth at the points of $s(S)$. Now, if $m:G\times_SG\to G$ is the morphism which defines the group prescheme structure of $G$, $m\circ(s\times1_G):S\times_SG\to G\times_SG\to G$ is an

\oldpage[IV-4]{89}

$S$-isomorphism of the prescheme $G$, the ``left translation'' $\tau_s$, transforming the points of the unit section of $G$ into the points of $s(S)$; we conclude from this that $G$ is smooth over $S$ at the points of $s(S)$.
\end{examples}

\begin{remark}[17.12.8]
\label{IV.17.12.8}
An $S$-group prescheme $G$, of finite type (and even finite) over a locally Noetherian prescheme $S$, can be differentially smooth over $S$ without being smooth (or even \emph{flat}) over $S$. Take for example for $S$ the spectrum of the algebra of dual numbers $D=k[T]/T^2k[T]$ over a field $k$, and for $G$ the spectrum of the direct product $D$-algebra $E=D\oplus k$. We define on $G$ a structure of $S$-group scheme by defining a ``diagonal map'' $\delta$, a homomorphism from $E$ to $E\otimes_DE$: if $e'$, $e''$ are the idempotents that are the canonical images of the unit element $1$ of $D$ in the factors $D$ and $k$ of $E$, one verifies without difficulty that such a diagonal map is obtained by taking
\[
\delta(e')=e'\otimes e'+e''\otimes e'',\qquad
\delta(e'')=e'\otimes e''+e''\otimes e'.
\]
The unit section of the group scheme $G$ corresponds to the homomorphism $E\to D$ which is the identity on $D$ and $0$ on $k$; it is an isomorphism of $S$ onto a connected component of $G$, and \emph{a fortiori} $G$ is differentially smooth over $S$ (17.12.6), but it is clear that $G$ is not $S$-flat.
\end{remark}


\subsection{Transversal morphisms.}
\EGAReaderTarget{ega4.c04.subsec.17-13}{subsection:IV.17.13}
\label{IV.17.13}

\begin{env}[17.13.1]
\label{IV.17.13.1}
Let $S$ be a prescheme, $X$, $Y$, $X'$ three $S$-preschemes, $i:Y\to X$ an $S$-immersion, $f:X'\to X$ an $S$-morphism. Set $Y'=Y\times_XX'$, and let $g:Y'\to Y$, $j:Y'\to X'$ be the canonical projections, so that we have the commutative diagram
\[
\label{IV.17.13.1.1}
\tag{17.13.1.1}
\xymatrix{
  Y \ar[d]_i & Y' \ar[l]_g \ar[d]^j\\
  X & X' \ar[l]^f
}
\]
and $j$ is an $S$-immersion. We then have a commutative diagram of quasi-coherent $\sh{O}_{Y'}$-modules
\[
\label{IV.17.13.1.2}
\tag{17.13.1.2}
\begin{tikzcd}
  g^*(\sh{N}_{Y/X}) \ar{r} \ar{d}{\operatorname{gr}_1(g)} & f^*(\Omega^1_{X/S})\otimes_{\sh{O}_{X'}}\sh{O}_{Y'} \ar{r} \ar{d} & g^*(\Omega^1_{Y/S}) \ar{r} \ar{d} & 0\\
  \sh{N}_{Y'/X'} \ar{r} & \Omega^1_{X'/S}\otimes_{\sh{O}_{X'}}\sh{O}_{Y'} \ar{r} & \Omega^1_{Y'/S} \ar{r} & 0
\end{tikzcd}
\]
where the lower row is the exact sequence (16.4.21) applied to $j$, and the upper row comes from the same exact sequence for $i$ by application of the right exact functor $g^*$ (which therefore leaves it exact); $\operatorname{gr}_1(g)$ is defined in (16.2.1), and the commutativity of the two squares follows from (0, 20.5.7.3) and (0, 20.5.11.3).

\oldpage[IV-4]{90}

We have moreover seen (16.2.2, (iii)) that $\operatorname{gr}_1(g)$ is surjective here, hence we deduce from (17.13.1.2) the exact sequence
\[
\label{IV.17.13.1.3}
\tag{17.13.1.3}
g^*(\sh{N}_{Y/X})\xrightarrow{\alpha}\Omega^1_{X'/S}\otimes_{\sh{O}_{X'}}\sh{O}_{Y'}\to\Omega^1_{Y'/S}\to0.
\]
\end{env}

\begin{proposition}[17.13.2]
\label{IV.17.13.2}
The notation being that of (17.13.1), let $x'$ be a point of $Y'$, $x=g(x')$ its image in $Y$; suppose $X$ and $Y$ smooth over $S$ at the point $x$, and $X'$ smooth over $S$ at the point $x'$. Let $m$, $m-c$ be the relative dimensions of $X$ and $Y$ over $S$ at the point $x$ (17.10.1), and $n$ that of $X'$ over $S$ at the point $x'$. The following conditions are equivalent:
\begin{enumerate}
  \item[a)] $Y'$ is smooth over $S$ at the point $x'$ and of relative dimension $n-c$.
  \item[b)] The homomorphism $\alpha\otimes1:g^*(\sh{N}_{Y/X})\otimes_{\sh{O}_{Y'}}k(x')\to\Omega^1_{X'/S}\otimes_{\sh{O}_{X'}}k(x')$ deduced from the homomorphism $\alpha$ of (17.13.1.3) is injective.
\end{enumerate}
Let $s$ be the canonical image of $x$ in $S$, and suppose that $x'$ is rational over $k(s)$. Then conditions a) and b) are also equivalent to the following:
\begin{enumerate}
  \item[b')] The homomorphism transpose to $\alpha\otimes1$
  \[
  \mathrm{T}_{X'/S}(x')\to(\sh{N}_{Y/X}\otimes_{\sh{O}_Y}k(x'))^\vee=\mathrm{T}_{X/S}(x)/\mathrm{T}_{Y/S}(x)
  \]
  (cf. (16.5.12)) is surjective.
\end{enumerate}
Moreover, when conditions a) and b) are satisfied at $x'$, they are satisfied in a neighbourhood of $x'$ in $Y'$; replacing $X'$ by a neighbourhood of $x'$ if necessary, the homomorphism
\[
\operatorname{gr}_1(g):g^*(\sh{N}_{Y/X})\to\sh{N}_{Y'/X'}
\]
is bijective, and the sequence
\[
\label{IV.17.13.2.1}
\tag{17.13.2.1}
0\to g^*(\sh{N}_{Y/X})\xrightarrow{\alpha}\Omega^1_{X'/S}\otimes_{\sh{O}_{X'}}\sh{O}_{Y'}\to\Omega^1_{Y'/S}\to0
\]
obtained by adding a $0$ to (17.13.1.3) is exact.
\end{proposition}

The fact that if the equivalent conditions a) and b) are satisfied at the point $x'$, they are also satisfied in a neighbourhood of $x'$ in $Y'$, follows from the fact that the set of points where a morphism is smooth is open (17.3.7) and from (17.10.2).

By virtue of (17.12.1) applied to $X'$ and $Y'$, to say that $Y'$ is smooth over $S$ at the point $x'$ is equivalent to saying that the homomorphism
\[
\delta\otimes1:\sh{N}_{Y'/X'}\otimes_{\sh{O}_{Y'}}k(x')\to\Omega^1_{X'/S}\otimes_{\sh{O}_{X'}}k(x')
\]
is injective, taking into account (0, 19.1.12) and the fact that $\Omega^1_{X'/S}$ is an $\sh{O}_{X'}$-module locally free at the point $x'$ (17.2.3); since then $\Omega^1_{Y'/S}$ is also an $\sh{O}_{Y'}$-module locally free in a neighbourhood of $x'$ and the sequence
\[
\label{IV.17.13.2.2}
\tag{17.13.2.2}
0\to\sh{N}_{Y'/X'}\to\Omega^1_{X'/S}\otimes_{\sh{O}_{X'}}\sh{O}_{Y'}\to\Omega^1_{Y'/S}\to0
\]
is exact (17.2.5), to say that $Y'$ is of relative dimension $n-c$ over $S$ at the point $x'$ means, by (17.10.2), that $\sh{N}_{Y'/X'}$ (which is locally free in a neighbourhood of $x'$) is of rank $c$ at the point $x'$. But by virtue of the hypotheses on $X$ and $Y$ at the point $x$, and by the same reasoning, $\sh{N}_{Y/X}$ is locally free and of rank $c$ at the point $x$, hence $g^*(\sh{N}_{Y/X})$ is also locally free and of rank $c$ at the point $x'$. Since the homomorphism $\operatorname{gr}_1(g):g^*(\sh{N}_{Y/X})\to\sh{N}_{Y'/X'}$ is surjective, the preceding conditions are equivalent to saying that this homomorphism is

\oldpage[IV-4]{91}

bijective at the point $x'$ (Bourbaki, \emph{Alg. comm.}, chap. II, § 3, n\textsuperscript{o} 2, cor. of prop. 6), hence also in a neighbourhood of $x'$ (0$_{\mathrm I}$, 5.2.7); this evidently implies b), as well as the last assertion of the statement, by virtue of the exactness of (17.13.2.2). Conversely, since $\alpha\otimes1$ factors as
\[
g^*(\sh{N}_{Y/X})\otimes k(x')\xrightarrow{\operatorname{gr}_1(g)\otimes1}\sh{N}_{Y'/X'}\otimes k(x')\xrightarrow{\delta\otimes1}\Omega^1_{X'/S}\otimes k(x')
\]
and $\operatorname{gr}_1(g)$ is surjective, to say that $\alpha\otimes1$ is injective implies that $\delta\otimes1$ is injective and that $\operatorname{gr}_1(g)\otimes1$ is bijective. We conclude from this (17.12.1) that $Y'$ is smooth over $S$ at the point $x'$, and that $\operatorname{gr}_1(g)$ is bijective in a neighbourhood of $x'$ (Bourbaki, \emph{loc. cit.}). Moreover, since the sequence (17.13.2.2) is then exact, and $\sh{N}_{Y'/X'}$, being isomorphic to $g^*(\sh{N}_{Y/X})$, is of rank $c$ at the point $x'$, $\Omega^1_{Y'/S}$ is of rank $n-c$ at this point, which completes the proof of the equivalence of a) and b), by virtue of (17.10.2).

It remains to show the equivalence of b) and b') when $x'$ is rational over $k(s)$ (which implies that the same holds for $x$); then $g^*(\sh{N}_{Y/X})\otimes_{\sh{O}_{Y'}}k(x')$ is identified with $\sh{N}_{Y/X}\otimes_{\sh{O}_Y}k(x')$, and since the sequence
\[
0\to\sh{N}_{Y/X}\to\Omega^1_{X/S}\otimes_{\sh{O}_X}\sh{O}_Y\to\Omega^1_{Y/S}\to0
\]
is exact at the point $x$ and formed of locally free $\sh{O}_Y$-modules, the dual of $\sh{N}_{Y/X}\otimes_{\sh{O}_Y}k(x)$ is identified with the quotient space $\mathrm{T}_{X/S}(x)/\mathrm{T}_{Y/S}(x)$ (16.5.12); whence the equivalence of b) and b').

\begin{definition}[17.13.3]\label{IV.17.13.3}
\EGAReaderTarget{ega4.c04.def.17-13-3}{IV.17.13.3}
With the notation of (17.13.1), we say that the morphism $f$ is \emph{transversal} to $Y$ at the point $x'$, relative to $S$, if $X$ and $Y$ are smooth over $S$ at the point $x$, $X'$ smooth over $S$ at the point $x'$, and if the equivalent conditions a), b) of (17.13.2) are satisfied.
\end{definition}

When no confusion is to be feared, we omit mention of the prescheme $S$ and simply say that $f$ is transversal to $Y$ at the point $x'$.

\begin{remarks}[17.13.4]
\label{IV.17.13.4}
(i) Suppose that $X$, $Y$ and $X'$ are \emph{flat and locally of finite presentation} over $S$. For every $s\in S$, denote by $X_s$, $Y_s$, $X'_s$, $Y'_s$ the fibres of $X$, $Y$, $X'$, $Y'$ at the point $s$, and by $f_s:X'_s\to X_s$ the morphism deduced from $f$ by base change. It then follows from (17.5.9) that for $f$ to be a morphism transversal to $Y$ at the point $x'$, it is necessary and sufficient that, if $s$ is the image of $x$ in $S$, $f_s$ be transversal to $Y_s$ at the point $x'$, relative to $k(s)$.

(ii) We have seen in (17.13.2) that the set of points $x'\in Y'$ where $f$ is transversal to $Y$ is \emph{open} in $Y'$; if moreover $Y'$ is \emph{proper} over $S$, we deduce that the set of $s\in S$ such that $f$ is transversal to $Y$ at all the points of $Y'_s$, relative to $S$, is \emph{open} in $S$.

When $X$, $Y$ and $X'$ are flat and locally of finite presentation over $S$, it follows from (i) that the set of $x'\in Y'$ such that $f_s$ is transversal to $Y_s$ at the point $x'$ ($s$ the image of $x'$ in $S$), relative to $k(s)$, is \emph{open} in $Y'$. If moreover $Y'$ is proper over $S$, the set of $s\in S$ such that $f_s$ is transversal to $Y_s$ at all the points of $Y'_s$ (relative to $k(s)$) is \emph{open} in $S$.

(iii) Since the property of being smooth over $S$, as well as the notion of relative dimension over $S$ at a point, are stable under every base change $S'\to S$ ((17.3.3) and (4.2.7)), the same holds for the property, for a morphism $f:X'\to X$, of being transversal to a sub-prescheme of $X$ at a point, relative to $S$.

\oldpage[IV-4]{92}

(iv) Condition b) of (17.13.2) also expresses that the homomorphism $\alpha:g^*(\sh{N}_{Y/X})\to\Omega^1_{X'/S}\otimes_{\sh{O}_{X'}}\sh{O}_{Y'}$ is \emph{universally injective} relative to $Y'$ (II, 9.18).
\end{remarks}

\begin{env}[17.13.5]
\label{IV.17.13.5}
Consider now a prescheme $S$, three $S$-preschemes $X$, $Y$, $Z$, and two $S$-morphisms $f:Y\to X$, $g:Z\to X$; set $T=Y\times_XZ$; we then know (I, 5.3.5) that we have a commutative diagram
\[
\label{IV.17.13.5.1}
\tag{17.13.5.1}
\xymatrix{
  X \ar[d]^{\Delta} & T=Y\times_X Z \ar[l] \ar[d]^{j}\\
  X\times_S X & Y\times_S Z \ar[l]^{u}
}
\]
making $T$ the product of the $(X\times_SX)$-preschemes $X$ and $Y\times_SZ$, where $u=f\times_Sg$. Since $\Delta$ is an $S$-immersion (I, 5.3.9), we are in the situation of the diagram (17.13.1.1); what corresponds to $\Omega^1_{X'/S}$ in (17.13.1) is then $(\Omega^1_{Y/S}\otimes_{\sh{O}_Y}\sh{O}_{Y\times_SZ})\oplus(\Omega^1_{Z/S}\otimes_{\sh{O}_Z}\sh{O}_{Y\times_SZ})$ by virtue of (16.4.23). On the other hand, what corresponds to the $\sh{O}_Y$-module $\sh{N}_{Y/X}$ in (17.13.1) is here by definition $\Omega^1_{X/S}$ (16.3.1); to (17.13.1.3) there therefore corresponds an exact sequence
\[
\label{IV.17.13.5.2}
\tag{17.13.5.2}
\Omega^1_{X/S}\otimes_{\sh{O}_X}\sh{O}_T\xrightarrow{\rho}
(\Omega^1_{Y/S}\otimes_{\sh{O}_Y}\sh{O}_T)\oplus
(\Omega^1_{Z/S}\otimes_{\sh{O}_Z}\sh{O}_T)
\xrightarrow{\sigma}\Omega^1_{T/S}\to0
\]
where it remains to make the homomorphisms $\rho$ and $\sigma$ explicit. Taking into account first (16.4.23) and (0, 20.5.2), we see that if $p:T\to Y$, $q:T\to Z$ are the canonical projections, we have, with the notation of (16.4.19),
\[
\label{IV.17.13.5.3}
\tag{17.13.5.3}
\sigma=p_{T/Y/S}+q_{T/Z/S}.
\]
On the other hand, to evaluate $\rho$, let us use the commutativity of the left-hand square in (17.13.1.2), which, in the present case, comes down first to making explicit the canonical homomorphism
\[
\rho':\Omega^1_{X/S}\to\Omega^1_{(X\times_SX)/S}\otimes_{\sh{O}_{X\times_SX}}\sh{O}_X
\]
defined in (16.4.21) applied to the immersion $\Delta$. We may restrict to the case where $S=\Spec(A)$, $X=\Spec(B)$ are affine, and then $\rho'$ corresponds to the homomorphism $\delta$ of (0, 20.5.11.2), where $B$ must be replaced by $B\otimes_AB$ and $C$ by $B=(B\otimes_AB)/\mathfrak J_{B/A}$. We then see that $\delta$ maps the class of $x\otimes1-1\otimes x\pmod{\mathfrak J_{B/A}^2}$ (for an $x\in B$) to the image of
\[
(x\otimes1-1\otimes x)\otimes(1\otimes1)-(1\otimes1)\otimes(x\otimes1-1\otimes x)
\]
in $\Omega^1_{(B\otimes_AB)/A}\otimes_{B\otimes_AB}B$; but the preceding element can be written
\[
\begin{split}
&((x\otimes1)\otimes(1\otimes1)-(1\otimes1)\otimes(x\otimes1))\\
&\quad-((1\otimes x)\otimes(1\otimes1)-(1\otimes1)\otimes(1\otimes x))
\end{split}
\]
and we therefore see that $\rho'$ is the difference of the two homomorphisms $\pi^{(1)}$ and $\pi^{(2)}$ from $\Omega^1_{X/S}$ to $\Omega^1_{(X\times_SX)/S}\otimes_{\sh{O}_{X\times_SX}}\sh{O}_X$ corresponding respectively to the first and to the second projection of $X\times_SX$ to $X$, by (16.4.3.3). To obtain $\rho$, we must first consider the homomorphism
\[
\rho'':\Omega^1_{(X\times_SX)/S}\otimes_{\sh{O}_{X\times_SX}}\sh{O}_{Y\times_SZ}\to\Omega^1_{(Y\times_SZ)/S}
\]
corresponding by (16.4.3.3) to the morphism $u$ of (17.13.5.1), then, after tensoring with $\sh{O}_T$, form the composite $(\rho''\otimes1)\circ(\rho'\otimes1)$; it follows from the above that we have, with the notation of (16.4.18),
\[
\label{IV.17.13.5.4}
\tag{17.13.5.4}
\rho=(f_{Y/X/S}\otimes1_{\sh{O}_T})\oplus(-g_{Z/X/S}\otimes1_{\sh{O}_T}).
\]

\oldpage[IV-4]{93}

This being said, the application of (17.13.2) to the situation of the diagram (17.13.5.1) (taking into account (17.3.3, (iv)), which implies that $X\times_SX$ is smooth over $S$ at $x$ if $X$ is) gives the following corollary.
\end{env}

\begin{corollary}[17.13.6]\label{IV.17.13.6}
\EGAReaderTarget{ega4.c04.cor.17-13-6}{IV.17.13.6}
Let $S$ be a prescheme, $X$, $Y$, $Z$ three $S$-preschemes, $f:Y\to X$, $g:Z\to X$ two $S$-morphisms; set $T=Y\times_XZ$, and let $p:T\to Y$, $q:T\to Z$ be the canonical projections. Let $t$ be a point of $T$, $y=p(t)$, $z=q(t)$, $x=f(y)=g(z)$. Suppose that $X$ is smooth over $S$ at the point $x$, of relative dimension $m$, and $Y$ (resp. $Z$) smooth over $S$ at the point $y$ (resp. $z$), of relative dimension $m+a$ (resp. $m+b$), $a$ and $b$ being positive or negative. Then the following conditions are equivalent:
\begin{enumerate}
  \item[a)] $T$ is smooth over $S$ at the point $t$, of relative dimension $m+a+b$.
  \item[b)] The homomorphism
  \[
  \rho\otimes1:\Omega^1_{X/S}\otimes_{\sh{O}_X}k(t)\to
  (\Omega^1_{Y/S}\otimes_{\sh{O}_Y}k(t))\oplus
  (\Omega^1_{Z/S}\otimes_{\sh{O}_Z}k(t))
  \]
  where $\rho$ is given by (17.13.5.4), is injective.
  \item[c)] The morphism $T=Y\times_SZ\to X\times_SX$ is transversal to the diagonal of $X\times_SX$ at the point $t$.
\end{enumerate}
If $t$ is rational over the residue field of its image $s$ in $S$, these conditions are also equivalent to the following:
\begin{enumerate}
  \item[b')] The homomorphism
  \[
  \label{IV.17.13.6.1}
  \tag{17.13.6.1}
  \mathrm{T}_y(f)-\mathrm{T}_z(g):\mathrm{T}_{Y/S}(y)\oplus\mathrm{T}_{Z/S}(z)\to\mathrm{T}_{X/S}(x)
  \]
  (cf. (16.5.12.5)) is surjective.
\end{enumerate}
Moreover, when conditions a) and b) are satisfied at $t$, they are satisfied in a neighbourhood of $t$ in $T$, and, on restricting $T$ to such a neighbourhood, the sequence
\[
\label{IV.17.13.6.2}
\tag{17.13.6.2}
0\to\Omega^1_{X/S}\otimes_{\sh{O}_X}\sh{O}_T\xrightarrow{\rho}
(\Omega^1_{Y/S}\otimes_{\sh{O}_Y}\sh{O}_T)\oplus
(\Omega^1_{Z/S}\otimes_{\sh{O}_Z}\sh{O}_T)
\xrightarrow{\sigma}\Omega^1_{T/S}\to0
\]
(where $\sigma$ is given by (17.13.5.3)) is exact.

The only point that remains to be proved in (17.13.6) concerns b'), and follows from the fact that, under the hypothesis that $t$ is rational over $k(s)$, the homomorphism $\rho\otimes1$ is the \emph{transpose} of $\mathrm{T}_y(f)-\mathrm{T}_z(g)$, by virtue of (17.13.5.4).

When the equivalent conditions of (17.13.6) are satisfied, we say that $f$ and $g$ form a \emph{pair of morphisms transversal at the point $t$, relative} to $S$; here again, mention of $S$ is often omitted.
\end{corollary}

\begin{env}[17.13.7]\label{IV.17.13.7}
\EGAReaderTarget{ega4.c04.statement.17-13-7}{IV.17.13.7}
Consider in particular the case where $Y$ and $Z$ are \emph{sub-preschemes} of $X$, $f$ and $g$ being the canonical injections; $T$ is then the ``intersection'' sub-prescheme

\oldpage[IV-4]{94}

$\inf(Y,Z)$ of $X$ (I, 4.4.3), and we have $t=y=z=x$. Instead of saying that the pair $(f,g)$ is transversal at the point $x$, we then say that $Y$ and $Z$ \emph{intersect transversally at the point $x$} (relative to $S$). Denoting by $X_s$, $Y_s$, $Z_s$, $T_s$ the fibres of $X$, $Y$, $Z$, $T$ at the point $s$, the image of $x$ in $S$, and taking into account (5.2.3), (5.1.9) and (0, 16.5.12), we see that for this to be the case (when $X$, $Y$, $Z$ are smooth over $S$ at the point $x$), it is necessary and sufficient that $T$ be smooth over $S$ at the point $x$ and that we have the relation
\[
\label{IV.17.13.7.1}
\tag{17.13.7.1}
\operatorname{codim}_x(T_s,X_s)=\operatorname{codim}_x(Y_s,X_s)+\operatorname{codim}_x(Z_s,X_s).
\]
\end{env}

\begin{proposition}[17.13.8]
\label{IV.17.13.8}
Let $S$ be a prescheme, $X$ an $S$-prescheme locally of finite presentation over $S$, $Y$, $Z$ two sub-preschemes of $X$, $T=\inf(Y,Z)$ the ``intersection'' sub-prescheme of $Y$ and $Z$, and $x$ a point of $T$. The following conditions are equivalent:
\begin{enumerate}
  \item[a)] The canonical injection $i:Y\to X$ is a morphism transversal to $Z$ at the point $x$, relative to $S$ (17.13.3).
  \item[a')] The canonical injection $j:Z\to X$ is a morphism transversal to $Y$ at the point $x$, relative to $S$.
  \item[a'')] $Y$ and $Z$ intersect transversally at the point $x$, relative to $S$.
  \item[b)] $X$, $Y$, $Z$ are smooth over $S$ at the point $x$, and the homomorphism $\rho$ (17.13.5.4) is such that
  \[
  \rho\otimes1:\Omega^1_{X/S}\otimes_{\sh{O}_X}k(x)\to
  (\Omega^1_{Y/S}\otimes_{\sh{O}_Y}k(x))\oplus
  (\Omega^1_{Z/S}\otimes_{\sh{O}_Z}k(x))
  \]
  is injective.
\end{enumerate}
When $x$ is rational over $k(s)$ ($s$ the image of $x$ in $S$), these conditions are also equivalent to the following:
\begin{enumerate}
  \item[b')] The homomorphism
  \[
  \mathrm{T}_x(i)-\mathrm{T}_x(j):\mathrm{T}_{Y/S}(x)\oplus\mathrm{T}_{Z/S}(x)\to\mathrm{T}_{X/S}(x)
  \]
  is surjective.
\end{enumerate}
Moreover, when the equivalent conditions a) to b) are satisfied at the point $x$, they are satisfied in a neighbourhood of $x$ in $T$, and, on restricting $X$ to a neighbourhood of $x$, the sequence (17.13.6.2) is exact, and we have a canonical isomorphism
\[
\label{IV.17.13.8.1}
\tag{17.13.8.1}
\sh{N}_{T/X}\xrightarrow{\sim}
(\sh{N}_{Y/X}\otimes_{\sh{O}_Y}\sh{O}_T)\oplus
(\sh{N}_{Z/X}\otimes_{\sh{O}_Z}\sh{O}_T).
\]
\end{proposition}

Conditions a), a'), a'') all imply that $X$, $Y$, $Z$ are smooth over $S$ at the point $x$. Moreover, let $m$, $m-a$, $m-b$ be the relative dimensions of $X$, $Y$, $Z$ over $S$ at the point $x$. It then follows from (17.13.2), applied with $X'$ replaced by $Z$ and $Y'$ by $T=Y\times_XZ$, that conditions a), a') are both equivalent to saying that $T$ is smooth over $S$ at the point $x$ and of relative dimension $m-a-b$ at this point; but by virtue of (17.13.7), this means precisely that condition a'') is satisfied, whence the equivalence of a), a') and a''). The equivalence of a'') and b) (or b') when $x$ is rational over $k(s)$) was proved in (17.13.6), as was the fact that if these conditions are satisfied at the point $x$, they are satisfied in a neighbourhood of $x$, and the exactness of the sequence (17.13.6.2) in such a neighbourhood. It remains to define the canonical isomorphism (17.13.8.1).

\oldpage[IV-4]{95}

Denote by $\alpha$ and $-\beta$ the homomorphisms appearing on the right-hand side of (17.13.5.4), and by $\gamma$ and $\delta$ those appearing on the right-hand side of (17.13.5.2). To say that the sequence (17.13.6.2)
\[
0\to\Omega^1_{X/S}\otimes_{\sh{O}_X}\sh{O}_T
\xrightarrow{\alpha\oplus(-\beta)}
(\Omega^1_{Y/S}\otimes_{\sh{O}_Y}\sh{O}_T)\oplus
(\Omega^1_{Z/S}\otimes_{\sh{O}_Z}\sh{O}_T)
\xrightarrow{\gamma+\delta}\Omega^1_{T/S}\to0
\]
is exact means that, in the category of $\sh{O}_T$-modules, $\Omega^1_{X/S}\otimes_{\sh{O}_X}\sh{O}_T$ is canonically identified with the \emph{fibre product} of $\Omega^1_{Y/S}\otimes_{\sh{O}_Y}\sh{O}_T$ and $\Omega^1_{Z/S}\otimes_{\sh{O}_Z}\sh{O}_T$ over $\Omega^1_{T/S}$, for the homomorphisms $\gamma$ and $\delta$. The same reasoning as in (0, 18.1.2 and 18.1.3), where rings and two-sided ideals are replaced respectively by modules and submodules, yields a commutative diagram
{\scriptsize
\[
\begin{tikzcd}[column sep=small,row sep=large]
0 \ar[dr] & & 0 \ar[d] & 0 \ar[d] & \\
& (\sh{N}_{Y/X}\otimes\sh{O}_T)\oplus(\sh{N}_{Z/X}\otimes\sh{O}_T)
  \ar[r] \ar[d] \ar[dr,"\iota" description]
& \sh{N}_{Z/X}\otimes\sh{O}_T \ar[r,"\sim"] \ar[d]
& \sh{N}_{T/Y} \ar[d] & \\
0 \ar[r] & \sh{N}_{Y/X}\otimes\sh{O}_T \ar[r] \ar[d]
& \Omega^1_{X/S}\otimes\sh{O}_T \ar[r,"\alpha"] \ar[d,"\beta"] \ar[dr,"\varepsilon" description]
& \Omega^1_{Y/S}\otimes\sh{O}_T \ar[r] \ar[d,"\gamma"] & 0\\
0 \ar[r] & \sh{N}_{T/Z} \ar[r]
& \Omega^1_{Z/S}\otimes\sh{O}_T \ar[r,"\delta"] \ar[d]
& \Omega^1_{T/S} \ar[r] \ar[d] \ar[dr] & 0\\
& & 0 & 0 & 0
\end{tikzcd}
\]
}
where the 3\textsuperscript{rd} and 4\textsuperscript{th} rows and columns are exact by virtue of the smoothness hypotheses. The fact that the composite homomorphism $\varepsilon=\gamma\circ\alpha=\delta\circ\beta$ is the homomorphism corresponding to the canonical injection $T\to X$ follows from (0, 20.5.2.7). Since the diagonal of the preceding diagram is exact, $\operatorname{Ker}(\varepsilon)$ is canonically identified with $\sh{N}_{T/X}$, taking into account the fact that $T$ is smooth over $S$ (17.2.5); hence we have a canonical isomorphism (17.13.8.1), inverse to $\iota$.

\begin{env}[17.13.9]
\label{IV.17.13.9}
We can generalize the results of (17.13.5) and (17.13.6) to an arbitrary finite number of $S$-morphisms $f_i:Y_i\to X$ ($i\in I$, $I$ a finite set). For this, denote by $X^I$ the product of the family $(X_i)_{i\in I}$ of $S$-preschemes all identical to $X$ (I, 3.3.5), and let $p_i$ ($i\in I$) be the canonical projections. The \emph{diagonal morphism} $\Delta:X\to X^I$ is defined (as in (I, 5.3.1)) as the unique $S$-morphism such that $p_i\circ\Delta=1_X$ for every $i\in I$. It is an $X$-section of $X^I$ for the morphism $p_1$, hence (I, 5.3.11) an \emph{immersion}.

\oldpage[IV-4]{96}

Let then $Y$ be the product of the $S$-preschemes $Y_i$ (for the composite morphisms $Y_i\xrightarrow{f_i}X\to S$), and $T$ the product of the $X$-preschemes $Y_i$ (for the morphisms $f_i$). We prove as in (I, 5.3.5) that we have a commutative diagram
\[
\label{IV.17.13.9.1}
\tag{17.13.9.1}
\xymatrix{
  X \ar[d]^{\Delta} & T \ar[l]_-v \ar[d]^{j}\\
  X^I & Y \ar[l]^{u}
}
\]
(where $u$ is the product (over $S$) of the $f_i$), which makes $T$ the product of the $X^I$-preschemes $X$ and $Y$.

By induction on the number of elements of $I$, it follows from (16.4.23) that $\Omega^1_{X^I/S}$ (resp. $\Omega^1_{Y/S}$) is canonically identified with the direct sum of the $p_i^*(\Omega^1_{X/S})$ (resp. of the $q_i^*(\Omega^1_{Y_i/S})$, if $q_i:Y\to Y_i$ are the canonical projections).

Suppose now that $X$ is \emph{smooth} over $S$ at a point $x$ (hence in a neighbourhood of this point). The same then holds for $X^I$ at this point (17.3.3, (iv)), and, restricting $X$ to a neighbourhood of $x$, we therefore have (17.2.5) an exact sequence of locally free $\sh{O}_X$-modules
\[
0\to\operatorname{gr}_1(\Delta)\to(\Omega^1_{X/S})^I\to\Omega^1_{X/S}\to0.
\]
Set $\sh{J}=\operatorname{gr}_1(\Delta)$; from the preceding sequence we deduce an exact sequence of locally free $\sh{O}_T$-modules
\[
\label{IV.17.13.9.1-seq}
\tag{17.13.9.1}
0\to\sh{J}\otimes_{\sh{O}_X}\sh{O}_T\to(\Omega^1_{X/S})^I\otimes_{\sh{O}_X}\sh{O}_T\xrightarrow{\sigma}\Omega^1_{X/S}\otimes_{\sh{O}_X}\sh{O}_T\to0
\]
which corresponds to the first row of the diagram (17.13.1.2), and where $\sigma$ is none other than the canonical homomorphism which, to each family of sections $(t_i)_{i\in I}$ of $\Omega^1_{X/S}\otimes_{\sh{O}_X}\sh{O}_T$ over an open subset of $T$, assigns its \emph{sum}.

On the other hand, what corresponds here to the second vertical arrow of the diagram (17.13.1.2) is the homomorphism
\[
\label{IV.17.13.9.2}
\tag{17.13.9.2}
\tau:(\Omega^1_{X/S})^I\otimes_{\sh{O}_X}\sh{O}_T\to\bigoplus_{i\in I}(\Omega^1_{Y_i/S}\otimes_{\sh{O}_{Y_i}}\sh{O}_T)
\]
which, to every family $(t'_i\otimes1)_{i\in I}$, where here the $t'_i$ are sections of $\Omega^1_{X/S}$ over an open subset of $X$, associates the sum of the $((f_i)_{Y_i/X/S}(t'_i))\otimes1$, with the notation of (16.4.18). The homomorphism $\rho$ which corresponds to the homomorphism $\alpha$ of (17.13.1.3) is therefore the \emph{restriction} to $\sh{J}\otimes_{\sh{O}_X}\sh{O}_T$ of the preceding homomorphism $\tau$.

These hypotheses and notation having been specified, we can apply prop. (17.13.2) to the situation of (17.13.9.1), which gives the
\end{env}

\begin{corollary}[17.13.10]\label{IV.17.13.10}
\EGAReaderTarget{ega4.c04.cor.17-13-10}{IV.17.13.10}
Under the general hypotheses of (17.13.9), let $t$ be a point of $T$, $y_i$ ($i\in I$) its projections in the $Y_i$, and $x$ the point of $X$ equal to each of the $f_i(y_i)$. Suppose that $X$ is smooth over $S$ at the point $x$ and of relative dimension $d$, and that each of the $Y_i$ is smooth over $S$ at the point $y_i$; let $d+c_i$ ($c_i$ a positive or negative integer) be the relative dimension of $Y_i$ at the point $y_i$. Then the following conditions are equivalent:

\oldpage[IV-4]{97}

\begin{enumerate}
  \item[a)] $T$ is smooth over $S$ at the point $t$, of relative dimension $d+\sum_{i\in I}c_i$.
  \item[b)] The homomorphism
  \[
  \rho\otimes1:\sh{J}\otimes_{\sh{O}_X}k(t)\to\bigoplus_{i\in I}(\Omega^1_{Y_i/S}\otimes_{\sh{O}_{Y_i}}k(t))
  \]
  is injective.
\end{enumerate}
When $t$ is rational over the field $k(s)$ (where $s$ is the image of $t$ in $S$), these conditions are also equivalent to the following:
\begin{enumerate}
  \item[b')] The homomorphism transpose to $\rho\otimes1$
  \[
  \bigoplus_{i\in I}\mathrm{T}_{Y_i/X}(y_i)\to(\mathrm{T}_{X/S}(x))^I/\delta(\mathrm{T}_{X/S}(x))
  \]
  (where $\delta$ is the diagonal map) is surjective.
\end{enumerate}
It suffices to remark that $X^I$ is smooth over $S$ and of relative dimension $d\cdot\operatorname{Card}(I)$, and $Y$ smooth over $S$ of relative dimension $d\cdot\operatorname{Card}(I)+\sum_{i\in I}c_i$.

When the equivalent conditions of (17.13.10) are satisfied, we again say that the $f_i$ form a \emph{family of morphisms transversal at the point $t$, relative} to $S$. We again see that the set of points $t\in T$ where this holds is \emph{open} in $T$. Remark (17.13.4, (iii)) then shows that if $X$ and the $Y_i$ are flat and locally of finite presentation over $S$, and if moreover $T$ is \emph{proper} over $S$, the set of $s\in S$ such that the $f_i$ form a family of transversal morphisms at \emph{all} the points of $T_s$, relative to $S$, is \emph{open} in $S$.
\end{corollary}

\begin{env}[17.13.11]\label{IV.17.13.11}
\EGAReaderTarget{ega4.c04.statement.17-13-11}{IV.17.13.11}
Consider in particular the case, generalizing (17.13.7), where the $Y_i$ ($i\in I$) are \emph{sub-preschemes} of $X$, the $f_i$ being the canonical injections, so that $T=\inf_{i\in I}(Y_i)$ is again the ``intersection'' sub-prescheme of the $Y_i$, and $t=y_i=x$ for every $i$; instead of saying that the $f_i$ form a family of transversal morphisms at the point $x$, we again say that \emph{the $Y_i$ intersect transversally at the point $x$} (relative to $S$). Condition a) of (17.13.10) is again expressed by the relation, generalizing (17.13.7.1),
\[
\label{IV.17.13.11.1}
\tag{17.13.11.1}
\operatorname{codim}_x(T_s,X_s)=\sum_{i\in I}\operatorname{codim}_x((Y_i)_s,X_s).
\]
Moreover, we have the following property, which extends (17.13.8.1) and gives another proof of it when $I$ has 2 elements:
\end{env}

\begin{corollary}[17.13.12]
\label{IV.17.13.12}
When the $Y_i$ intersect transversally at the point $x$, we have a canonical isomorphism
\[
\label{IV.17.13.12.1}
\tag{17.13.12.1}
\sh{N}_{T/X}\xrightarrow{\sim}\bigoplus_{i\in I}(\sh{N}_{Y_i/X}\otimes_{\sh{O}_{Y_i}}\sh{O}_T).
\]
\end{corollary}

We may restrict to the case where the $Y_i$ are closed sub-preschemes of $X$, defined by quasi-coherent ideals $\sh{J}_i$, so that $T$ is defined by the ideal $\sh{J}=\sum_i\sh{J}_i$. By definition of the conormal sheaf of an immersion (16.1.3), the canonical homomorphism $\bigoplus_i\sh{J}_i\to\sh{J}$ gives, by passage to quotients, a \emph{surjective} homomorphism
\[
\label{IV.17.13.12.2}
\tag{17.13.12.2}
\bigoplus_{i\in I}(\sh{N}_{Y_i/X}\otimes_{\sh{O}_{Y_i}}\sh{O}_T)\to\sh{N}_{T/X}.
\]

\oldpage[IV-4]{98}

But here the $\sh{O}_T$-modules on the two sides of (17.13.12.2) are locally free and \emph{of the same rank} $\sum_i c_i$ (if $c_i$ is the rank of $\sh{N}_{Y_i/X}$), by virtue of (17.2.5) and of condition a) of (17.13.10); we therefore conclude from Bourbaki, \emph{Alg. comm.}, chap. II, § 3, n\textsuperscript{o} 2, cor. of prop. 6, that (17.13.12.2) is bijective, and (17.13.12.1) is the inverse isomorphism.

\subsection{Local and infinitesimal characterizations of smooth morphisms, unramified morphisms and étale morphisms.}
\EGAReaderTarget{ega4.c04.subsec.17-14}{subsection:IV.17.14}

\begin{proposition}[17.14.1]
\label{IV.17.14.1}
Let $f:X\to Y$ be a morphism locally of finite presentation, $x$ a point of $X$, $y=f(x)$. For $f$ to be smooth (resp. unramified, resp. étale) at the point $x$, it is necessary and sufficient that the following condition be satisfied:

For every local scheme $Y'=\operatorname{Spec}(A')$ with closed point $y'$, every morphism $h:Y'\to Y$ such that $h(y')=y$, every closed subscheme $Y'_0=\operatorname{Spec}(A'/\mathfrak{J}')$ of $Y'$, where the ideal $\mathfrak{J}'$ is of square zero, and every $Y$-morphism $g_0:Y'_0\to X$ such that $g_0(y')=x$, there exists at least one (resp. at most one, resp. one and only one) $Y$-morphism $g:Y'\to X$ whose restriction to $Y'_0$ is $g_0$.
\end{proposition}

Taking into account the definitions (17.1.1 and 17.3.1), we have to show that the condition of the statement is \emph{sufficient} for $f$ to be smooth (resp. unramified) at the point $x$.

\emph{(i) Case of smooth morphisms.} We may restrict to the case where $Y=\operatorname{Spec}(A)$ and $X=\operatorname{Spec}(C)$ are affine, with $C=B/\mathfrak{R}$, where $B=A[T_1,\ldots,T_n]$ is a polynomial $A$-algebra and $\mathfrak{R}$ an ideal of finite type of $B$. To prove that $f$ is smooth at the point $x$, it suffices to establish that $\sh{O}_{X,x}$ is a formally smooth $\sh{O}_{Y,y}$-algebra for the discrete topologies (17.5.1). Now we have $\sh{O}_{X,x}=B_{\mathfrak q}/\mathfrak{R}B_{\mathfrak q}$, where $\mathfrak q$ is a prime ideal of $B$, and if $\mathfrak p$ is the inverse image of $\mathfrak q$ in $A$, $B_{\mathfrak q}$ is a formally smooth $A_{\mathfrak p}$-algebra for the discrete topologies (0, 19.3.2 and 19.3.5). By application of the Jacobian criterion (0, 22.6.1 and 20.5.12) it therefore suffices to see that $B_{\mathfrak q}/\mathfrak{R}^2B_{\mathfrak q}$ is an $A_{\mathfrak p}$-trivial extension of $B_{\mathfrak q}/\mathfrak{R}B_{\mathfrak q}$ by $\mathfrak{R}B_{\mathfrak q}/\mathfrak{R}^2B_{\mathfrak q}$. But this follows precisely from the hypothesis applied to $A'=B_{\mathfrak q}/\mathfrak{R}^2B_{\mathfrak q}$ and $\mathfrak{J}'=\mathfrak{R}B_{\mathfrak q}/\mathfrak{R}^2B_{\mathfrak q}$ and to the morphism $g_0$ corresponding to the identity automorphism of $A'/\mathfrak{J}'=\sh{O}_{X,x}$.

\emph{(ii) Case of unramified morphisms.} Here set $A=\sh{O}_{Y,y}$, $B=\sh{O}_{X,x}$, and note that by virtue of (17.4.1, c')), it suffices to show that $\Omega^1_{B/A}=(\Omega^1_{X/Y})_x=0$. Now, with the notation of (0, 20.4.1), the ring $C=\sh{P}^1_{B/A}=(B\otimes_A B)/\mathfrak{J}^2$ is local, since the same holds for $C/(\mathfrak{J}/\mathfrak{J}^2)$, which is isomorphic to $B$. The hypothesis applied to $A'=C$ and $\mathfrak{J}'=\mathfrak{J}/\mathfrak{J}^2$ shows by definition that the map $d_{B/A}:B\to\Omega^1_{B/A}$ is zero (0, 20.4.6), hence that $\Omega^1_{B/A}=0$ by (0, 20.4.7).

\begin{proposition}[17.14.2]
\label{IV.17.14.2}
Let $Y$ be a locally Noetherian prescheme, $f:X\to Y$ a morphism locally of finite type, $x$ a point of $X$, $y=f(x)$. For $f$ to be smooth (resp. unramified, resp. étale) at the point $x$, it is necessary and sufficient that it satisfy the condition of (17.14.1) where we suppose moreover that the local ring $A'$ is Artinian, that $\mathfrak m'\mathfrak{J}'=0$, where $\mathfrak m'$ is the maximal ideal of $A'$, and that the residue field $A'/\mathfrak m'$ of $A'$ is equal to $k(x)$.
\end{proposition}

We again have to show that these conditions are sufficient in the case of smooth morphisms and in the case of unramified morphisms.

\oldpage[IV-4]{99}

\emph{(i) Case of smooth morphisms.} If we set $A=\sh{O}_{Y,y}$ and $B=\sh{O}_{X,x}$, we have to prove here, taking into account (17.5.3), that $B$ is a formally smooth $A$-algebra for the preadic topologies. Now, by virtue of (0, 22.1.4), this follows from the hypothesis when the condition $\mathfrak m'\mathfrak{J}'=0$ in it is replaced by $\mathfrak{J}'\subset\mathfrak m'$. But since $\mathfrak m'$ is nilpotent, we see that the condition of the statement is already sufficient by considering the rings $A'/\mathfrak m'^{,j}\mathfrak{J}'$ and reasoning by induction on $j$ as in (0, 19.4.3).

\emph{(ii) Case of unramified morphisms.} Let us take up again the notation of the proof of (17.14.1, (ii)); to prove that the ideal $\mathfrak{J}/\mathfrak{J}^2$ of $C$ is zero, note that $C$ is then a Noetherian local ring, being a local ring of the ring of an affine open subset of a sub-prescheme of $X\times_YX$, which is locally of finite type over $X$, hence also over $Y$. If $\mathfrak r$ is the maximal ideal of $C$, it therefore suffices to verify that we have $\mathfrak{J}/\mathfrak{J}^2\subset\mathfrak r^n$ for every $n$ (0$_{\mathrm I}$, 7.3.5), or equivalently that $\mathfrak r^n=\mathfrak r^n+(\mathfrak{J}/\mathfrak{J}^2)$. With the notation of (0, 20.4.6), this again means that for every $n$, the two composite maps $B\xrightarrow{p_1}C\to C/\mathfrak r^n$ and $B\xrightarrow{p_2}C\to C/\mathfrak r^n$ are identical; but since $C/\mathfrak r^n$ is Artinian, this identity follows from the hypothesis applied, by induction on $n$, to $A'=C/\mathfrak r^n$ and $\mathfrak{J}'=(\mathfrak r^n+(\mathfrak{J}/\mathfrak{J}^2))/\mathfrak r^n$.

\begin{remark}[17.14.3]
\label{IV.17.14.3}
Note that, under the hypotheses of (17.14.2), if moreover the field $k(x)$ is a \emph{finite} extension of $k(y)$, then the $A'$-modules $\mathfrak m'^{,j}/\mathfrak m'^{,j+1}$ are $k(y)$-vector spaces of finite rank, hence \emph{a fortiori} $A'$ is a finite $\sh{O}_{Y,y}$-algebra.
\end{remark}


\subsection{Case of preschemes over a base field.}
\EGAReaderTarget{ega4.c04.subsec.17-15}{subsection:IV.17.15}

Let us first recall (6.7.7, 6.7.8 and 6.8.1) the following

\begin{proposition}[17.15.1]
\label{IV.17.15.1}
Let $k$ be a field and $X$ a prescheme locally of finite type over $k$. For $X$ to be smooth at a point $x$, it is necessary and sufficient that for every radicial extension $k'$ of $k$ (or only for every finite radicial extension of $k$), the local ring $\sh{O}_{X,x}\otimes_k k'$ be regular. If $k(x)$ is a separable extension of $k$ (in particular if $k$ is perfect), it amounts to the same to say that $\sh{O}_{X,x}$ is regular.
\end{proposition}

\begin{corollary}[17.15.2]
\label{IV.17.15.2}
Under the conditions of (17.15.1), for $X$ to be smooth over $k$, it is necessary and sufficient that $X$ be geometrically regular over $k$ (6.7.6), which implies that $X$ is regular. If $k$ is perfect, then $X$ is smooth over $k$ if and only if $X$ is regular.
\end{corollary}

\begin{proposition}[17.15.3]
\label{IV.17.15.3}
Let $k$ be a field, $X$ a prescheme locally of finite type over $k$, $x$ a point of $X$, and $n=\dim_x(X)$. Let $(g_i)_{1\leq i\leq n}$ be a family of $n$ sections of $\sh{O}_X$ over $X$, and let $f:X\to Y=\operatorname{Spec}(k[T_1,\ldots,T_n])$ be the morphism corresponding to the $k$-homomorphism $k[T_1,\ldots,T_n]\to\Gamma(X,\sh{O}_X)$ that sends $T_i$ to $g_i$ (I, 2.2.4). The following conditions are equivalent (and imply that $X$ is smooth over $k$ at the point $x$, and consequently regular at the point $x$):
\begin{enumerate}
  \item[a)] $f$ is étale at the point $x$.
  \item[b)] The images of the $d_{X/k}g_i$ form a basis of the $\sh{O}_{X,x}$-module $(\Omega^1_{X/k})_x$.
  \item[c)] The images of the $d_{X/k}g_i$ generate the $\sh{O}_{X,x}$-module $(\Omega^1_{X/k})_x$.
\end{enumerate}
\end{proposition}

If $f$ is étale at the point $x$, then $X$ is smooth over $k$ at the point $x$, since $Y=\operatorname{Spec}(k[T_1,\ldots,T_n])$ is smooth over $k$ (17.3.8); the fact that a) implies b) is therefore a particular case of (17.11.3). Since b) trivially implies c), it remains to see that c) implies a). Let us first note

\oldpage[IV-4]{100}

that the hypothesis implies that the homomorphism $(f^*(\Omega^1_{Y/k}))_x\to(\Omega^1_{X/k})_x$ is surjective; hence, replacing $X$ by an open neighbourhood of $x$, we may suppose that the homomorphism $f^*(\Omega^1_{Y/k})\to\Omega^1_{X/k}$ is surjective (Bourbaki, \emph{Alg. comm.}, chap. II, § 5, n\textsuperscript{o} 1, prop. 2); consequently $f$ is \emph{unramified} (17.2.2). We will first see that we may restrict to the case where $x$ is \emph{rational over $k$}. Indeed, if we set $k'=k(x)$, and $X'=X\otimes_k k'$, $Y'=Y\otimes_k k'=\operatorname{Spec}(k'[T_1,\ldots,T_n])$, there exists a point $x'\in X'$ above $x$ such that $k(x')=k'$. To prove that $f$ is étale at the point $x$, it suffices to show that $f'=f_{(k')}:X'\to Y'$ is étale at the point $x'$ (17.7.1, (ii)); moreover, $f'$ is unramified (17.3.3, (iii)) and we have $\dim_{x'}(X')=n$ (4.2.7). In the same way, replacing $k'$ by an algebraically closed extension of $k'$, we may suppose that $k$ is \emph{algebraically closed}. Then set $y=f(x)$, $A=\sh{O}_{Y,y}$, $B=\sh{O}_{X,x}$; since the residue field of $B$ is equal to $k$, the same holds for that of $A$, so $x$ (resp. $y$) is a closed point of $X$ (resp. $Y$) (I, 6.4.2), and consequently we have $\dim(A)=\dim(B)=n$. Since $f$ is unramified, the homomorphism $\widehat A\to\widehat B$ is surjective (17.4.4, f'')); but since $\dim(\widehat A)=\dim(\widehat B)=n$, and $\widehat A$, being a regular ring, is integral, the homomorphism $\widehat A\to\widehat B$ is also injective (0, 16.3.10). This homomorphism is therefore bijective, which completes the proof that $f$ is étale at the point $x$ (17.6.3, e'')).

\begin{corollary}[17.15.4]
\label{IV.17.15.4}
Under the general hypotheses of (17.15.3), suppose moreover that $k(x)$ is a finite and separable extension of $k$ and that the germs $(g_i)_x$ belong to $\mathfrak m_x$. Then the conditions a), b), c) of (17.15.3) are also equivalent to each of the following:
\begin{enumerate}
  \item[d)] The $n$ germs $(g_i)_x$ generate the maximal ideal $\mathfrak m_x$ of $\sh{O}_{X,x}$.
  \item[d')] The ring $\sh{O}_{X,x}$ is regular and the $(g_i)_x$ form a regular system of parameters for this ring (0, 17.1.6).
\end{enumerate}
\end{corollary}

Indeed, d') trivially implies d). On the other hand, since $k(x)$ is a finite and separable extension of $k$, we have $\Omega^1_{k(x)/k}=0$ (0, 20.6.20), and $k(x)=\sh{O}_{X,x}/\mathfrak m_x$ is a formally smooth $k$-algebra for the discrete topologies (0, 19.6.1); hence the exact sequence (0, 20.5.14.1) applies to $A=k$, $B=\sh{O}_{X,x}$, $\mathfrak{R}=\mathfrak m_x$, and yields a canonical isomorphism
\[
\delta:\mathfrak m_x/\mathfrak m_x^2\xrightarrow{\sim}(\Omega^1_{X/k})_x\otimes_{\sh{O}_{X,x}}k(x).
\]
From this we first deduce the equivalence of conditions c) and d), taking into account Nakayama's lemma. On the other hand, if $f$ is étale at the point $x$, the ring $\sh{O}_{X,x}$ is regular and of dimension $n$, since $x$ is a closed point of $X$ (I, 6.4.2), and $n$ elements of $\mathfrak m_x$ that generate this maximal ideal then necessarily form a regular system of parameters for $\sh{O}_{X,x}$ (0, 17.1.6), which proves that a) implies d').

\begin{proposition}[17.15.5]
\label{IV.17.15.5}
Let $k$ be a field, $X$ a prescheme locally of finite type over $k$, $x$ a point of $X$, and $n=\dim_x(X)$. The following conditions are equivalent:
\begin{enumerate}
  \item[a)] $X$ is smooth over $k$ at the point $x$.
  \item[b)] $X$ is differentially smooth over $k$ at the point $x$.
  \item[c)] $(\Omega^1_{X/k})_x$ is a free $\sh{O}_{X,x}$-module of rank $n$.
  \item[d)] $(\Omega^1_{X/k})_x$ is an $\sh{O}_{X,x}$-module admitting a system of $n$ generators.

\oldpage[IV-4]{101}

  \item[e)] There exist a perfect extension $k'$ of $k$ and an open neighbourhood $U$ of $x$ in $X$ such that the prescheme $U\otimes_k k'$ is regular.
\end{enumerate}
\end{proposition}

The fact that $X$ is smooth over $k$ at the point $x$ implies the existence of an open neighbourhood $U$ of $x$ that is smooth over $k$; hence $U\otimes_k k'$ is regular for every extension $k'$ of $k$ (17.15.2), which proves that a) implies e). Conversely, e) implies a), since then $U\otimes_k k'$ is smooth over $k'$ (17.15.2), hence $U$ is smooth over $k$ (17.7.1, (ii)).

We have already proved that a) implies c) (17.10.2); c) implies d) trivially, and the fact that d) implies a) follows from (0, 20.4.7) and (17.15.3, c)).

Finally, we have already seen that a) implies b) (17.12.4). Conversely, to prove that b) implies a), we may restrict to the case where $x$ is rational over $k$, by considering, as in the proof of (17.15.3), a point $x'$ of $X'=X\otimes_k k'$ above $x$ such that $k(x')=k'=k(x)$, and using again (17.7.1, (ii)) and the fact that if $X$ is differentially smooth at the point $x$, then $X'$ is differentially smooth at the point $x'$ (16.10.4). Supposing therefore that $x$ is rational over $k$, the fact that $f$ is smooth at $x$ then follows from hypothesis b) and from (17.12.5) applied to the $k$-section $u:\operatorname{Spec}(k)\to X$ such that $u(\operatorname{Spec}(k))=\{x\}$.

\begin{corollary}[17.15.6]
\label{IV.17.15.6}
Let $X$ be a prescheme of finite type over a field $k$. For $X$ to be smooth over $k$, it is necessary and sufficient that the $\sh{O}_X$-module $\Omega^1_{X/k}$ be locally free, and that the local rings at the maximal points of $X$ be fields that are separable extensions of $k$ (this last condition being automatically satisfied if $k$ is a perfect field and $X$ a reduced prescheme).
\end{corollary}

The conditions are necessary, for if $X$ is smooth over $k$, it follows from (17.10.2) that $\Omega^1_{X/k}$ is locally free; on the other hand, $X$ is regular, hence \emph{a fortiori} reduced, so at every maximal point $x$ of $X$, $\sh{O}_{X,x}$ is a field, which must be a formally smooth $k$-algebra for the discrete topologies (17.5.1), hence a separable extension of $k$ (0, 19.6.1).

The conditions are sufficient. Indeed, we may restrict to the case where $X$ is connected, so that $\Omega^1_{X/k}$ is locally free of constant rank $n$. For every maximal point $x$ of $X$, we then have $\operatorname{rg}_{k(x)}\Omega^1_{k(x)/k}=n$ since $\sh{O}_{X,x}=k(x)$; since by hypothesis $k(x)$ is a separable extension of $k$, we have $\Upsilon_{k(x)/k}=0$ (0, 20.6.3), hence $\deg.\operatorname{tr}_k k(x)=n$ by Cartier's equality (0, 21.7.1). All the irreducible components of $X$ therefore have the same dimension $n$ (5.2.1), and we conclude that $X$ is smooth over $k$ at every point by virtue of the fact that c) implies a) in (17.15.5).

\begin{corollary}[17.15.7]
\label{IV.17.15.7}
If $k$ is of characteristic 0, then, for $X$ to be smooth over $k$ at the point $x$, it is necessary and sufficient that $(\Omega^1_{X/k})_x$ be a free $\sh{O}_{X,x}$-module.
\end{corollary}

This follows from (16.12.2) and from the equivalence of b) and a) in (17.15.5).

\begin{proposition}[17.15.8]
\label{IV.17.15.8}
Let $k$ be a field, $X$ a prescheme locally of finite type over $k$, and $x$ a point of $X$; set $n=\dim_x(X)$, $r=\dim(\sh{O}_{X,x})$, so that $n-r=\deg.\operatorname{tr}_k k(x)$ (5.2.3). Let $(g_i)_{1\leq i\leq n}$ be a family of $n$ sections of $\sh{O}_X$ over $X$ such that $(g_i)_x\in\mathfrak m_x$ for $1\leq i\leq r$; let $f:X\to Y=\operatorname{Spec}(k[T_1,\ldots,T_n])$ be the morphism corresponding to the $k$-homomorphism $k[T_1,\ldots,T_n]\to\Gamma(X,\sh{O}_X)$ that sends $T_i$ to $g_i$. The following conditions are equivalent (and imply that $X$ is smooth over $k$ at the point $x$ and that $k(x)$ is a separable extension of $k$):
\begin{enumerate}
  \item[a)] $f$ is étale at the point $x$.

\oldpage[IV-4]{102}

  \item[b)] The $(g_i)_x$ such that $1\leq i\leq r$ generate $\mathfrak m_x$ (and consequently form a regular system of parameters for $\sh{O}_{X,x}$ (0, 17.1.6)), and the images in $\Omega^1_{k(x)/k}$ of the elements $(d_{X/k}g_i)_x$ for $r+1\leq i\leq n$ generate $\Omega^1_{k(x)/k}$.
\end{enumerate}
\end{proposition}

Indeed, we have (0, 20.5.12.1) the exact sequence of $k(x)$-modules
\[
\label{IV.17.15.8.1}
\tag{17.15.8.1}
\mathfrak m_x/\mathfrak m_x^2\to(\Omega^1_{X/k})_x\otimes_{\sh{O}_{X,x}}k(x)\to\Omega^1_{k(x)/k}\to0
\]
and condition b) consequently implies that the $(d_{X/k}g_i)_x$ generate $(\Omega^1_{X/k})_x$, taking into account Nakayama's lemma; the fact that b) implies a) therefore follows from (17.15.3). Conversely, if a) holds, the $(d_{X/k}g_i)_x$ for $1\leq i\leq n$ form a basis of $(\Omega^1_{X/k})_x$ by virtue of (17.15.3). If $t_i$ ($1\leq i\leq n$) is the image of $(g_i)_x$ in $k(x)$, we conclude from the above that the $dt_i$ generate $\Omega^1_{k(x)/k}$ for $1\leq i\leq n$, and since by hypothesis $t_i=0$ for $1\leq i\leq r$, the $dt_i$ for $r+1\leq i\leq n$ already generate $\Omega^1_{k(x)/k}$. Since $\deg.\operatorname{tr}_k k(x)=n-r$, it follows from Cartier's equality (0, 21.7.1) that $\Upsilon_{k(x)/k}=0$, hence $k(x)$ is a separable extension of $k$ (0, 20.6.3), and the sequence of $k(x)$-vector spaces
\[
\label{IV.17.15.8.2}
\tag{17.15.8.2}
0\to\mathfrak m_x/\mathfrak m_x^2\to(\Omega^1_{X/k})_x\otimes_{\sh{O}_{X,x}}k(x)\to\Omega^1_{k(x)/k}\to0
\]
is exact ((0, 20.5.14) and (0, 19.6.1)); moreover, the $dt_i$ for $r+1\leq i\leq n$ form a basis of $\Omega^1_{k(x)/k}$, so none of the $t_i$ such that $r+1\leq i\leq n$ can be zero. This shows that the images of the $(g_i)_x$ in $\mathfrak m_x/\mathfrak m_x^2$ for $1\leq i\leq r$ necessarily generate $\mathfrak m_x/\mathfrak m_x^2$, hence the $(g_i)_x$ for $1\leq i\leq r$ generate $\mathfrak m_x$ by virtue of Nakayama's lemma, which completes the proof that a) implies b).

\begin{corollary}[17.15.9]
\label{IV.17.15.9}
Let $X$ be a prescheme locally of finite type over a field $k$, $x$ a point of $X$, and set $n=\dim_x(X)$, $r=\dim(\sh{O}_{X,x})$. The following conditions are equivalent:
\begin{enumerate}
  \item[a)] The ring $\sh{O}_{X,x}$ is regular and $k(x)$ is a separable extension of $k$.
  \item[b)] $X$ is smooth over $k$ at the point $x$, and the canonical homomorphism
  \[
  \mathfrak m_x/\mathfrak m_x^2\to(\Omega^1_{X/k})_x\otimes_{\sh{O}_{X,x}}k(x)
  \]
  is injective.
  \item[c)] There exist $n$ sections $g_i$ of $\sh{O}_X$ over an open neighbourhood $U$ of $x$ such that $(g_i)_x\in\mathfrak m_x$ for $1\leq i\leq r$ and such that the morphism $U\to\operatorname{Spec}(k[T_1,\ldots,T_n])$ corresponding to the $g_i$ (cf. (17.15.8)) is étale at the point $x$.
  \item[d)] There exist $n$ sections $g_i$ ($1\leq i\leq n$) of $\sh{O}_X$ over an open neighbourhood $U$ of $x$ such that the $(g_i)_x$ for $1\leq i\leq r$ generate $\mathfrak m_x$ and the images in $\Omega^1_{k(x)/k}$ of the $(d_{U/k}g_i)_x$ for $r+1\leq i\leq n$ generate $\Omega^1_{k(x)/k}$.
\end{enumerate}
\end{corollary}

The equivalence of c) and d) follows from (17.15.8). Moreover, we have seen in the proof of (17.15.8) that condition c) implies that $X$ is smooth over $k$ at the point $x$ and that the sequence (17.15.8.2) is exact; hence c) implies b). Condition b) implies that the ring $\sh{O}_{X,x}$ is regular, so $\mathfrak m_x/\mathfrak m_x^2$ is of rank $r$; moreover, $(\Omega^1_{X/k})_x$ is then a free $\sh{O}_{X,x}$-module of rank $n$ (17.10.2); since the sequence (17.15.7.2) is exact by hypothesis, $\Omega^1_{k(x)/k}$ is of rank $n-r=\deg.\operatorname{tr}_k k(x)$, and Cartier's equality shows that $\Upsilon_{k(x)/k}=0$, so $k(x)$ is separable over $k$ (0, 20.6.3); thus b) implies a). Finally, if a)

\oldpage[IV-4]{103}

holds, we again deduce from Cartier's equality that $\Omega^1_{k(x)/k}$ is of rank $n-r$. Since, on the other hand, the ring $\sh{O}_{X,x}$ is regular, the existence of the $g_i$ satisfying the conditions of d) is immediate.

\begin{remarks}[17.15.10]
\label{IV.17.15.10}
\begin{enumerate}
  \item[(i)] For a prescheme $X$ of finite type over $k$ to be smooth over $k$, it does not suffice that $\Omega^1_{X/k}$ be locally free, as is shown by the example where $X=\operatorname{Spec}(K)$, $K$ being a finite extension of $k$ that is not separable.
  \item[(ii)] When $k$ is not perfect, it can happen that $X$ is smooth over $k$ without $k(x)$ being separable over $k$. We obtain an example by taking $X=\operatorname{Spec}(k[T])$ and for $x$ the point corresponding to the principal prime ideal $(T^p-\lambda)$ ($p>0$ the characteristic of $k$, $\lambda\in k-k^p$).
  \item[(iii)] However, if $X$ is a prescheme smooth over $k$, the set of closed points of $X$ such that $k(x)$ is separable (and finite) over $k$ is dense in $X$. Indeed, let $f:X\to\operatorname{Spec}(k)$ be the structural morphism; for every $x_0\in X$, there are an open neighbourhood $U$ of $x_0$ in $X$ and a factorization of $f|U:U\xrightarrow{g}\operatorname{Spec}(k[T_1,\ldots,T_n])\xrightarrow{h}\operatorname{Spec}(k)$ where $g$ is étale (17.11.4). Since we may restrict to the case where $k$ is not perfect, hence infinite, the set of points of the open subset $g(U)$ that are rational over $k$ is non-empty; if $y$ is such a point and $x\in U$ a point above $y$, then $x$ is closed in $X$ and $k(x)$ is separable over $k(y)=k$ (17.6.2).
\end{enumerate}
\end{remarks}

\begin{proposition}[17.15.11]
\label{IV.17.15.11}
Let $X$ be a prescheme of finite type over a field $k$. The following conditions are equivalent:
\begin{enumerate}
  \item[a)] $X$ is étale over $k$.
  \item[b)] $X$ is unramified over $k$.
  \item[c)] $X$ is isomorphic to $\operatorname{Spec}(A)$, where $A$ is a finite and separable $k$-algebra.
\end{enumerate}
\end{proposition}

This follows from (17.6.2) and (17.4.2), taking into account the fact that condition b) implies here that $X$ is finite over $k$.

\begin{proposition}[17.15.12]
\label{IV.17.15.12}
Let $k$ be a field and $X$ a prescheme locally of finite type over $k$. For there to exist an everywhere dense open subset $U$ of $X$ that is smooth over $k$, it is necessary and sufficient that for every maximal point $x$ of $X$, $X$ be reduced at this point and $k(x)$ be a separable extension of $k$. For there to exist an everywhere dense open subset $U$ of $X$ such that $U_{\mathrm{red}}$ is smooth over $k$, it is necessary and sufficient that for every maximal point $x$ of $X$, $k(x)$ be a separable extension of $k$.
\end{proposition}

The second assertion evidently follows from the first. To say that there exists an everywhere dense open subset $U$ smooth over $k$ means that $X$ is smooth over $k$ at each of its maximal points $x$. For this it is necessary that $\sh{O}_{X,x}$ be regular (17.15.1) and \emph{a fortiori} reduced, hence a field, equal to $k(x)$; moreover (17.15.1), $k(x)\otimes_k k'$ must be regular for every radicial extension $k'$ of $k$, so $k(x)$ must be a separable extension of $k$ (4.3.5); the converse is immediate, by virtue of (17.15.1).

\begin{corollary}[17.15.13]
\label{IV.17.15.13}
Let $X$ be an algebraic prescheme over a field $k$. Then there exist an everywhere dense open subset $U$ in $X$ and a finite radicial extension $k'$ of $k$ such that $(U_{(k')})_{\mathrm{red}}$ is smooth over $k'$.
\end{corollary}

Indeed, there is such an extension $k'$ for which $(X_{(k')})_{\mathrm{red}}$ is geometrically reduced over $k'$ (4.6.6), which is equivalent to saying (4.6.1) that for every maximal point $x'$ of $X_{(k')}$, $k(x')$ is a separable extension of $k'$. We therefore conclude from (17.15.12) that there is an everywhere dense open subset $U'$ of $X_{(k')}$ such that $U'_{\mathrm{red}}$ is smooth over $k'$. But the \oldpage[IV-4]{104}morphism $X_{(k')}\to X$ is a homeomorphism (2.4.5), hence $U'$ is of the form $U_{(k')}$, where $U$ is an everywhere dense open subset in $X$.

\begin{proposition}[17.15.14]
\label{IV.17.15.14}
Let $X$ be an algebraic prescheme over $k$, of dimension $\leq1$. Then there exists a finite radicial extension $k'$ of $k$ such that the normalization (II, 6.3.8) of $(X_{(k')})_{\mathrm{red}}$ is smooth over $k'$.
\end{proposition}

Let us first prove the following two lemmas:

\begin{lemma}[17.15.14.1]
\label{IV.17.15.14.1}
Let $X$ be a reduced prescheme whose set of irreducible components is locally finite, and $f:Y\to X$ a morphism. For $Y$ to be $X$-isomorphic to the normalization $X'$ of $X$ (II, 6.3.8), it is necessary and sufficient that the following conditions be satisfied:
\begin{enumerate}
  \item[(i)] $Y$ is normal;
  \item[(ii)] $f$ is an integral and birational morphism.
\end{enumerate}
When there exists in $X$ a dense and normal open subset (which is always the case when $X$ is excellent (7.8.3), in particular when $X$ is locally of finite type over a field), we may replace condition (ii) by the following:
\begin{enumerate}
  \item[(ii bis)] $f$ is integral and there exists a dense open subset $U$ in $X$ such that $f^{-1}(U)$ is dense in $Y$ and the restriction $f^{-1}(U)\to U$ of $f$ is an isomorphism.
\end{enumerate}
\end{lemma}

The question being local on $X$, we may suppose that $X$ has only finitely many irreducible components; let $X_i$ ($1\leq i\leq m$) be the reduced sub-preschemes of $X$ having these components as underlying spaces. We know (II, 6.3.6) that $X'$ is the sum of the normalizations $X'_i$ of the $X_i$; each of the structural morphisms $f_i:X'_i\to X_i$ is therefore integral and birational, hence $f$ is integral and birational; it is immediate that (ii) implies (ii bis) when there exists a dense and normal open subset $V$ in $X$, by taking $U=V$. Conversely, if conditions (i) and (ii) are fulfilled, $Y$ has only finitely many irreducible components $Y_i$ ($1\leq i\leq m$), $Y_i$ dominating $X_i$ for each index $i$ (6.15.4); $Y$ is the sum of the $Y_i$ since it is normal, and the morphism $g_i:Y_i\to X_i$ through which the restriction $Y_i\to X$ of $f$ factors (I, 5.2.2) is birational; since $X_i$ and $Y_i$ are integral, it then follows from the fact that $g_i$ is integral and $Y_i$ normal that, for every affine open subset $V$ of $X_i$, $\Gamma(g_i^{-1}(V),\sh{O}_{Y_i})$ is the integral closure of $\Gamma(V,\sh{O}_{X_i})$, hence $Y$ is canonically identified with $X'$ (II, 6.3.4). Finally, if we suppose condition (ii bis) satisfied, we may restrict to the case where $U$ is the union of disjoint irreducible open subsets $U_i$ ($1\leq i\leq m$), so that $f^{-1}(U)$ is the disjoint union of the $V_i=f^{-1}(U_i)$. Since the $V_i$ are irreducible and $f^{-1}(U)$ is dense in $Y$, the $\overline{V_i}$ are the irreducible components of $Y$, and consequently $f$ is birational.

\begin{lemma}[17.15.14.2]
\label{IV.17.15.14.2}
Let $k$ be a field and $X$ an algebraic prescheme over $k$. Then there exists a finite radicial extension $k'$ of $k$ such that $(X_{(k')})_{\mathrm{red}}$ is geometrically reduced over $k'$ and its normalization is geometrically normal over $k'$.
\end{lemma}

Taking into account (4.6.6) and (I, 5.1.8), we may restrict to the case where $X$ is already geometrically reduced over $k$. Let $p$ be the characteristic exponent of $k$ and $k_1=k^{p^{-\infty}}$ the smallest perfect extension of $k$; $k_1$ is thus the inductive limit of the finite radicial extensions of $k$. Set $X_1=X_{(k_1)}$, which is reduced by hypothesis, and let $Y_1$ be its normalization; if $f_1:Y_1\to X_1$ is the structural morphism, $f_1$ is therefore finite (7.8.3, (vi)) and surjective, and there exists a dense open subset $U_1$ in $X_1$ such that $f_1^{-1}(U_1)=V_1$ is dense in $Y_1$ and \oldpage[IV-4]{105}the morphism $V_1\to U_1$, restriction of $f_1$, is an isomorphism (17.15.14.1). Applying (8.9.1) and (8.10.5, (vi) and (x)), we therefore first see that there exist a finite radicial extension $k''$ of $k$ and a finite surjective morphism $Y''\to X''=X_{(k'')}$ such that $Y_1=Y''\times_{X''}X_1=Y''\otimes_{k''}k_1$; since $Y_1$ is normal and $k_1$ is perfect, it follows from (6.7.7) that $Y''$ is geometrically normal over $k''$. Since the projections $X_1\to X''$ and $Y_1\to Y''$ are integral, surjective and radicial morphisms, they are homeomorphisms (2.4.5), and if $U''$ and $V''$ are the images of $U_1$ and $V_1$ in $X''$ and $Y''$ respectively, these are dense open subsets of $X''$ and $Y''$ respectively such that $V''=f''^{-1}(U'')$, where $f'':Y''\to X''$ is the structural morphism. Since we have $U_1=U''\otimes_{k''}k_1$, it follows from (8.10.5, (i)) that there exists a finite radicial extension $k'$ of $k''$ such that, if we set $X'=X\otimes_k k'=X''\otimes_{k''}k'$, $Y'=Y''\otimes_{k''}k'$, $f'=f''_{(k')}:Y'\to X'$, and if $U'$ and $V'$ are the images of $U''$ and $V''$ in $X'$ and $Y'$ respectively, the restriction $V'\to U'$ of $f'$ is an isomorphism. Since $Y'$ is normal and $f'$ is integral and birational, we conclude from (17.15.14.1) that $Y'$ is isomorphic to the normalization of $X'$, which proves the lemma, since $Y'$ is geometrically normal.

Let us now return to the proof of (17.15.14), and apply Lemma (17.15.14.2) to $k$ and $X$; since $\dim(X)\leq1$, we also have $\dim(X_{(k')})\leq1$ (4.1.4), and the normalization $Y'$ of $X_{(k')}$ is also of dimension $\leq1$ (5.4.2 and II, 7.4.6). To say that $Y'$ is geometrically normal over $k'$ then amounts to saying, by virtue of the definitions and of (II, 7.4.5), that $Y'$ is geometrically regular over $k'$, hence smooth over $k'$ (17.5.1), which proves (17.15.14).

\begin{proposition}[17.15.15]
\label{IV.17.15.15}
Let $f:X\to Y$ be a morphism locally of finite presentation. For $f$ to be smooth at a point $x\in X$, it is necessary and sufficient that $f$ be flat at the point $x$ and that $\Omega^1_f$ be an $\sh{O}_X$-module that is locally free in a neighbourhood of $x$, of rank at $x$ equal to $\dim_x f$.
\end{proposition}

The necessity of the conditions follows from (17.5.1) and (17.10.2). Conversely, if these conditions are satisfied, and if $y=f(x)$, it suffices to show (17.5.1) that $f^{-1}(y)$ is smooth over $k(y)$ at the point $x$; but this follows from the definition of $\dim_x f$ (17.10.1), from (16.4.5) and from (17.15.5).

\subsection{Quasi-sections of flat or smooth morphisms.}
\EGAReaderTarget{ega4.c04.subsec.17-16}{subsection:IV.17.16}

The statements of this subsection supplement those of (14.5), under flatness hypotheses.

\begin{proposition}[17.16.1]
\label{IV.17.16.1}
Let $f:X\to S$ be a flat morphism locally of finite presentation. Let $s\in S$, and let $x$ be a closed point of $X_s$ such that $\sh{O}_{X_s,x}$ is a Cohen--Macaulay ring; then there exist an open neighbourhood $U$ of $s$ in $S$ and a sub-prescheme $X'\subset f^{-1}(U)$ of $X$ such that $x\in X'$ and the morphism $X'\to U$, restriction of $f$, is flat, quasi-finite and of finite presentation.
\end{proposition}

The question being local on $X$ and $Y$, we may suppose $f$ of finite presentation. Let $(t_i)_{1\leq i\leq m}$ be a system of parameters of the local ring $\sh{O}_{X_s,x}$; there exist an affine open neighbourhood $V=\Spec(A)$ of $x$ in $X$ and $m$ sections $g_i$ of $\sh{O}_X$ over $V$ such that

\oldpage[IV-4]{106}

the images of the $(g_i)_x\in\sh{O}_{X,x}$ in $\sh{O}_{X_s,x}=\sh{O}_{X,x}/\mathfrak m_s\sh{O}_{X,x}$ are equal to the $t_i$. Let $X'$ be the closed sub-prescheme of $V$ defined by the ideal of $A$ generated by the $g_i$; since the sequence $(t_i)$ is regular by hypothesis (0, 16.5.7), it follows from (11.3.8) that, replacing $V$ if necessary by a smaller open neighbourhood, we may suppose that the morphism $X'\to S$, restriction of $f$, is flat and of finite presentation. On the other hand, since the $t_i$ form a system of parameters of $\sh{O}_{X_s,x}$, the ring $\sh{O}_{X'_s,x}$ is Artinian by definition, and since $x$ is closed in $X'_s$, we conclude that it is isolated in $X'_s$. Replacing $V$ if necessary by a smaller neighbourhood of $x$, we conclude, thanks to (13.1.4), that the morphism $X'\to S$ is quasi-finite.

\begin{corollary}[17.16.2]
\label{IV.17.16.2}
Let $f:X\to S$ be a morphism that is faithfully flat and locally of finite presentation. Then there exists a morphism $g:S'\to S$ that is faithfully flat, locally of finite presentation and locally quasi-finite, such that there exists an $S$-morphism $S'\to X$ (in other words, such that there exists an $S'$-section of $X'=X\times_S S'$ (I, 3.3.14)). If $S$ is quasi-compact (resp. quasi-compact and quasi-separated), we may suppose $S'$ affine (resp. $S'$ affine and the morphism $g$ quasi-finite).
\end{corollary}

For every $s\in S$, the fibre $X_s$ is non-empty by hypothesis and is a prescheme locally of finite type over $k(s)$; the set $U_s$ of points $x$ of $X_s$ at which $\sh{O}_{X_s,x}$ is a Cohen--Macaulay ring is open in $X_s$ (6.11.3) and is non-empty, since it contains the maximal points of $X_s$ (0, 16.5.1); consequently it contains a point $x_s$ closed in $X_s$ (10.4.7). Let $X'(s)$ be an affine sub-prescheme of $X$ containing $x_s$ and satisfying the conditions of (17.16.1). To obtain a prescheme $S'$ satisfying the conditions of the statement, it suffices to take the sum of the $X'(s)$, where $s$ runs through $S$. Since the morphism $X'(s)\to S$ is flat and locally of finite presentation, the image $U(s)$ of $X'(s)$ is open in $S$ (2.4.6); when $S$ is quasi-compact, we can therefore cover $S$ by finitely many $U(s_i)$, and the prescheme $S'$, sum of the $X'(s_i)$, still answers the question and is affine. If moreover $S$ is quasi-separated, we may suppose the open immersions $U(s_i)\to S$ quasi-compact (I, 2.7), hence of finite presentation (I, 6.2), and then the morphisms $X'(s_i)\to S$ are of finite presentation (I, 6.2), and consequently so is the morphism $S'\to S$.

\begin{corollary}[17.16.3]
\label{IV.17.16.3}
\begin{enumerate}
  \item[(i)] Let $f:X\to S$ be a smooth morphism. Let $s\in S$, and let $x$ be a closed point of $X_s$ such that the residue field $k(x)$ is separable over $k(s)$; then, in the conclusion of (17.16.1), we can take $X'$ such that the morphism $X'\to U$, restriction of $f$, is étale.
  \item[(ii)] Let $f:X\to S$ be a smooth surjective morphism. Then, in the conclusions of (17.16.2), we may moreover suppose that $g:S'\to S$ is étale.
\end{enumerate}
\end{corollary}

It is clear that (ii) is deduced from (i) as (17.16.2) was deduced from (17.16.1), taking into account the fact that, for every $s\in S$, since $X_s$ is non-empty and smooth over $k(s)$, there exists a closed point $x\in X_s$ such that $k(x)$ is separable over $k(s)$ (17.15.10, (iii)). It therefore suffices to prove (i). Note that the ring $\sh{O}_{X_s,x}$ is then regular; if we repeat the construction made in (17.16.1), taking for $(t_i)$ a regular system of parameters of $\sh{O}_{X_s,x}$, the ring $\sh{O}_{X'_s,x}$ is a field isomorphic to $k(x)$, hence separable over $k(s)$ by hypothesis. The conclusion then follows from (17.6.1, $c'$)).

\begin{proposition}[17.16.4]
\label{IV.17.16.4}
Let $S$ be a quasi-compact and quasi-separated prescheme, and $f:X\to S$ a surjective morphism locally of finite presentation. Then there exists a finite family $(S_i)_{i\in I}$

\oldpage[IV-4]{107}

of affine sub-preschemes of $S$, of finite presentation over $S$, pairwise disjoint, with union $S$, and having the following property: for each $i\in I$, there exist two finite surjective morphisms of finite presentation
\[
  S_i''\xrightarrow{g_i}S_i'\xrightarrow{h_i}S_i,
\]
where $h_i$ is étale and $g_i$ is flat and radicial, and an $S$-morphism $S_i''\to X$ (in other words, an $S_i''$-section of $X_i''=X\times_S S_i''$).
\end{proposition}

There is a finite covering $(U_i)_{1\leq i\leq n}$ of $S$ by affine open subsets; for every $i$, the intersection
\[
  W_i=U_i\cap\left(\bigcup_{j<i}U_j\right)
\]
is quasi-compact (I, 2.7), so there exists a closed sub-prescheme $V_i$ of $U_i$ having $U_i-W_i$ as underlying space and defined by an ideal of finite type of the ring of $U_i$; hence $V_i$ is an affine sub-prescheme of $S$ that is of finite presentation over $S$ (I, 6.2). It is clear that we may restrict to proving the proposition with $S$ replaced by $V_i$ and $X$ by $X\times_S V_i$. In other words, we may already suppose $S$ affine. On the other hand, $X$ is the union of affine open subsets $W_\alpha$, and the $f(W_\alpha)$ are constructible in $S$ (I, 8.4) and form a covering of $S$; consequently (I, 9.9) there exists a finite subfamily $(W_{\alpha_j})_{1\leq j\leq m}$ such that the $f(W_{\alpha_j})$ already form a covering of $S$. Then let $X'$ be the prescheme sum of the sub-preschemes induced on the open subsets $W_{\alpha_j}$ of $X$; it is immediate that it suffices to prove the proposition with $X$ replaced by $X'$, since an $S$-morphism $S_i''\to X'$ gives by composition an $S$-morphism $S_i''\to X'\to X$.

We may therefore suppose $X$ affine and $f$ of finite presentation. We can then (8.9.1) write $X$ in the form $X_0\times_{S_0}S$, where $S_0$ is Noetherian and $X_0\to S_0$ is a morphism of finite type, which we may suppose surjective (8.10.5, (vi)). If the proposition is proved for this morphism, it will follow immediately for $X$ by the base change $S\to S_0$. We may therefore suppose moreover that $S$ is Noetherian and $f$ is of finite type.

Let $s_h$ $(1\leq h\leq r)$ be the maximal points of $S$. Let us show that it suffices to prove the statement with $S$ replaced by a sufficiently small open neighbourhood $U_h$ of each of the $s_h$, and $X$ by the inverse image of $U_h$ in $X$. Indeed, suppose the proposition established in this case, and let us argue by Noetherian induction ($0_{\mathrm I}$, 2.2.2), supposing the statement established for every closed sub-prescheme of $S$ having an underlying space $\neq S$. We may suppose that the $U_h$ pairwise have no point in common; if $T$ is a closed sub-prescheme having $S-(\bigcup_h U_h)$ as underlying space, the induction hypothesis implies that the statement is true for $T$; since it is also true for each of the $U_h$, it is evidently true for $S$.

Since we may evidently (replacing $S$ by $S_{\mathrm{red}}$ and $X$ by $X\times_S S_{\mathrm{red}}$) suppose $S$ reduced, each of the $\sh{O}_{S,s_h}$ is a field $k(s_h)$. Suppose that we have proved the proposition when $S$ is the spectrum of a field and $f$ is of finite type. Then the existence of the $U_h$ follows from the method of (8.1.2, $a$)) and from (8.8.2, (i) and (ii)), (8.10.5, (vi), (vii) and (x)), (11.2.6, (ii)) and (17.7.8, (ii)).

Suppose therefore that $S=\Spec(k)$, where $k$ is a field, $X$ being of finite type over $k$ and $X\neq\varnothing$. Since $X$ is Noetherian, there exists in $X$ a closed point $x$ ($0_{\mathrm I}$, 2.1.3), so $k(x)$ is a finite extension of $k$ (I, 6.4.2). Consequently there is a finite separable extension $k'$ of $k$ such that $k''=k(x)$ is a finite radicial extension of $k'$. We then answer the question by taking $I$ reduced to one element, $S_i'=\Spec(k')$, $S_i''=\Spec(k'')$.

\oldpage[IV-4]{108}

\begin{corollary}[17.16.5]
\label{IV.17.16.5}
Let $f:X\to S$ be a surjective morphism locally of finite presentation. Then there exists a surjective morphism $g:S'\to S$, locally of finite presentation and locally quasi-finite, such that there exists an $S$-morphism $S'\to X$ (in other words, such that there exists an $S'$-section of $X'=X\times_S S'$). If $S$ is quasi-compact (resp. quasi-compact and quasi-separated), we may suppose $S'$ affine (resp. $S'$ affine and the morphism $g$ of finite presentation and quasi-finite).
\end{corollary}

It suffices to prove the corollary supposing $S$ affine: indeed, $S$ is the union of a family $(S_\alpha)$ of affine open subsets, and if for each $\alpha$, $S'_\alpha$ is affine and the morphism $g_\alpha:S'_\alpha\to S_\alpha$ answers the question and is of finite presentation and quasi-finite, then, taking for $S'$ the prescheme sum of the $S'_\alpha$, with $g:S'\to S$ coinciding with $g_\alpha$ on each of the $S'_\alpha$, this morphism answers the question; moreover, if $S$ is quasi-compact, we may suppose the family $(S_\alpha)$ finite, hence $S'$ affine; if in addition $S$ is quasi-separated, the immersions $S_\alpha\to S$ are of finite presentation (I, 6.2), hence so is $g$.

Consider therefore the case where $S$ is affine: we then form the finite morphisms of finite presentation $S_i''\to S_i$ of (17.16.4); since the immersions $S_i\to S$ are of finite presentation (the $S_i$ being affine), the morphisms $g_i'':S_i''\to S_i\to S$ are of finite presentation and quasi-finite, and we answer the question by taking $S'$ equal to the sum of the $S_i''$ and $g$ coinciding with $g_i''$ on each of the $S_i''$.

\begin{proposition}[17.16.6]
\label{IV.17.16.6}
Let $S$ be a quasi-compact and quasi-separated prescheme, and $f:X\to S$ a morphism of finite presentation such that, for every $s\in S$, $X_s=f^{-1}(s)$ is a prescheme proper over $k(s)$. Then there exists a finite family $(S_i)_{i\in I}$ of affine sub-preschemes of $S$, of finite presentation over $S$, pairwise disjoint, with union $S$, and such that, for every $i$, the morphism $X_i=X\times_S S_i\to S_i$ is proper and flat. If moreover, for every $s\in S$, $X_s$ is finite over $k(s)$, we can take the $S_i$ such that each of the morphisms $X_i\to S_i$ factors as
\[
  X_i\xrightarrow{u_i}X_i'\xrightarrow{h_i}S_i
\]
where $u_i$ and $h_i$ are finite and locally free (18.2.7), $h_i$ is étale, and $u_i$ is radicial and surjective.
\end{proposition}

We follow a course analogous to that of (17.16.4). We first reduce to the case where $S$ is affine and $f$ is flat, using the generic flatness theorem (8.9.5) (observe that the sub-preschemes $S_i'$ defined in the proof of (8.9.5) are affine). Since $f$ is of finite presentation, we can then write $X=X_0\times_{S_0}S$, where $S_0$ is Noetherian and $X_0\to S_0$ is a flat morphism of finite type (11.2.7); moreover, each fibre $(X_0)_{s_0}$ is again proper over $k(s_0)$, as follows from (2.7.1, (vii)), and we are therefore reduced to the case where $S$ is Noetherian. By Noetherian induction and application of the procedure of (8.1.2, $a$)) (this time also using (8.10.5, (xii))), we are finally reduced, in order to prove the first assertion, to the case where $S$ is the spectrum of a field, which follows trivially from the hypothesis (with $S_i=S$). We treat in the same way the case where $X_s$ is supposed finite over $k(s)$ for every $s\in S$ (this time one must use (2.7.1, (xv)), (8.10.5, (x), (iv), (vi) and (vii)), (17.7.8) and (2.1.12)), and we are reduced to proving the last assertion of the statement when $S=\Spec(k)$ is the spectrum of a field. But then $X=\Spec(A)$, where $A$ is a finite $k$-algebra (I, 6.4.4); so $A$ is the direct product of finite $k$-algebras $A_j$ that are local rings $(1\leq j\leq m)$; since $A_j$ is an Artinian ring and a

\oldpage[IV-4]{109}

$k$-algebra, it contains a subfield $K'_j$ canonically isomorphic to the residue field of $A_j$ (0, 19.6.2). We then take $X_i=X$ and $X'_i$ the sum of the $\Spec(K'_j)$, and it is clear that this answers the question, since for every $j$ we have two homomorphisms $K'_j\to A_j\to k(A_j)$ whose composite is an isomorphism.
